\documentclass[11pt]{book}

\usepackage[T1]{fontenc}
\usepackage[margin=1.1in]{geometry}
\usepackage{amsmath,amssymb,amsthm,mathtools}
\usepackage{bm}
\usepackage{booktabs}
\usepackage{longtable,array}
\usepackage{enumitem}
\usepackage{placeins}
\usepackage{xcolor}
\usepackage{tikz}
\usepackage{pgfplots}
\usepackage{hyperref}

\pgfplotsset{compat=1.18}
\usetikzlibrary{arrows.meta,calc,decorations.markings,decorations.pathreplacing,patterns,positioning}

\hypersetup{
  pdftitle={Mechanics},
  pdfauthor={GPT 5.5 and Xi Yin},
  colorlinks=true,
  linkcolor=blue,
  urlcolor=blue,
  citecolor=blue
}

\newtheorem{definition}{Definition}[chapter]
\newtheorem{proposition}{Proposition}[chapter]
\newtheorem{example}{Example}[chapter]
\newtheorem{remark}{Remark}[chapter]
\newtheorem{assumption}{Assumption}[chapter]
\newtheorem{derivation}{Derivation}[chapter]

\setlist[itemize]{itemsep=0.25em, topsep=0.35em}
\setlist[enumerate]{itemsep=0.25em, topsep=0.35em}

\newcommand{\R}{\mathbb{R}}
\newcommand{\N}{\mathbb{N}}
\newcommand{\Z}{\mathbb{Z}}
\newcommand{\dd}{\mathrm{d}}
\newcommand{\D}{\mathrm{D}}
\newcommand{\pd}{\partial}
\newcommand{\TQ}{TQ}
\newcommand{\TstarQ}{T^{*}Q}
\newcommand{\Lie}{\mathcal{L}}
\newcommand{\Ham}{H}
\newcommand{\Lag}{L}
\newcommand{\action}{\mathcal{S}}
\newcommand{\om}{\omega}
\newcommand{\J}{\mathcal{J}}
\newcommand{\eps}{\varepsilon}
\newcommand{\vect}[1]{\bm{#1}}
\newcommand{\inner}[2]{\left\langle #1,#2\right\rangle}
\newcommand{\norm}[1]{\left\lVert #1\right\rVert}
\newcommand{\Poisson}[2]{\left\{#1,#2\right\}}
\DeclareMathOperator{\Ad}{Ad}
\DeclareMathOperator{\ad}{ad}
\DeclareMathOperator{\diag}{diag}
\DeclareMathOperator{\id}{id}
\DeclareMathOperator{\SO}{SO}
\DeclareMathOperator{\SE}{SE}
\DeclareMathOperator{\tr}{tr}
\DeclareMathOperator{\grad}{grad}
\DeclareMathOperator{\curl}{curl}
\DeclareMathOperator{\divergence}{div}

\newcommand{\symbolentry}[3]{#1 & #2 & #3\\}


\title{Mechanics}
\author{GPT 5.5 and Xi Yin}
\date{\today}

\begin{document}

\frontmatter
\maketitle
\tableofcontents

\chapter*{Preface}

These notes present mechanics as a sequence of structures rather than as a
catalog of equations. The starting point is the configuration space and the
action. From there, the notes pass to phase space, symplectic geometry,
integrability, perturbation, chaos, fluids, elasticity, and gauge kinematics of
deformation. The examples are not side calculations: Kepler motion, the rigid
body, the standard map, the restricted Sun-Jupiter-asteroid model, laminar
flows, point vortices, and deforming bodies are the places where the abstract
language earns its keep.

The intended reading style is active. At each derivation, identify the space,
the variables, the assumptions, the conserved quantities, and the geometric
object being preserved or varied. At each figure, ask which equation or
reduction the picture is meant to make visible. At each numerical demo, ask
which continuous structure the discretization respects and which quantities
must be monitored.

The exposition uses a repeated rhythm. First state the model in words. Then
name the mathematical space on which the model lives. Then derive the equation
from a principle or balance law. Only after that should one choose coordinates,
plot a phase portrait, or run a simulation. This order is deliberate: it keeps
the reader from mistaking a convenient coordinate formula for the mechanics
itself.

\paragraph{How the chapters fit together.}
The notes are meant to be read linearly, but each chapter also has a precise
role in the overall architecture:
\[
\begin{array}{ccl}
\text{notation and assumptions} &\longrightarrow&
\text{well-defined models and symbols},\\
\text{Lagrangian mechanics} &\longrightarrow&
\text{equations from variation and symmetry},\\
\text{Hamiltonian mechanics} &\longrightarrow&
\text{phase space and symplectic structure},\\
\text{rigid body and tori} &\longrightarrow&
\text{integrability made concrete},\\
\text{perturbation and chaos} &\longrightarrow&
\text{what survives after integrability is broken},\\
\text{fluids and elastic bodies} &\longrightarrow&
\text{continuum mechanics and gauge kinematics},\\
\text{worked laboratories} &\longrightarrow&
\text{slow examples that test the formalism}.
\end{array}
\]
This dependency map is also a warning about shortcuts. For example, KAM theory
uses action-angle variables, which require the Liouville-torus construction,
which in turn rests on the Hamiltonian and symplectic formalism. Similarly, the
gauge theory of deforming bodies is clearest only after one has separated
configuration, symmetry, momentum, and connection in the rigid-body chapters.

\mainmatter
\chapter{Notation, Assumptions, And Reading Rules}
\label{chap:notation}

\section{Purpose Of These Notes}
\label{sec:notation-purpose}

These notes are written as lecture notes rather than as a compressed reference.
Each chapter states its assumptions, defines its symbols before use whenever
possible, and derives the central equations from a clearly named structure:
action principles, symplectic geometry, symmetry, or continuum balance laws.

The guiding convention is simple: when an equation appears, the reader should
know what space it lives on, which variables are independent, which assumptions
make it valid, and what physical information it carries.

The notes are meant to be self-contained at the level of a first serious pass
through advanced mechanics. When a result is used later, the earlier occurrence
should have supplied the model, the variables, the assumptions, and the
derivation strategy. The reader should never have to guess whether an equation
is a definition, an equation of motion, a constraint, a conservation law, or a
coordinate representation of a geometric object.

Self-contained does not mean that every calculation starts from elementary
calculus. It means that every new structure is introduced with enough context
to be used responsibly. For example, when the notes invoke a symplectic form,
they state the phase space and the sign convention; when they invoke a stress
tensor, they state the balance law and the constitutive assumption; when they
invoke a connection, they say which frame choice is being compared.

The recurring reading questions are:
\[
\begin{gathered}
  \text{What is the configuration or phase space?}\\
  \text{What structure is being preserved or varied?}\\
  \text{Which assumptions make the calculation legal?}\\
  \text{What physical or geometric information survives a change of coordinates?}
\end{gathered}
\]
These questions are deliberately repeated in different contexts. They are the
thread connecting particles, rigid bodies, celestial mechanics, fluids, and
deforming continua.

\section{Exposition Standard Used In The Notes}
\label{sec:exposition-standard}

Each major derivation is written to answer five questions in order:
\[
\begin{array}{rcl}
\text{model} &:& \text{what physical idealization is being made?}\\
\text{space} &:& \text{where do the variables live?}\\
\text{principle} &:& \text{what variational, Hamiltonian, or balance statement is used?}\\
\text{calculation} &:& \text{which algebraic or analytic step produces the equation?}\\
\text{interpretation} &:& \text{what does the result say invariantly?}
\end{array}
\]
When a derivation feels long, this list is the intended guide. A displayed
equation should not be read as an isolated fact; it is either part of the
model, a definition of notation, a consequence of a principle, or an
interpretation of the result. The surrounding prose is meant to identify which
role the equation is playing.

The notes also distinguish three kinds of equality. A definition introduces a
symbol, as in \(p_i=\partial L/\partial\dot q^i\). An identity holds because of
previous definitions, as in \(\omega=-\dd\Theta\). An equation of motion holds
only along physical trajectories, as in \(\dot p_i=-\partial H/\partial q^i\).
Many confusions in mechanics come from silently moving an equation from one
category to another.

\section{How Notation Is Controlled}
\label{sec:notation-control}

Notation in mechanics is necessarily overloaded: \(I\) can mean a time
interval, an inertia tensor, or an action variable; \(\omega\) can mean a
symplectic form, an angular velocity, a frequency vector, or vorticity. The
notes therefore use the following discipline.

\begin{enumerate}
\item A symbol has no permanent meaning unless it appears in the global ledger
below or in the notation table at the beginning of a chapter.
\item If a standard symbol is reused with a different meaning, the local
definition overrides the global convention only inside that section or chapter.
\item Bold or vector notation indicates Euclidean vectors only when the ambient
space and metric have been named.
\item Upper/lower indices in finite-dimensional mechanics indicate coordinate
indices on a manifold and its cotangent bundle. In continuum chapters, capital
indices such as \(A,B\) label material coordinates and lower-case indices such
as \(i,j\) label spatial coordinates.
\item Every theorem, derivation, example, and demo should be readable together
with the nearest notation table and the assumptions stated immediately before
the calculation.
\end{enumerate}

\section{Standing Mathematical Assumptions}
\label{sec:standing-assumptions}

\begin{assumption}[Smoothness]
Unless explicitly stated otherwise, all functions, curves, vector fields, and
maps are assumed smooth enough for the differentiations and integrations being
performed. For most finite-dimensional chapters, \(C^2\) regularity is enough.
For continuum chapters, additional differentiability is assumed locally.
\end{assumption}

\begin{assumption}[Finite-dimensional mechanics]
When the configuration space is a finite-dimensional manifold \(Q\), a motion is
a curve \(q:I\to Q\) on a time interval \(I=[t_0,t_1]\). The velocity
\(\dot q(t)\) lies in the tangent space \(T_{q(t)}Q\). Local coordinates are
written \(q^1,\ldots,q^n\).
\end{assumption}

\begin{assumption}[Endpoint variations]
For Hamilton's principle in its basic form, variations of the path satisfy
\(\delta q(t_0)=\delta q(t_1)=0\). Free-endpoint and boundary-force variants are
important, but they are not the default convention.
\end{assumption}

\begin{assumption}[Regular Lagrangians]
The Legendre transform from velocities to momenta is assumed locally invertible
whenever we pass from \(L(q,\dot q,t)\) to \(H(q,p,t)\). In coordinates this
means the matrix
\[
  \frac{\partial^2L}{\partial \dot q^i\partial \dot q^j}
\]
is nonsingular. Singular Lagrangians lead to constraints and gauge theories in
the Dirac sense; those require a separate treatment.
\end{assumption}

\begin{assumption}[Einstein summation]
Repeated upper/lower coordinate indices are summed unless stated otherwise:
\[
  p_i\dot q^i=\sum_{i=1}^{n}p_i\dot q^i.
\]
In Euclidean vector formulas, dots and cross products use the standard metric
and orientation.
\end{assumption}

\section{Global Notation Ledger}
\label{sec:global-notation-ledger}

This table is a reference ledger, not a replacement for local definitions. The
``first reference'' column points to the chapter where the symbol family is
used most systematically.

\small
\begin{longtable}{p{0.18\textwidth}p{0.52\textwidth}p{0.22\textwidth}}
\toprule
Symbol & Meaning and assumptions & First reference\\
\midrule
\endhead
\bottomrule
\endfoot
\(t\) & time variable; dots denote \(\dd/\dd t\) unless a section says otherwise & all chapters\\
\(I=[t_0,t_1]\) & time interval for a path or field evolution & Chapter~\ref{chap:lag-ham}\\
\(Q\) & finite-dimensional configuration manifold & Chapter~\ref{chap:lag-ham}\\
\(q(t)\) & configuration path \(I\to Q\) & Chapter~\ref{chap:lag-ham}\\
\(q^i\) & local coordinates on \(Q\), \(i=1,\ldots,n\) & Chapter~\ref{chap:lag-ham}\\
\(\dot q^i\) & generalized velocities \(\dd q^i/\dd t\) & Chapter~\ref{chap:lag-ham}\\
\(TQ\) & tangent bundle, the space of \((q,\dot q)\) & Chapter~\ref{chap:lag-ham}\\
\(T^*Q\) & cotangent bundle, the phase space of \((q,p)\) & Chapter~\ref{chap:lag-ham}\\
\(L\) & Lagrangian, usually \(T-V\) but not assumed to be natural unless stated & Chapter~\ref{chap:lag-ham}\\
\(S[q]\) & configuration-space action functional \(\int L\,\dd t\) & Chapter~\ref{chap:lag-ham}\\
\(\eta\) & variation vector field \(\delta q\) along a path & Chapter~\ref{chap:lag-ham}\\
\(p_i\) & canonical momentum \(\partial L/\partial \dot q^i\) & Chapter~\ref{chap:lag-ham}\\
\(H\) & Hamiltonian on phase space; in autonomous conservative examples it is energy & Chapter~\ref{chap:lag-ham}\\
\(\Theta\) & canonical one-form on \(T^*Q\), \(\Theta=p_i\,\dd q^i\) & Chapter~\ref{chap:lag-ham}\\
\(\omega\) & symplectic form in Hamiltonian chapters; other uses are local and stated & Chapter~\ref{chap:lag-ham}\\
\(\{f,g\}\) & Poisson bracket of observables \(f,g\) & Chapter~\ref{chap:lag-ham}\\
\(X_H\) & Hamiltonian vector field determined by \(\iota_{X_H}\omega=\dd H\) & Chapter~\ref{chap:lag-ham}\\
\(G\) & Lie group in symmetry sections; gravitational constant in celestial sections & Chapters~\ref{chap:lag-ham}, \ref{chap:perturbation-chaos}\\
\(\mathfrak g\) & Lie algebra of a Lie group \(G\) & Chapters~\ref{chap:lag-ham}, \ref{chap:elastic-gauge}\\
\(\xi_Q\) & infinitesimal generator on \(Q\) associated with \(\xi\in\mathfrak g\) & Chapter~\ref{chap:lag-ham}\\
\(J_\xi\) & Noether conserved quantity associated with \(\xi\) & Chapter~\ref{chap:lag-ham}\\
\(R\) & rotation matrix in rigid-body chapters; radius or gas constant only when locally declared & Chapter~\ref{chap:rigid-integrability}\\
\(\SO(3)\) & group of orientation-preserving orthogonal \(3\times3\) matrices & Chapter~\ref{chap:rigid-integrability}\\
\(\Omega\) & body angular velocity vector; not the same as symplectic \(\omega\) & Chapter~\ref{chap:rigid-integrability}\\
\(\widehat\Omega\) & skew matrix satisfying \(\widehat\Omega v=\Omega\times v\) & Chapter~\ref{chap:rigid-integrability}\\
\(\omega_{\mathrm{space}}\) & spatial angular velocity when needed to avoid conflict with symplectic \(\omega\) & Chapter~\ref{chap:rigid-integrability}\\
\(I\) & inertia tensor or action variable depending on context; local table fixes meaning & Chapters~\ref{chap:rigid-integrability}, \ref{chap:perturbation-chaos}\\
\(M\) & body angular momentum in rigid-body chapters & Chapter~\ref{chap:rigid-integrability}\\
\(J\) & spatial angular momentum in rigid-body/celestial chapters; Jacobian determinant in continuum chapters & Chapters~\ref{chap:rigid-integrability}, \ref{chap:fluid}\\
\(F_i\) & commuting first integrals in Liouville theory; deformation gradient \(F^i{}_A\) in continua & Chapters~\ref{chap:rigid-integrability}, \ref{chap:elastic-gauge}\\
\(\gamma_j\) & homology cycle on an invariant torus; material loop in fluid contexts & Chapters~\ref{chap:rigid-integrability}, \ref{chap:fluid}\\
\(I_j\) & action variable \((2\pi)^{-1}\oint_{\gamma_j}p_i\,\dd q^i\) & Chapter~\ref{chap:rigid-integrability}\\
\(\theta^j\) & angle variable conjugate to \(I_j\); ordinary polar/Euler angles are locally defined & Chapters~\ref{chap:rigid-integrability}, \ref{chap:perturbation-chaos}\\
\(\mathbb T^n\) & \(n\)-torus, usually \((\R/2\pi\Z)^n\) & Chapters~\ref{chap:rigid-integrability}, \ref{chap:perturbation-chaos}\\
\(\varepsilon\) & small perturbation parameter & Chapter~\ref{chap:perturbation-chaos}\\
\(H_0,H_1\) & integrable Hamiltonian and perturbing Hamiltonian & Chapter~\ref{chap:perturbation-chaos}\\
\(\omega(I)\) & frequency vector in action-angle variables & Chapter~\ref{chap:perturbation-chaos}\\
\(k\) & integer Fourier/resonance vector in KAM sections; spring/gravity constants are locally defined & Chapter~\ref{chap:perturbation-chaos}\\
\(\alpha,\tau\) & Diophantine constants; \(\alpha\) also appears as a wave number in stability sections only locally & Chapter~\ref{chap:perturbation-chaos}\\
\(K\) & standard-map kick strength & Chapter~\ref{chap:perturbation-chaos}\\
\(\rho\) & mass density in fluids and continua; radial coordinate only when locally stated & Chapters~\ref{chap:fluid}, \ref{chap:elastic-gauge}\\
\(u(x,t)\) & Eulerian fluid velocity field & Chapter~\ref{chap:fluid}\\
\(p(x,t)\) & fluid pressure in fluid chapters; canonical momentum elsewhere & Chapter~\ref{chap:fluid}\\
\(\sigma\) & Cauchy stress tensor in continuum chapters & Chapters~\ref{chap:fluid}, \ref{chap:elastic-gauge}\\
\(\mu,\nu\) & dynamic and kinematic viscosity in fluids; reduced mass \(\mu\) in Kepler sections & Chapters~\ref{chap:fluid}, \ref{chap:perturbation-chaos}\\
\(\mathrm{Re}\) & Reynolds number \(UL/\nu\) & Chapter~\ref{chap:fluid}\\
\(B\) & reference body/material domain in elasticity; bending stiffness in rod sections when locally defined & Chapter~\ref{chap:elastic-gauge}\\
\(X^A\) & material coordinates on \(B\) & Chapter~\ref{chap:elastic-gauge}\\
\(x^i\) & spatial coordinates in \(\R^3\) & Chapter~\ref{chap:elastic-gauge}\\
\(\varphi(X,t)\) & deformation map from reference body to space & Chapters~\ref{chap:fluid}, \ref{chap:elastic-gauge}\\
\(F^i{}_A\) & deformation gradient \(\partial_A\varphi^i\) & Chapter~\ref{chap:elastic-gauge}\\
\(C_{AB},E_{AB}\) & right Cauchy-Green tensor and Green strain & Chapter~\ref{chap:elastic-gauge}\\
\(P^A{}_i\) & first Piola-Kirchhoff stress & Chapter~\ref{chap:elastic-gauge}\\
\(\mathcal S\) & shape space for deforming bodies & Chapter~\ref{chap:elastic-gauge}\\
\(\mathcal A\) & mechanical/gauge connection on shape space & Chapter~\ref{chap:elastic-gauge}\\
\(\mathcal F\) & curvature of the connection \(\mathcal A\) & Chapter~\ref{chap:elastic-gauge}\\
\end{longtable}
\normalsize

\section{Common Overloads And How To Read Them}
\label{sec:overloaded-symbols}

\begin{center}
\begin{tabular}{p{0.16\textwidth}p{0.76\textwidth}}
\toprule
Symbol & Reading rule\\
\midrule
\(\omega\) & In symplectic sections it is the symplectic form. In rigid body
sections the spatial angular velocity is written in prose as spatial angular
velocity or \(\omega_{\mathrm{space}}\) when ambiguity matters. In fluids it is
vorticity. In KAM sections \(\omega(I)\) is a frequency vector.\\
\(I\) & In Chapter~\ref{chap:notation} and Chapter~\ref{chap:lag-ham}, \(I\)
may be a time interval. In rigid body mechanics \(I\) is an inertia tensor or
principal moment. In action-angle theory \(I_j\) is an action variable.\\
\(J\) & In mechanics \(J\) is usually angular momentum or a momentum map. In
continuum mechanics \(J=\det F\) is the deformation Jacobian.\\
\(p\) & In Hamiltonian mechanics \(p_i\) is canonical momentum. In fluid
mechanics \(p(x,t)\) is pressure.\\
\(F\) & In integrability \(F_i\) are first integrals. In elasticity
\(F^i{}_A\) is the deformation gradient.\\
\(A\) & In Kepler theory \(A\) may denote the Laplace-Runge-Lenz vector. In
gauge theory \(A\) or \(\mathcal A\) denotes a connection.\\
\(\lambda\) & A Lagrange multiplier in constrained mechanics, a second
viscosity coefficient in fluids, or a Lame modulus in elasticity; the section
defines which one.\\
\bottomrule
\end{tabular}
\end{center}

\section{A Map Of The Course}
\label{sec:course-map}

The course moves through four enlargements of the same idea. The diagram is not
a dependency graph; it is a loop that the notes revisit at increasing levels of
structure. A variational principle produces equations, Hamiltonian geometry
organizes their phase space, perturbation theory tests the stability of that
organization, and continuum mechanics returns to the same questions with
fields instead of finitely many coordinates.

\begin{figure}[htbp]
\centering
\begin{tikzpicture}[
  node distance=1.1cm and 1.35cm,
  box/.style={draw, rounded corners=2pt, align=center, minimum width=3.3cm,
  minimum height=0.92cm, fill=blue!4},
  smallbox/.style={draw, rounded corners=2pt, align=center, minimum width=2.8cm,
  minimum height=0.75cm, fill=gray!8},
  arr/.style={-{Latex[length=2.5mm]}, thick},
  softarr/.style={-{Latex[length=2.2mm]}, thick, gray!70}
]
\node[box] (lag) {Lagrangian mechanics\\configuration, action, variation};
\node[box, right=of lag] (ham) {Hamiltonian mechanics\\phase space, symplectic form};
\node[box, below=of ham] (chaos) {Integrability and chaos\\tori, resonances, transport};
\node[box, left=of chaos] (cont) {Fluids and elasticity\\fields, stress, frames};
\node[smallbox, below=of lag, xshift=-0.35cm] (examples) {Kepler and rigid body\\first exact models};
\node[smallbox, above=of ham, xshift=0.35cm] (demos) {Demos and diagnostics\\drift, sections, probability};
\draw[arr] (lag) -- (ham);
\draw[arr] (ham) -- (chaos);
\draw[arr] (chaos) -- (cont);
\draw[arr] (cont) -- (lag);
\draw[softarr] (examples) -- (lag);
\draw[softarr] (examples) -- (chaos);
\draw[softarr] (demos) -- (ham);
\draw[softarr] (demos) -- (chaos);
\end{tikzpicture}
\caption{The course loop. The central theories are connected by examples and
computational diagnostics: exact models teach the geometry, while numerical
demos test which structures survive discretization or perturbation.}
\label{fig:course-loop}
\end{figure}

\section{Units And Dimensions}
\label{sec:units-dimensions}

Dimensional consistency is treated as a mathematical constraint. In celestial
mechanics demos we often use astronomical units, years, and solar masses, so
\[
  G = 4\pi^2
\]
when the central mass is measured in solar masses. In rigid body mechanics,
moments of inertia have dimensions \(ML^2\), angular velocity has dimension
\(T^{-1}\), angular momentum has dimension \(ML^2T^{-1}\), and kinetic energy
has dimension \(ML^2T^{-2}\).

\section{How To Read The Figures}
\label{sec:read-figures}

Figures in the notes are schematic unless explicitly generated from a numerical
demo. Their purpose is to identify variables, spaces, and geometric relations.
For quantitative plots, the corresponding Python or Mathematica script should be
treated as part of the figure caption: the model assumptions and numerical
parameters matter.

\section{How To Read Long Derivations}
\label{sec:read-long-derivations}

Long derivations in mechanics usually have two layers. The first layer is
formal: differentiate an action, apply a balance law, invert a Legendre
transform, or change coordinates. The second layer is structural: identify the
conserved quantity, the reduced phase space, the constraint force, or the
connection hidden in the algebra. The notes spell out both layers whenever the
second layer is easy to miss.

A useful way to read each derivation is to keep a small ledger:
\[
\begin{array}{c|c}
\text{input} & \text{output}\\
\hline
\text{model assumptions} & \text{equation of motion or constraint}\\
\text{symmetry} & \text{conserved quantity or reduced variable}\\
\text{constitutive law} & \text{closed continuum PDE}\\
\text{choice of frame} & \text{connection and curvature}
\end{array}
\]
Most mistakes in mechanics are ledger mistakes: using an equation outside the
assumptions that produced it, changing a frame without transforming the
components, or interpreting a reduced equation as if no reduction had occurred.

When a derivation contains a local coordinate calculation, the conclusion should
be read as coordinate-independent only after the geometric object has been
identified. For example, the symbols \(\Gamma^\ell{}_{ij}\), \(\lambda_a\),
\(A_A\), and \(\omega^a{}_b\) are coordinate or frame components; the geodesic
equation, constraint reaction, gauge connection, and torsion are the underlying
objects.

\chapter{Geometric Lagrangian And Hamiltonian Mechanics}
\label{chap:lag-ham}

\section{Mechanical Data}
\label{sec:lag-ham-mechanical-data}

The first question in any mechanics problem is not ``what is the force?'' but
``what is the space of possible configurations?'' Once that space is named, the
rest of the theory becomes coordinate-independent.

\begin{definition}[Finite-dimensional mechanical system]
A regular finite-dimensional Lagrangian mechanical system consists of:
\begin{enumerate}
\item a smooth \(n\)-dimensional configuration manifold \(Q\);
\item a time interval \(I=[t_0,t_1]\);
\item a Lagrangian \(L:TQ\times I\to \R\);
\item a class of admissible curves \(q:I\to Q\), possibly satisfying constraints.
\end{enumerate}
The action of an admissible curve is
\[
  S[q]=\int_{t_0}^{t_1} L(q(t),\dot q(t),t)\,\dd t.
\]
\end{definition}

\begin{assumption}[Coordinate patch for derivations]
Derivations in this chapter are written in one coordinate chart
\((q^1,\ldots,q^n)\). The equations obtained are tensorial or intrinsic when
the objects used to build them are intrinsic. Coordinates are a calculation
device, not an additional physical structure.
\end{assumption}

\begin{table}[htbp]
\centering
\caption{Core notation for Chapter~\ref{chap:lag-ham}. Global conventions and
overloaded symbols are listed in Chapter~\ref{chap:notation}.}
\label{tab:lag-ham-notation}
\small
\begin{tabular}{p{0.18\textwidth}p{0.45\textwidth}p{0.28\textwidth}}
\toprule
Symbol & Definition & Interpretation or first use\\
\midrule
\(Q\) & smooth configuration manifold & space of positions\\
\(I=[t_0,t_1]\) & time interval & domain of a path\\
\(q:I\to Q\) & configuration path & physical motion candidate\\
\(q^i(t)\) & local coordinates of \(q(t)\) & generalized positions\\
\(\dot q^i(t)\) & \(\dd q^i/\dd t\) & generalized velocities\\
\(\eta^i(t)\) & \(\delta q^i(t)\) & infinitesimal virtual displacement\\
\(L:TQ\times I\to\R\) & Lagrangian & input to Hamilton's principle\\
\(S[q]\) & \(\int L\,\dd t\) & action functional\\
\(p_i\) & \(\partial L/\partial \dot q^i\) & canonical momentum\\
\(E_L\) & \(p_i\dot q^i-L\) & Lagrangian energy\\
\(g_{ij}\) & mass metric & natural Lagrangian \(T-V\)\\
\(\Gamma^\ell{}_{jk}\) & Christoffel symbols of \(g\) & inertial terms in coordinates\\
\(\Phi^a(q,t)\) & holonomic constraints & admissible submanifold\\
\(\lambda_a\) & constraint multipliers & reaction forces\\
\(G,\mathfrak g\) & symmetry group and Lie algebra & Noether theorem\\
\(\xi_Q\) & infinitesimal generator & vector field on \(Q\)\\
\(J_\xi\) & \(p_i\xi_Q^i\) & Noether charge\\
\(r,\rho,\theta\) & Kepler vector, radius, polar angle & central-force example\\
\(\ell\) & \(\mu\rho^2\dot\theta\) & angular momentum magnitude\\
\(R\) & attitude matrix & rigid-body example\\
\(\Omega,M\) & body angular velocity and body angular momentum & \(M=I\Omega\)\\
\(H\) & \(p_i\dot q^i-L\) after Legendre transform & Hamiltonian\\
\(\Theta\) & \(p_i\,\dd q^i\) & canonical one-form\\
\(\omega\) & \(-\dd\Theta=\dd q^i\wedge\dd p_i\) & symplectic form\\
\(X_H\) & \(\iota_{X_H}\omega=\dd H\) & Hamiltonian vector field\\
\(\Phi_H^t\) & Hamiltonian flow map & phase-space evolution\\
\(\Psi_{\Delta t}\) & numerical one-step map & integrator output\\
\bottomrule
\end{tabular}
\end{table}

\paragraph{Roadmap.}
The chapter has three jobs. First it derives the Euler-Lagrange equations from
Hamilton's principle, including boundary terms, constraints, and Noether
charges. Second it treats Kepler motion and the rigid body as the first serious
examples of the formalism. Third it passes through the Legendre transform to
Hamiltonian mechanics, where phase space, symplectic form, Poisson brackets,
and canonical transformations become the basic language. The point is not to
replace Newton's law by harder notation; it is to make clear which parts of
mechanics are coordinate choices and which parts are geometric structure.

The chapter should be read as a change in viewpoint. In the Lagrangian picture,
the primitive object is a path in configuration space and the variational
principle selects the physical path. In the Hamiltonian picture, the primitive
object is a point in phase space and the symplectic form turns the Hamiltonian
function into a vector field. The two pictures are equivalent only under the
regularity assumption stated in Chapter~\ref{chap:notation}; this is why the
Legendre transform is a theorem-sized step, not a change of notation.

\paragraph{Derivation checkpoints.}
Whenever this chapter differentiates an action, check three things before
following the algebra: which variables are varied, which endpoint conditions
are imposed, and which boundary terms survive. Whenever it changes from
\((q,\dot q)\) to \((q,p)\), check the velocity Hessian. Whenever it invokes a
symmetry, check whether the symmetry acts on configuration space, phase space,
or both. These checks are small, but they prevent the most common mistakes:
using an inadmissible virtual displacement, assuming a singular Legendre map is
invertible, or calling a quantity conserved before verifying the symmetry
hypothesis.

\section{Variations And The Euler-Lagrange Equations}
\label{sec:euler-lagrange}

Let \(q_\eps(t)\) be a smooth one-parameter family of curves with
\(q_0(t)=q(t)\). The variation vector field along \(q\) is
\[
  \eta(t)=\left.\frac{\partial q_\eps(t)}{\partial \eps}\right|_{\eps=0},
  \qquad
  \eta^i(t)=\left.\frac{\partial q_\eps^i(t)}{\partial \eps}\right|_{\eps=0}.
\]
For fixed-endpoint variations we impose
\[
  \eta(t_0)=0,\qquad \eta(t_1)=0.
\]
The variational calculation below has one essential idea. The path \(q(t)\) is
a critical point of the action when every infinitesimal deformation of the path
changes the action only at second order. Differentiating the action produces a
term proportional to \(\dot\eta\); integration by parts moves the derivative
onto the momentum, and the endpoint condition kills the boundary term. What is
left must vanish for every interior virtual displacement \(\eta\), which forces
the Euler-Lagrange equations pointwise in time.

\begin{figure}[htbp]
\centering
\begin{tikzpicture}[scale=1.0, >=Latex]
  \draw[->] (-0.3,0) -- (7.2,0) node[right] {\(t\)};
  \draw[->] (0,-0.2) -- (0,3.0) node[above] {configuration coordinate};
  \draw[densely dotted] (0.5,0) -- (0.5,2.75);
  \draw[densely dotted] (6.5,0) -- (6.5,2.75);
  \draw[thick, blue] plot[smooth] coordinates {(0.5,0.7) (1.8,1.25) (3.2,1.0) (4.7,2.0) (6.5,2.35)};
  \draw[thick, red, dashed] plot[smooth] coordinates {(0.5,0.7) (1.8,1.55) (3.2,1.25) (4.7,1.75) (6.5,2.35)};
  \draw[->, red] (3.2,1.0) -- (3.2,1.25) node[midway,right] {\(\eta(t)\)};
  \draw[->, red!70!black] (4.7,2.0) -- (4.7,1.75);
  \fill (0.5,0.7) circle (2pt) node[below right] {\(q(t_0)\)};
  \fill (6.5,2.35) circle (2pt) node[above left] {\(q(t_1)\)};
  \node[below] at (0.5,0) {\(t_0\)};
  \node[below] at (6.5,0) {\(t_1\)};
  \node[blue] at (2.3,0.6) {\(q(t)\)};
  \node[red] at (4.25,1.35) {\(q_\eps(t)=q(t)+\eps\eta(t)+O(\eps^2)\)};
\end{tikzpicture}
\caption{A fixed-endpoint variation changes the path in the interior while
holding the initial and final configurations fixed. The infinitesimal change is
the vector field \(\eta(t)=\delta q(t)\) along the path.}
\label{fig:fixed-endpoint-variation}
\end{figure}

\begin{derivation}[Euler-Lagrange equations]
The first variation of the action is
\[
\begin{aligned}
  \delta S
  &=
  \left.\frac{\dd}{\dd\eps}\right|_{\eps=0}
  \int_{t_0}^{t_1}L(q_\eps,\dot q_\eps,t)\,\dd t\\
  &=
  \int_{t_0}^{t_1}
  \left(
    \frac{\partial L}{\partial q^i}\eta^i
    +
    \frac{\partial L}{\partial \dot q^i}\dot\eta^i
  \right)\dd t.
\end{aligned}
\]
Integrate the second term by parts:
\[
  \int_{t_0}^{t_1}
  \frac{\partial L}{\partial \dot q^i}\dot\eta^i\,\dd t
  =
  \left[
    \frac{\partial L}{\partial \dot q^i}\eta^i
  \right]_{t_0}^{t_1}
  -
  \int_{t_0}^{t_1}
  \frac{\dd}{\dd t}\left(\frac{\partial L}{\partial \dot q^i}\right)\eta^i
  \,\dd t.
\]
The boundary term vanishes because \(\eta(t_0)=\eta(t_1)=0\). Therefore
\[
  \delta S
  =
  \int_{t_0}^{t_1}
  \left[
    \frac{\partial L}{\partial q^i}
    -
    \frac{\dd}{\dd t}\left(\frac{\partial L}{\partial \dot q^i}\right)
  \right]\eta^i\,\dd t.
\]
The fundamental lemma of the calculus of variations says that if this integral
vanishes for all compactly supported variations \(\eta\), then the coefficient
of every \(\eta^i\) vanishes. Hence
\[
  \boxed{
  \frac{\dd}{\dd t}\left(\frac{\partial L}{\partial \dot q^i}\right)
  -
  \frac{\partial L}{\partial q^i}=0
  }.
\]
\end{derivation}

The boxed equation is local in time because the variations are arbitrary in the
interior of the interval. This is the conceptual content of the fundamental
lemma: a single global statement, \(\delta S=0\), becomes a pointwise
differential equation only because the virtual displacement \(\eta(t)\) can be
chosen independently near each time. Boundary conditions and endpoint terms are
therefore not afterthoughts; they decide which variations are admissible and
which physical quantities appear at the endpoints.

\section{Boundary Terms And Forces}
\label{sec:boundary-forces}

The integration-by-parts calculation also explains momenta. Without imposing
fixed endpoints, the boundary contribution is
\[
  \left[p_i\eta^i\right]_{t_0}^{t_1},
  \qquad
  p_i=\frac{\partial L}{\partial \dot q^i}.
\]
Thus \(p_i\) is the coefficient of endpoint displacement in the first variation
of the action. This is a sharper definition than ``mass times velocity'', which
is only true in special coordinates for special Lagrangians.

If nonconservative generalized forces \(Q_i(q,\dot q,t)\) are present, the
Lagrange-d'Alembert principle is
\[
  \delta S+\int_{t_0}^{t_1}Q_i\eta^i\,\dd t=0,
\]
which gives
\[
  \frac{\dd}{\dd t}\left(\frac{\partial L}{\partial \dot q^i}\right)
  -
  \frac{\partial L}{\partial q^i}
  =
  Q_i.
\]
The present course focuses mainly on conservative and Hamiltonian systems, but
this formula is useful when comparing with fluids and continua.

\section{Natural Lagrangians And The Metric Form}
\label{sec:natural-lagrangians}

Many mechanical systems have a Lagrangian of the form
\[
  L(q,\dot q)=T(q,\dot q)-V(q)
  =
  \frac{1}{2}g_{ij}(q)\dot q^i\dot q^j - V(q),
\]
where \(g_{ij}\) is a positive-definite mass metric on \(Q\), and \(V\) is
potential energy.

\begin{derivation}[Geodesic-plus-force equation]
For the natural Lagrangian,
\[
  \frac{\partial L}{\partial \dot q^i}=g_{ij}\dot q^j,
\]
so
\[
  \frac{\dd}{\dd t}\frac{\partial L}{\partial \dot q^i}
  =
  g_{ij}\ddot q^j
  +
  \frac{\partial g_{ij}}{\partial q^k}\dot q^k\dot q^j.
\]
Also
\[
  \frac{\partial L}{\partial q^i}
  =
  \frac{1}{2}\frac{\partial g_{jk}}{\partial q^i}\dot q^j\dot q^k
  -
  \frac{\partial V}{\partial q^i}.
\]
The Euler-Lagrange equation becomes
\[
  g_{ij}\ddot q^j
  +
  \left(
    \frac{\partial g_{ij}}{\partial q^k}
    -
    \frac{1}{2}\frac{\partial g_{jk}}{\partial q^i}
  \right)\dot q^k\dot q^j
  =
  -\frac{\partial V}{\partial q^i}.
\]
After symmetrizing the velocity terms and raising an index with \(g^{\ell i}\),
\[
  \boxed{
  \ddot q^\ell+\Gamma^\ell_{jk}\dot q^j\dot q^k
  =
  -g^{\ell i}\frac{\partial V}{\partial q^i}
  },
\]
where
\[
  \Gamma^\ell_{jk}
  =
  \frac{1}{2}g^{\ell i}
  \left(
    \frac{\partial g_{ij}}{\partial q^k}
    +
    \frac{\partial g_{ik}}{\partial q^j}
    -
    \frac{\partial g_{jk}}{\partial q^i}
  \right)
\]
are the Christoffel symbols of the mass metric.
\end{derivation}

Thus even ``force-free'' motion on a curved configuration space need not look
like \(\ddot q=0\) in coordinates. The Christoffel terms are inertial terms
created by the geometry of \(Q\) and the chosen coordinates.

\section{Holonomic Constraints}
\label{sec:holonomic-constraints}

A holonomic constraint is a condition
\[
  \Phi^a(q,t)=0,\qquad a=1,\ldots,k,
\]
that cuts out an admissible submanifold of configuration space. If one works in
the ambient coordinates \(q^i\), the constrained equations take the multiplier
form
\[
  \frac{\dd}{\dd t}\frac{\partial L}{\partial \dot q^i}
  -
  \frac{\partial L}{\partial q^i}
  =
  \lambda_a\frac{\partial \Phi^a}{\partial q^i}.
\]
The multipliers \(\lambda_a\) are not arbitrary forces. They are whatever
constraint reactions are needed to keep the motion on \(\Phi^a=0\). If instead
one parametrizes the constraint surface directly, the \(\lambda_a\) disappear
and the same physics is expressed with fewer coordinates.

\section{Symmetry And Noether's Theorem}
\label{sec:noether}

Noether's theorem is the first place where the variational viewpoint pays a
large dividend. A symmetry is not merely a visual invariance of a picture; it
is a transformation of the configuration space that leaves the Lagrangian, and
therefore the action, unchanged. The infinitesimal version of that
transformation is a vector field \(\xi_Q\) along which the action has zero
first variation. The conserved quantity is the momentum paired with that vector
field.

\begin{definition}[Infinitesimal generator]
Let a Lie group \(G\) act on \(Q\). For \(\xi\in\mathfrak g\), the infinitesimal
generator is the vector field
\[
  \xi_Q(q)=\left.\frac{\dd}{\dd s}\right|_{s=0}\exp(s\xi)\cdot q.
\]
In coordinates, write \(\xi_Q=\xi_Q^i(q)\partial/\partial q^i\).
\end{definition}

\begin{assumption}[Lagrangian symmetry]
The action of \(G\) is a symmetry if
\[
  L(g\cdot q,\,T_qg\cdot \dot q,\,t)=L(q,\dot q,t)
\]
for all \(g\in G\), \(q\in Q\), and \(\dot q\in T_qQ\).
\end{assumption}

\begin{proposition}[Noether theorem]
For every \(\xi\in\mathfrak g\), the quantity
\[
  J_\xi(q,\dot q)
  =
  \frac{\partial L}{\partial \dot q^i}\xi_Q^i(q)
  =
  p_i\xi_Q^i(q)
\]
is conserved along solutions of the Euler-Lagrange equations.
\end{proposition}

\begin{proof}
Differentiate the symmetry condition along \(g=\exp(s\xi)\) at \(s=0\):
\[
  0=
  \frac{\partial L}{\partial q^i}\xi_Q^i
  +
  \frac{\partial L}{\partial \dot q^i}
  \frac{\dd}{\dd t}\xi_Q^i(q(t)).
\]
Along an Euler-Lagrange solution,
\[
  \frac{\partial L}{\partial q^i}
  =
  \frac{\dd}{\dd t}\frac{\partial L}{\partial \dot q^i}.
\]
Substitute this into the previous equation:
\[
  0=
  \frac{\dd}{\dd t}\left(\frac{\partial L}{\partial \dot q^i}\right)\xi_Q^i
  +
  \frac{\partial L}{\partial \dot q^i}
  \frac{\dd}{\dd t}\xi_Q^i
  =
  \frac{\dd}{\dd t}
  \left(
    \frac{\partial L}{\partial \dot q^i}\xi_Q^i
  \right).
\]
\end{proof}

\begin{example}[Translation and rotation]
For a particle in \(\R^3\) with \(L=\frac12m\norm{\dot x}^2-V(x)\), translation
symmetry in direction \(a\) gives \(J_a=m\dot x\cdot a\). If \(V\) is invariant
under all translations, linear momentum is conserved. Rotational symmetry about
axis \(a\) gives
\[
  J_a=(x\times m\dot x)\cdot a,
\]
so angular momentum is conserved when \(V\) is rotationally invariant.
\end{example}

\section{Basic Example: The Kepler Problem}
\label{sec:kepler-lagrangian}

The Kepler problem should appear immediately after the Lagrangian formalism,
because it is the simplest serious payoff of cyclic coordinates, Noether's
theorem, and reduction by symmetry.

\begin{assumption}[Kepler model]
We study a particle of reduced mass \(\mu\) moving in \(\R^3\setminus\{0\}\)
under an attractive inverse-square central potential. The Lagrangian is
\[
  L(r,\dot r)
  =
  \frac12\mu\norm{\dot r}^2+\frac{k}{\norm r},
  \qquad
  k=Gm_1m_2>0.
\]
The coordinate \(r=x_1-x_2\) is the relative separation of the two bodies.
\end{assumption}

Rotational invariance gives conservation of angular momentum
\[
  J=r\times p=\mu r\times\dot r.
\]
Since \(J\cdot r=0\) and \(J\cdot \dot r=0\), the motion stays in the plane
perpendicular to \(J\). In polar coordinates \((\rho,\theta)\) in that plane,
\[
  L=\frac12\mu(\dot\rho^2+\rho^2\dot\theta^2)+\frac{k}{\rho}.
\]
The angle \(\theta\) is cyclic, hence
\[
  \ell=\frac{\partial L}{\partial\dot\theta}=\mu\rho^2\dot\theta
\]
is conserved. The energy is
\[
  E=\frac12\mu\dot\rho^2+\frac{\ell^2}{2\mu\rho^2}-\frac{k}{\rho}.
\]
Thus the radial motion is one-dimensional motion in the effective potential
\[
  V_{\mathrm{eff}}(\rho)=\frac{\ell^2}{2\mu\rho^2}-\frac{k}{\rho}.
\]
In this subsection \(\ell\) is the ordinary, massful angular momentum of the
relative two-body problem. In Chapter~\ref{chap:perturbation-chaos}, where
asteroid motion is written per unit asteroid mass, the same orbit equation is
often written using the specific angular momentum
\(\ell_{\mathrm{spec}}=\ell/\mu=\rho^2\dot\theta\). The detailed central-force
derivation in Section~\ref{subsec:detailed-binet} keeps an arbitrary particle
mass \(m\), making the conversion between the two conventions explicit.

\begin{figure}[htbp]
\centering
\begin{tikzpicture}[scale=1.0, >=Latex]
  \draw[->] (0,0) -- (6.5,0) node[right] {\(\rho\)};
  \draw[->] (0,-1.7) -- (0,3.1) node[above] {\(V_{\mathrm{eff}}\)};
  \fill[blue!8] (0.624,-0.7) rectangle (2.52,0);
  \draw[blue, thick, domain=0.36:6.1, samples=160]
    plot (\x,{1.1/(\x*\x)-2.2/\x});
  \draw[dashed] (0,-0.7) -- (6.2,-0.7) node[right] {bound \(E<0\)};
  \draw[dashed] (0,0.45) -- (6.2,0.45) node[right] {scattering \(E>0\)};
  \node[blue] at (4.45,-1.22) {\(V_{\mathrm{eff}}\)};
  \fill (1.0,-1.1) circle (1.4pt) node[below] {minimum};
  \fill (0.624,-0.7) circle (1.3pt);
  \fill (2.52,-0.7) circle (1.3pt);
  \draw[<->, blue!70!black] (0.624,-1.42) -- (2.52,-1.42)
    node[midway, below] {allowed radial interval};
  \node[align=center] at (2.55,2.05) {centrifugal barrier\\near \(\rho=0\)};
\end{tikzpicture}
\caption{The Kepler problem reduces to one-dimensional radial motion in an
effective potential. For a fixed negative energy, the radial coordinate is
trapped between two turning points; for positive energy there is only a closest
approach before scattering.}
\label{fig:kepler-effective-potential}
\end{figure}

\begin{derivation}[Orbit equation]
Set \(u=1/\rho\). Since \(\ell=\mu\rho^2\dot\theta\),
\[
  \frac{\dd}{\dd t}=\dot\theta\frac{\dd}{\dd\theta}
  =
  \frac{\ell}{\mu}\,u^2\frac{\dd}{\dd\theta}.
\]
The radial Euler-Lagrange equation is
\[
  \mu(\ddot\rho-\rho\dot\theta^2)=-\frac{k}{\rho^2}.
\]
In terms of \(u(\theta)\) this becomes
\[
  \frac{\dd^2u}{\dd\theta^2}+u=\frac{\mu k}{\ell^2}.
\]
Therefore
\[
  u(\theta)=\frac{\mu k}{\ell^2}\left[1+e\cos(\theta-\theta_0)\right],
\]
or
\[
  \boxed{
  \rho(\theta)=
  \frac{\ell^2/(\mu k)}{1+e\cos(\theta-\theta_0)}
  }.
\]
The parameter \(e\) is the eccentricity. The conic is an ellipse, parabola, or
hyperbola according as \(e<1\), \(e=1\), or \(e>1\).
\end{derivation}

The Kepler problem also has the Laplace-Runge-Lenz vector
\[
  A=p\times J-\mu k\frac{r}{\norm r},
\]
which is conserved and points toward periapse. Its existence is special: it is
the extra conservation law behind the absence of perihelion precession in the
ideal inverse-square problem.

\section{Basic Example: Rigid Body As A Lagrangian System}
\label{sec:rigid-body-lagrangian}

The second basic example is the rigid body. It is best introduced directly after
Lagrangian mechanics, before using it later as a reduced Hamiltonian and
integrable system.

\begin{assumption}[Free rigid body]
A material point \(X\) in the body frame is located at
\[
  x(t)=R(t)X,
  \qquad R(t)\in\SO(3),
\]
with the center of mass fixed at the origin. The body is rigid, so all dynamics
is in the attitude \(R(t)\). No external torque acts.
\end{assumption}

The body angular velocity is defined by
\[
  R^{-1}\dot R=\widehat\Omega,
  \qquad
  \widehat\Omega v=\Omega\times v.
\]
The kinetic energy is
\[
  L(R,\dot R)=T
  =
  \frac12\Omega^TI\Omega,
\]
where \(I\) is the inertia tensor in the body frame. In principal axes,
\[
  I=\diag(I_1,I_2,I_3),
  \qquad
  L=\frac12(I_1\Omega_1^2+I_2\Omega_2^2+I_3\Omega_3^2).
\]

\begin{derivation}[Euler equation from constrained variations]
A variation inside \(\SO(3)\) has the form
\[
  \delta R=R\widehat\Sigma,
\]
where \(\Sigma(t_0)=\Sigma(t_1)=0\). This implies
\[
  \delta\Omega=\dot\Sigma+\Omega\times\Sigma.
\]
Let \(M=\partial L/\partial\Omega=I\Omega\). Then
\[
\begin{aligned}
  \delta S
  &=
  \int M\cdot(\dot\Sigma+\Omega\times\Sigma)\,\dd t\\
  &=
  \int(-\dot M+M\times\Omega)\cdot\Sigma\,\dd t.
\end{aligned}
\]
Stationarity gives
\[
  \dot M=M\times\Omega,
\]
or, with the hat convention \(\widehat\Omega v=\Omega\times v\) used
throughout,
\[
  \dot M+\Omega\times M=0.
\]
In principal axes this is
\[
\begin{aligned}
  I_1\dot\Omega_1 &= (I_2-I_3)\Omega_2\Omega_3,\\
  I_2\dot\Omega_2 &= (I_3-I_1)\Omega_3\Omega_1,\\
  I_3\dot\Omega_3 &= (I_1-I_2)\Omega_1\Omega_2.
\end{aligned}
\]
\end{derivation}

This derivation is deliberately placed here: the rigid body is not first of all
an abstract Lie-Poisson system. It is a Lagrangian system on \(\SO(3)\), and the
Lie-Poisson picture is what appears after symmetry reduction.

\section{Legendre Transform And Hamiltonian}
\label{sec:legendre-hamiltonian}

The Legendre transform sends \((q,\dot q)\in TQ\) to \((q,p)\in T^*Q\) by
\[
  p_i=\frac{\partial L}{\partial \dot q^i}.
\]
For a regular Lagrangian this can be inverted to express \(\dot q^i\) as a
function of \((q,p,t)\). Concretely, regularity means that the velocity Hessian
\[
  W_{ij}(q,\dot q,t)=
  \frac{\partial^2L}{\partial\dot q^i\partial\dot q^j}
\]
is nonsingular. The inverse function theorem then says that, locally in phase
space, the equations \(p_i=\partial L/\partial\dot q^i\) can be solved for
\(\dot q^i=\dot q^i(q,p,t)\).
The Hamiltonian is
\[
  H(q,p,t)=p_i\dot q^i(q,p,t)-L(q,\dot q(q,p,t),t).
\]

\begin{figure}[htbp]
\centering
\begin{tikzpicture}[scale=1.0, >=Latex]
  \node[draw, rounded corners=2pt, minimum width=2.9cm, minimum height=0.9cm,
    fill=blue!5] (tq) at (0,0) {\(TQ:\ (q,\dot q)\)};
  \node[draw, rounded corners=2pt, minimum width=2.9cm, minimum height=0.9cm,
    fill=red!5] (tsq) at (5.0,0) {\(T^*Q:\ (q,p)\)};
  \draw[->, thick] (tq) -- node[above]
    {\(\mathbb F L:\ p_i=\partial L/\partial\dot q^i\)} (tsq);
  \draw[->, thick, bend left=18] (tsq.north west) to node[above]
    {regular inverse} (tq.north east);
  \node[align=center, text width=0.82\textwidth] at (2.5,-1.25)
    {The Legendre map preserves \(q\) but replaces velocity by the endpoint
    covector \(p\).\\It is a coordinate change only when the velocity Hessian is
    nonsingular.};
\end{tikzpicture}
\caption{The Legendre transform from tangent-bundle data to cotangent-bundle
data. Singular Lagrangians require constraints; the Hamiltonian construction in
this chapter assumes the regular case.}
\label{fig:legendre-transform}
\end{figure}

\begin{derivation}[Hamilton's equations from the Legendre transform]
Differentiate \(H\):
\[
\begin{aligned}
  \dd H
  &=
  \dot q^i\,\dd p_i
  +
  p_i\,\dd \dot q^i
  -
  \frac{\partial L}{\partial q^i}\dd q^i
  -
  \frac{\partial L}{\partial \dot q^i}\dd \dot q^i
  -
  \frac{\partial L}{\partial t}\dd t\\
  &=
  \dot q^i\,\dd p_i
  -
  \frac{\partial L}{\partial q^i}\dd q^i
  -
  \frac{\partial L}{\partial t}\dd t,
\end{aligned}
\]
because \(p_i=\partial L/\partial \dot q^i\). Thus
\[
  \frac{\partial H}{\partial p_i}=\dot q^i,
  \qquad
  \frac{\partial H}{\partial q^i}=-\frac{\partial L}{\partial q^i}.
\]
Using the Euler-Lagrange equation \(\dot p_i=\partial L/\partial q^i\), we get
\[
  \boxed{
    \dot q^i=\frac{\partial H}{\partial p_i},
    \qquad
    \dot p_i=-\frac{\partial H}{\partial q^i}
  }.
\]
\end{derivation}

The Legendre transform changes the bookkeeping of initial data. A Lagrangian
state \((q,\dot q)\) records a position and a velocity; a Hamiltonian state
\((q,p)\) records a position and the covector that measures how the action
changes under endpoint displacement. The symplectic form below is built from
that endpoint one-form, so Hamiltonian mechanics remembers the variational
origin of the momentum rather than merely renaming velocity.

\section{Phase Space As The Space Of Solutions}
\label{sec:phase-space-solutions}

A useful slogan is that, for a deterministic Hamiltonian system, phase space is
the same thing as the space of solutions.
More precisely, if \(H\) is smooth and the Hamiltonian vector field is complete
on the time interval under study, then every point
\[
  z_0=(q_0,p_0)\in T^*Q
\]
determines a unique trajectory \(z(t)=\Phi^t_H(z_0)\), and every trajectory is
obtained this way from its value at any one time. Thus phase space is the set
of all possible instantaneous states, but it is also a coordinate system on the
family of all classical histories. The caveat is that this identification uses
existence and uniqueness of the equations of motion; singular potentials,
collisions, constraints, or finite-time blowup can invalidate the global
statement even when the local Hamiltonian equations make sense.

\section{The Symplectic Form}
\label{sec:symplectic-form}

The cotangent bundle \(T^*Q\) carries a canonical one-form
\[
  \Theta=p_i\,\dd q^i.
\]
The canonical symplectic form is
\[
  \omega=-\dd\Theta=\dd q^i\wedge\dd p_i.
\]
In arbitrary local phase-space coordinates
\[
  x^m,\qquad m=1,\ldots,2n,
\]
one can write
\[
  \omega=\frac12\omega_{mn}(x)\,\dd x^m\wedge\dd x^n,
\]
where \(\omega_{mn}=-\omega_{nm}\), \(\det(\omega_{mn})\neq0\), and
\[
  \dd\omega=0.
\]
The closedness condition is what makes \(\omega\) locally equivalent to the
canonical form, while nondegeneracy is what lets the Hamiltonian determine a
unique vector field. In canonical coordinates \(x=(q,p)\),
\[
  (\omega_{mn})=
  \begin{pmatrix}
  0 & I_n\\
  -I_n & 0
  \end{pmatrix}.
\]
The Poisson bracket can be expressed intrinsically by the bivector inverse to
the symplectic form, with the sign convention fixed by Hamilton's equations
below. The point is that the bracket is not extra structure once \(\omega\) is
given.

The two defining properties of \(\omega\) have different jobs. Nondegeneracy
lets one solve \(\iota_{X_H}\omega=\dd H\) for a unique vector field, just as a
mass matrix lets one solve force equals mass times acceleration. Closedness is
the integrability condition behind canonical coordinates, Stokes-type
arguments, and the preservation of symplectic area. Hamiltonian mechanics needs
both.

The Hamiltonian vector field \(X_H\) is defined by
\[
  \iota_{X_H}\omega=\dd H.
\]
Here \(\iota_X\) denotes interior contraction: if \(\omega\) is a two-form, then
\((\iota_X\omega)(Y)=\omega(X,Y)\) for every vector field \(Y\).
If
\[
  X_H=A^i\frac{\partial}{\partial q^i}
  +
  B_i\frac{\partial}{\partial p_i},
\]
then
\[
  \iota_{X_H}\omega=A^i\,\dd p_i-B_i\,\dd q^i.
\]
Equating this with
\[
  \dd H=\frac{\partial H}{\partial q^i}\dd q^i
  +
  \frac{\partial H}{\partial p_i}\dd p_i
\]
gives
\[
  A^i=\frac{\partial H}{\partial p_i},
  \qquad
  B_i=-\frac{\partial H}{\partial q^i}.
\]
Thus \(X_H\) is Hamilton's equation written geometrically.

\begin{figure}[htbp]
\centering
\begin{tikzpicture}[scale=1.0, >=Latex]
  \draw[->] (-0.25,0) -- (6.9,0) node[right] {\(q\)};
  \draw[->] (3.35,-2.35) -- (3.35,2.35) node[above] {\(p\)};
  \node[below] at (0.25,0) {\(-\pi\)};
  \node[below] at (3.35,0) {\(0\)};
  \node[below] at (6.45,0) {\(\pi\)};
  \draw[blue!70!black, thick] (3.35,0) ellipse (0.85 and 0.55);
  \draw[blue!70!black, thick] (3.35,0) ellipse (1.45 and 0.95);
  \draw[blue!70!black, thick] (3.35,0) ellipse (2.05 and 1.35);
  \draw[red!70!black, thick, domain=-3.12:3.12, samples=180]
    plot ({\x+3.35},{1.62*cos(0.5*\x r)});
  \draw[red!70!black, thick, domain=-3.12:3.12, samples=180]
    plot ({\x+3.35},{-1.62*cos(0.5*\x r)});
  \draw[gray!65, thick, domain=0.45:6.25, samples=160]
    plot (\x,{1.9+0.18*cos(120*\x)});
  \draw[gray!65, thick, domain=0.45:6.25, samples=160]
    plot (\x,{-1.9+0.18*cos(120*\x)});
  \draw[->, thick] (4.15,0.78) -- (4.45,0.62) node[right] {\(X_H\)};
  \node[blue!70!black] at (1.45,1.75) {libration levels};
  \node[red!70!black] at (5.15,-1.9) {separatrix};
  \node[gray!65!black] at (1.15,-1.9) {rotation levels};
  \node at (3.35,-2.75) {Pendulum level sets: \(H(q,p)=\frac12p^2+1-\cos q\).};
\end{tikzpicture}
\caption{A schematic phase plane using the pendulum Hamiltonian. The
Hamiltonian vector field is tangent to energy level sets because
\(\dot H=\{H,H\}=0\).}
\label{fig:pendulum-phase-plane}
\end{figure}

\section{Poisson Brackets And Time Evolution}
\label{sec:poisson-time-evolution}

For \(f,g\in C^\infty(T^*Q)\), define
\[
  \Poisson{f}{g}
  =
  \frac{\partial f}{\partial q^i}
  \frac{\partial g}{\partial p_i}
  -
  \frac{\partial f}{\partial p_i}
  \frac{\partial g}{\partial q^i}.
\]
If \(z(t)=(q(t),p(t))\) follows Hamilton's equation, then
\[
  \frac{\dd}{\dd t}f(q(t),p(t),t)
  =
  \frac{\partial f}{\partial t}+\Poisson{f}{H}.
\]
A time-independent observable \(f\) is conserved exactly when
\[
  \Poisson{f}{H}=0.
\]

\begin{proposition}[Liouville volume preservation]
Hamiltonian flow preserves the phase-space volume form
\[
  \Omega=\frac{1}{n!}\omega^n.
\]
\end{proposition}

\begin{proof}
Cartan's formula gives
\[
  \Lie_{X_H}\omega=\dd(\iota_{X_H}\omega)+\iota_{X_H}\dd\omega.
\]
Since \(\iota_{X_H}\omega=\dd H\) and \(\dd\omega=0\),
\[
  \Lie_{X_H}\omega=\dd\dd H=0.
\]
Therefore \(\Lie_{X_H}\omega^n=0\), so the associated volume form is preserved.
\end{proof}

\section{Phase-Space Action}
\label{sec:phase-space-action}

Hamilton's equations can also be derived from the phase-space action
\[
  S_H[q,p]=\int_{t_0}^{t_1}\left(p_i\dot q^i-H(q,p,t)\right)\dd t.
\]
Varying \(q\) and \(p\) independently, with \(\delta q(t_0)=\delta q(t_1)=0\),
gives
\[
  \delta S_H
  =
  \int_{t_0}^{t_1}
  \left[
    \left(\dot q^i-\frac{\partial H}{\partial p_i}\right)\delta p_i
    -
    \left(\dot p_i+\frac{\partial H}{\partial q^i}\right)\delta q^i
  \right]\dd t.
\]
Since \(\delta q\) and \(\delta p\) are arbitrary in the interior, the
stationary curves are exactly Hamiltonian trajectories.

\section{Canonical Transformations}
\label{sec:canonical-transformations}

A diffeomorphism \(\Phi:T^*Q\to T^*Q\) is canonical, or symplectic, if
\[
  \Phi^*\omega=\omega.
\]
Canonical transformations preserve Poisson brackets:
\[
  \Poisson{f\circ\Phi}{g\circ\Phi}
  =
  \Poisson{f}{g}\circ\Phi.
\]
This is why canonical coordinates may change but Hamiltonian mechanics keeps
the same form. Later, action-angle variables will be canonical coordinates
adapted to integrable motion.

\begin{derivation}[Generating function for a canonical transformation]
Let \((q,p)\) be old canonical coordinates and \((Q,P)\) new ones. A convenient
type of generating function depends on the new momenta and old positions:
\[
  F=F(P,q).
\]
Impose
\[
  Q_i=\frac{\partial F}{\partial P_i}\bigg|_q,
  \qquad
  p_i=\frac{\partial F}{\partial q^i}\bigg|_P.
\]
Then
\[
  \dd F=Q_i\,\dd P_i+p_i\,\dd q^i.
\]
Taking another exterior derivative gives
\[
  0=\dd Q_i\wedge\dd P_i+\dd p_i\wedge\dd q^i.
\]
Therefore
\[
  \dd Q_i\wedge\dd P_i
  =
  \dd q^i\wedge\dd p_i.
\]
The canonical symplectic form is preserved, so the transformation is canonical.
In coordinates this is the same cancellation one sees by expanding
\[
  \dd P_i\wedge\dd Q_i
  =
  \dd P_i\wedge
  \left(
    \frac{\partial^2F}{\partial P_i\partial q^j}\dd q^j
    +
    \frac{\partial^2F}{\partial P_i\partial P_j}\dd P_j
  \right),
\]
where the pure \(\dd P_i\wedge\dd P_j\) term vanishes because second
derivatives are symmetric while the wedge product is antisymmetric.
\end{derivation}

\section{Numerical Lesson}
\label{sec:numerical-lesson}

The equations in this chapter are exact statements about smooth systems. A
numerical integrator creates a discrete map
\[
  z_{n+1}=\Psi_{\Delta t}(z_n).
\]
For long-time Hamiltonian problems, whether \(\Psi_{\Delta t}\) preserves a
symplectic form or a modified Hamiltonian can matter more than the local order
of the method. This is why the pendulum, standard map, and asteroid demos are
not merely programming exercises; they are experiments in preserving structure.

The practical lesson is diagnostic. A good mechanics simulation reports not
only a trajectory but also the quantities that the continuous problem says
should be preserved or nearly preserved. Energy drift, angular-momentum drift,
phase-space area, and event criteria are part of the model audit trail. Without
them, a numerical plot can look convincing while solving the wrong mechanical
problem.

\section{What To Take Forward}
\label{sec:lag-ham-take-forward}

The rest of the notes repeatedly reuse four lessons from this chapter.
\begin{itemize}
\item The configuration space \(Q\) is part of the model. Choosing bad
coordinates can obscure a simple system, but it cannot change the system.
\item Boundary terms are information, not clutter. They identify canonical
momenta, forces, and conserved quantities.
\item Symmetry produces conserved quantities through Noether's theorem, and
those quantities are often the coordinates of a reduced problem.
\item Hamiltonian mechanics is geometry on \(T^*Q\): the symplectic form
determines Hamilton's equations, volume preservation, Poisson brackets, and the
meaning of a canonical numerical method.
\end{itemize}

\chapter{Integrability And The Rigid Body}
\label{chap:rigid-integrability}

\section{Why The Rigid Body Is Central}
\label{sec:rigid-why-central}

The free rigid body is a nearly ideal teaching example. It is nonlinear, it
lives naturally on the Lie group \(\SO(3)\), it has a reduced Hamiltonian
description on angular momentum space, and it is completely integrable. At the
same time, small modifications such as gravity acting on a heavy top lead toward
nonintegrability and chaos.

\begin{assumption}[Rigid body model]
A rigid body is an ideal continuum whose internal distances do not change. A
material point is labelled by a body coordinate \(X\in B\subset\R^3\). Its
spatial position is
\[
  x(t)=r(t)+R(t)X,
\]
where \(r(t)\in\R^3\) is the center-of-mass position and
\(R(t)\in\SO(3)\) is the attitude matrix. The free rigid body chapter suppresses
translation by working in the center-of-mass frame.
\end{assumption}

\begin{table}[htbp]
\centering
\caption{Rigid-body and integrability notation. Symbols that are overloaded in
other chapters are cross-listed in Section~\ref{sec:overloaded-symbols}.}
\label{tab:rigid-integrability-notation}
\small
\begin{tabular}{p{0.18\textwidth}p{0.48\textwidth}p{0.25\textwidth}}
\toprule
Symbol & Meaning & Space or units\\
\midrule
\(R\) & attitude matrix & \(\SO(3)\)\\
\(\hat e_a\) & \(a\)-th body axis seen in space & unit vector in \(\R^3\)\\
\(\omega_{\mathrm{space}}\) & spatial angular velocity & \(\R^3\), \(T^{-1}\)\\
\(\Omega\) & body angular velocity & \(\R^3\), \(T^{-1}\)\\
\(\widehat\Omega\) & skew matrix with \(\widehat\Omega v=\Omega\times v\) &
\(\mathfrak{so}(3)\)\\
\(I\) & inertia tensor in body frame & symmetric positive matrix\\
\(I_1,I_2,I_3\) & principal moments of inertia & positive scalars\\
\(M\) & body angular momentum & \(I\Omega\)\\
\(J\) & spatial angular momentum & \(J=RM\)\\
\(H\) & kinetic energy Hamiltonian & scalar energy\\
\(C\) & squared angular momentum & Casimir on \(\mathfrak{so}(3)^*\)\\
\(\ell(\Omega)\) & reduced Lagrangian & function on \(\mathfrak{so}(3)\)\\
\(F_i\) & commuting first integrals & Liouville theorem\\
\(M_f\) & common regular level set \(F_i=f_i\) & invariant manifold\\
\(\Lambda\) & period lattice of commuting flows & subgroup of \(\R^n\)\\
\(\gamma_j\) & one-cycle on an invariant torus & element of \(H_1(\mathbb T^n)\)\\
\(I_j\) & action variable & \((2\pi)^{-1}\oint_{\gamma_j}p_i\dd q^i\)\\
\(\theta^j\) & angle variable conjugate to \(I_j\) & \(2\pi\)-periodic coordinate\\
\(S(I,q)\) & generating function for angles & locally multivalued\\
\bottomrule
\end{tabular}
\end{table}

\paragraph{Roadmap.}
The chapter moves from a concrete body to the general geometry of integrable
systems. The first half explains the kinematics and dynamics of a free rigid
body: rotations live on \(\SO(3)\), angular velocity lives in its Lie algebra,
and the reduced motion lives on angular momentum spheres. The second half
abstracts the lesson: enough independent commuting first integrals force
regular compact level sets to be tori, and on those tori the motion is linear.
The rigid body is therefore not just an example after the theorem; it is the
model that makes the theorem visible.

\paragraph{Three layers to keep separate.}
Rigid-body mechanics mixes three descriptions that are often compressed into
one line of notation. The kinematic layer says how \(R(t)\in\SO(3)\) defines
angular velocity. The dynamical layer says how the kinetic energy and the
inertia tensor produce Euler's equations. The reduced Hamiltonian layer says
why angular momentum spheres carry a Lie-Poisson flow. Each layer has its own
variables and its own conserved quantities. Confusing them is the fastest way
to lose a sign or to mistake body components for spatial components.

\section{The Rotation Group And Angular Velocity}
\label{sec:rotation-group-angular-velocity}

The rotation group is
\[
  \SO(3)=\{R\in\R^{3\times 3}:R^TR=I,\ \det R=1\}.
\]
It is useful to think of \(R\) as the matrix whose columns are the three
orthonormal body axes \(\hat e_1,\hat e_2,\hat e_3\) as seen in space:
\[
  R=(\hat e_1\ \hat e_2\ \hat e_3).
\]
Thus a point of configuration space is not a triple of unconstrained vectors;
it is an orthonormal oriented frame. The condition \(R^TR=I\) gives six
constraints on nine matrix entries, leaving three degrees of freedom.

The spatial angular velocity \(\omega_{\mathrm{space}}\) is defined by the fact that every body
axis instantaneously rotates according to
\[
  \frac{\dd \hat e_a}{\dd t}=\omega_{\mathrm{space}}\times\hat e_a,
\]
or, in components,
\[
  \dot R_{ia}=\epsilon_{ijk}(\omega_{\mathrm{space}})_jR_{ka}.
\]

Differentiate \(R^TR=I\):
\[
  \dot R^TR+R^T\dot R=0.
\]
Thus \(R^T\dot R\) is skew-symmetric. Every skew-symmetric \(3\times3\) matrix
is uniquely of the form
\[
  \widehat\Omega=
  \begin{pmatrix}
  0 & -\Omega_3 & \Omega_2\\
  \Omega_3 & 0 & -\Omega_1\\
  -\Omega_2 & \Omega_1 & 0
  \end{pmatrix},
\]
where \(\widehat\Omega v=\Omega\times v\). The body angular velocity is defined
by
\[
  R^T\dot R=\widehat\Omega.
\]
The spatial angular velocity \(\omega_{\mathrm{space}}\) satisfies
\[
  \dot R R^T=\widehat{\omega}_{\mathrm{space}},
  \qquad
  \omega_{\mathrm{space}}=R\Omega.
\]

This space/body distinction is one of the main sources of sign errors in rigid
body mechanics. The vector \(\omega_{\mathrm{space}}\) is expressed in fixed
space axes, while \(\Omega\) is expressed in the moving body axes. They describe
the same instantaneous rotation, but derivatives of body components must also
differentiate the moving frame. Euler's equation is exactly where this
bookkeeping becomes dynamical.

\section{\texorpdfstring{Euler Angles, Metric, And The Topology Of \(\SO(3)\)}{Euler Angles, Metric, And The Topology Of SO(3)}}
\label{sec:euler-angles-so3}

One can parametrize a dense open set of \(\SO(3)\) by Euler angles
\[
  (\phi,\theta,\psi),
  \qquad
  0\leq \phi<2\pi,\quad 0\leq \theta\leq \pi,\quad 0\leq \psi<2\pi.
\]
For the Euler-angle convention used below, the body axes are
\[
\begin{aligned}
\hat e_1 &=
\begin{pmatrix}
\cos\phi\cos\psi-\sin\phi\cos\theta\sin\psi\\
\sin\phi\cos\psi+\cos\phi\cos\theta\sin\psi\\
\sin\theta\sin\psi
\end{pmatrix},\\
\hat e_2 &=
\begin{pmatrix}
-\cos\phi\sin\psi-\sin\phi\cos\theta\cos\psi\\
-\sin\phi\sin\psi+\cos\phi\cos\theta\cos\psi\\
\sin\theta\cos\psi
\end{pmatrix},\\
\hat e_3 &=
\begin{pmatrix}
\sin\phi\sin\theta\\
-\cos\phi\sin\theta\\
\cos\theta
\end{pmatrix}.
\end{aligned}
\]
The exact convention is less important than the lesson: Euler angles are a
coordinate chart, not the space itself. The chart is singular at
\(\theta=0,\pi\), even though \(\SO(3)\) is a smooth homogeneous space.

A coordinate-independent line element on \(\SO(3)\) is
\[
  \dd s^2=-\frac12\tr\left(R^{-1}\dd R\right)^2.
\]
For this standard Euler-angle convention, this becomes
\[
  \dd s^2
  =
  \dd\phi^2+\dd\theta^2+\dd\psi^2
  +2\cos\theta\,\dd\phi\,\dd\psi.
\]
Introduce
\[
  \xi_1=\phi+\psi,\qquad \xi_2=\phi-\psi.
\]
Then
\[
  \dd s^2
  =
  \dd\theta^2
  +
  \cos^2\frac{\theta}{2}\,\dd\xi_1^2
  +
  \sin^2\frac{\theta}{2}\,\dd\xi_2^2.
\]
This formula shows the topology pictorially. At one endpoint of the
\(\theta\)-interval the \(\xi_2\)-circle shrinks; at the other endpoint the
\(\xi_1\)-circle shrinks. Gluing these solid tori gives \(S^3\), and the
remaining identification is the antipodal identification. Hence
\[
  \SO(3)\simeq S^3/\Z_2.
\]
This is why a rigid body admits local angle coordinates but no single global
nondegenerate Euler-angle chart.

\section{Kinetic Energy And The Inertia Tensor}
\label{sec:inertia-tensor}

Take the center of mass as origin in body coordinates:
\[
  \int_B X\,\rho(X)\,\dd^3X=0,
\]
where \(\rho(X)\) is the mass density. A point \(X\) has velocity
\[
  \dot x=\dot R X=R\widehat\Omega X=R(\Omega\times X).
\]
Since \(R\) preserves Euclidean lengths,
\[
  \norm{\dot x}^2=\norm{\Omega\times X}^2.
\]
The kinetic energy is
\[
  T=\frac12\int_B\rho(X)\norm{\Omega\times X}^2\,\dd^3X.
\]
Using
\[
  \norm{\Omega\times X}^2
  =
  \norm{\Omega}^2\norm{X}^2-(\Omega\cdot X)^2,
\]
we obtain
\[
  T=\frac12\Omega^T I\Omega,
\]
where
\[
  I_{ij}
  =
  \int_B\rho(X)
  \left(\norm{X}^2\delta_{ij}-X_iX_j\right)\dd^3X.
\]
The matrix \(I\) is the inertia tensor in the body frame. In principal axes,
\[
  I=\diag(I_1,I_2,I_3),
\]
and
\[
  T=\frac12(I_1\Omega_1^2+I_2\Omega_2^2+I_3\Omega_3^2).
\]
Equivalently, using the spatial angular velocity \(\omega_{\mathrm{space}}\) and the body axes
\(\hat e_a\),
\[
  L=\frac12\sum_{a=1}^3 I_a(\hat e_a\cdot\omega_{\mathrm{space}})^2.
\]
The spatial angular momentum is
\[
  J_i=\frac{\partial L}{\partial (\omega_{\mathrm{space}})_i}
  =
  \sum_{a=1}^3 I_a(\hat e_a)_i(\hat e_a\cdot\omega_{\mathrm{space}}).
\]
Its body components are
\[
  \widetilde J_a=\hat e_a\cdot J=I_a\Omega_a,
\]
and the Hamiltonian is
\[
  H=\sum_{a=1}^3\frac{\widetilde J_a^2}{2I_a}.
\]

\section{\texorpdfstring{Full Phase Space: \(T^*\SO(3)\)}{Full Phase Space: T*SO(3)}}
\label{sec:rigid-full-phase-space}

The Hamiltonian phase space of a finite-dimensional mechanical system with
configuration space \(Q\) is the cotangent bundle \(T^*Q\). Therefore the
unreduced phase space of a rigid body with fixed center of mass is
\[
  T^*\SO(3).
\]
If \(q^i\), \(i=1,2,3\), are any nondegenerate local coordinates on
\(\SO(3)\), the canonical momenta are
\[
  p_i=\frac{\partial L}{\partial \dot q^i}.
\]
Under a change of coordinates \(q^i=q^i(\tilde q)\), the momenta transform as a
one-form:
\[
  \tilde p_j
  =
  \frac{\partial q^i}{\partial \tilde q^j}p_i.
\]
This is the coordinate calculation behind the intrinsic statement
\((q,p)\in T^*Q\).

\begin{remark}[Angular momentum is not a canonical momentum]
Let \(M_a\) be the components of body angular momentum in principal axes:
\[
  M_a=I_a\Omega_a.
\]
The \(M_a\) are excellent coordinates on the fiber of \(T^*\SO(3)\) after a
choice of body frame, but they are not canonical momenta conjugate to Euler
angles. Locally one has a nondegenerate linear relation
\[
  p_i=A_{ia}(q)M_a,
\]
where the matrix \(A(q)\) depends on the chosen angle coordinates. The canonical
symplectic form is still
\[
  \omega=\dd q^i\wedge \dd p_i,
\]
but when it is expressed in the noncanonical variables \(M_a\), the Poisson
brackets become the rigid-body Lie-Poisson brackets rather than
\(\Poisson{M_a}{M_b}=0\).
Globally, because \(\SO(3)\) is a Lie group, left or right translation
trivializes the cotangent bundle:
\[
  T^*\SO(3)\simeq \SO(3)\times\mathfrak{so}(3)^*.
\]
This global trivialization is the precise sense in which the angular-momentum
components give global fiber coordinates, even though they are noncanonical.
\end{remark}

\section{\texorpdfstring{Variations On \(\SO(3)\)}{Variations On SO(3)}}
\label{sec:variations-so3}

A variation of \(R(t)\in\SO(3)\) must remain tangent to the group. Write
\[
  \delta R=R\widehat\Sigma,
\]
where \(\Sigma(t)\in\R^3\) vanishes at the time endpoints. The induced variation
of body angular velocity is
\[
  \delta\widehat\Omega
  =
  \delta(R^T\dot R)
  =
  \dot{\widehat\Sigma}+[\widehat\Omega,\widehat\Sigma].
\]
Under the hat isomorphism this becomes
\[
  \delta\Omega=\dot\Sigma+\Omega\times\Sigma.
\]

\begin{derivation}[Euler-Poincare equation for the free rigid body]
The reduced action is
\[
  S[R]=\int_{t_0}^{t_1}\ell(\Omega)\,\dd t,
  \qquad
  \ell(\Omega)=\frac12\Omega^TI\Omega.
\]
Let
\[
  M=\frac{\partial \ell}{\partial \Omega}=I\Omega.
\]
Then
\[
\begin{aligned}
  \delta S
  &=
  \int_{t_0}^{t_1} M\cdot(\dot\Sigma+\Omega\times\Sigma)\,\dd t\\
  &=
  \left[M\cdot\Sigma\right]_{t_0}^{t_1}
  -
  \int_{t_0}^{t_1}\dot M\cdot\Sigma\,\dd t
  +
  \int_{t_0}^{t_1}M\cdot(\Omega\times\Sigma)\,\dd t.
\end{aligned}
\]
The boundary term vanishes. Since
\[
  M\cdot(\Omega\times\Sigma)=(M\times\Omega)\cdot\Sigma,
\]
stationarity for arbitrary \(\Sigma\) gives
\[
  -\dot M+M\times\Omega=0.
\]
Equivalently,
\[
  \boxed{\dot M+\Omega\times M=0}
\]
with the hat convention \(\widehat\Omega v=\Omega\times v\) used throughout
these notes. In components this is
\[
  I_1\dot\Omega_1=(I_2-I_3)\Omega_2\Omega_3,
\]
and cyclic permutations.
\end{derivation}

\section{Euler Equations In Principal Axes}
\label{sec:euler-equations-principal}

With \(M_i=I_i\Omega_i\), Euler's equations are
\[
\begin{aligned}
  I_1\dot\Omega_1 &= (I_2-I_3)\Omega_2\Omega_3,\\
  I_2\dot\Omega_2 &= (I_3-I_1)\Omega_3\Omega_1,\\
  I_3\dot\Omega_3 &= (I_1-I_2)\Omega_1\Omega_2.
\end{aligned}
\]
These equations are autonomous and quadratic. They conserve the kinetic energy
\[
  H=\frac12(I_1\Omega_1^2+I_2\Omega_2^2+I_3\Omega_3^2)
  =
  \frac12\left(\frac{M_1^2}{I_1}+\frac{M_2^2}{I_2}+\frac{M_3^2}{I_3}\right)
\]
and the squared angular momentum
\[
  C=\norm{M}^2=M_1^2+M_2^2+M_3^2.
\]

\begin{proof}[Direct conservation check]
Since \(M=I\Omega\),
\[
  \dot H=\Omega\cdot \dot M.
\]
Using \(\dot M=-\Omega\times M\),
\[
  \dot H=-\Omega\cdot(\Omega\times M)=0.
\]
Similarly,
\[
  \dot C=2M\cdot\dot M=-2M\cdot(\Omega\times M)=0.
\]
\end{proof}

\section{Geometry Of The Reduced Motion}
\label{sec:rigid-reduced-geometry}

The reduced motion in \(M\)-space lies on the intersection of:
\[
  C(M)=\norm{M}^2=\text{constant}
\]
and
\[
  H(M)=\frac12M^TI^{-1}M=\text{constant}.
\]
Thus it is the intersection of an angular momentum sphere and an energy
ellipsoid.

\begin{figure}[htbp]
\centering
\begin{tikzpicture}[
    scale=1.05,
    >=Latex,
    x={(1cm,0cm)},
    y={(0.42cm,0.27cm)},
    z={(0cm,1cm)},
    line join=round
  ]
  \draw[->, gray!65] (0,0,0) -- (2.75,0,0) node[right] {\(M_1\)};
  \draw[->, gray!65] (0,0,0) -- (0,2.45,0) node[above right] {\(M_2\)};
  \draw[->, gray!65] (0,0,0) -- (0,0,3.15) node[above] {\(M_3\)};

  \draw[red!65!black, thick, domain=0:360, samples=120, variable=\t]
    plot ({2*cos(\t)}, {2*sin(\t)}, {0});
  \draw[red!65!black, dashed, domain=0:360, samples=120, variable=\t]
    plot ({2*cos(\t)}, {0}, {2*sin(\t)});
  \draw[red!65!black, dashed, domain=0:360, samples=120, variable=\t]
    plot ({0}, {2*cos(\t)}, {2*sin(\t)});

  \draw[blue!70!black, thick, domain=0:360, samples=140, variable=\t]
    plot ({sqrt(847/340)*cos(\t)}, {sqrt(2541/507)*sin(\t)}, {0});
  \draw[blue!70!black, thin, domain=0:360, samples=140, variable=\t]
    plot ({sqrt(847/340)*cos(\t)}, {0}, {2.75*sin(\t)});
  \draw[blue!70!black, thin, domain=0:360, samples=140, variable=\t]
    plot ({0}, {sqrt(2541/507)*cos(\t)}, {2.75*sin(\t)});

  \draw[purple!80!black, very thick, domain=0:360, samples=180, variable=\t]
    plot ({sqrt(1 + 0.75*sin(\t)*sin(\t))}, {1.732*cos(\t)}, {1.5*sin(\t)});
  \draw[purple!80!black, very thick, dashed, domain=0:360, samples=180, variable=\t]
    plot ({-sqrt(1 + 0.75*sin(\t)*sin(\t))}, {1.732*cos(\t)}, {1.5*sin(\t)});

  \fill[purple!80!black] ({sqrt(1 + 0.75*0*0)}, {1.732}, {0}) circle (1.3pt);
  \node[red!65!black, align=center] at (-2.2,-1.15,-0.65)
    {angular momentum\\sphere \(C=\norm M^2\)};
  \node[blue!70!black, align=center] at (1.95,1.55,2.95)
    {energy ellipsoid\\\(2H=M^TI^{-1}M\)};
  \node[purple!80!black] at (-0.55,-2.2,1.45)
    {intersection trajectories};
\end{tikzpicture}
\caption{Rigid-body reduction in angular-momentum space. Conservation of
\(\norm M^2\) confines the motion to a sphere, conservation of \(H\) confines it
to an ellipsoid, and the reduced trajectory is their intersection. The solid
and dashed purple curves are the two projected branches of the exact
three-dimensional intersection for the sphere \(\norm M^2=4\) and the
representative ellipsoid with semi-axes
\(\sqrt{847/340},\sqrt{2541/507},11/4\).}
\label{fig:rigid-body-reduced-geometry}
\end{figure}

\section{Stability Of Principal-Axis Rotations}
\label{sec:principal-axis-stability}

Assume
\[
  I_1<I_2<I_3.
\]
A steady rotation about a principal axis is a solution. For example,
\[
  \Omega(t)=(\Omega_0,0,0)
\]
is rotation about the first axis.

\begin{derivation}[Linear stability]
Perturb rotation about the first axis:
\[
  \Omega_1=\Omega_0+\alpha_1,\qquad
  \Omega_2=\alpha_2,\qquad
  \Omega_3=\alpha_3,
\]
and keep only linear terms. Then
\[
  \dot\alpha_1=0,
\]
while
\[
  I_2\dot\alpha_2=(I_3-I_1)\Omega_0\alpha_3,
  \qquad
  I_3\dot\alpha_3=(I_1-I_2)\Omega_0\alpha_2.
\]
Differentiate the first of these:
\[
  \ddot\alpha_2
  =
  \frac{(I_3-I_1)(I_1-I_2)}{I_2I_3}\Omega_0^2\alpha_2.
\]
Since \(I_3-I_1>0\) and \(I_1-I_2<0\), the coefficient is negative, so
\(\alpha_2\) oscillates. Rotation about the first axis is linearly stable.

For rotation about the second axis, the analogous coefficient is
\[
  \frac{(I_2-I_3)(I_1-I_2)}{I_1I_3}\Omega_0^2>0,
\]
because both factors in the numerator are negative. The perturbation therefore
has exponential solutions. Thus the intermediate axis is unstable. Rotation
about the third axis is stable.
\end{derivation}

The geometric reason is that the energy restricted to a fixed angular momentum
sphere has extrema at the smallest and largest inertia axes but a saddle at the
intermediate axis.

\section{Lie-Poisson Structure}
\label{sec:lie-poisson-rigid-body}

The reduced rigid body is Hamiltonian on \(\mathfrak{so}(3)^*\cong\R^3\), but
not with canonical coordinates. The Lie-Poisson bracket is
\[
  \Poisson{F}{G}(M)
  =
  -M\cdot\left(\nabla F(M)\times\nabla G(M)\right).
\]
For
\[
  H(M)=\frac12M\cdot I^{-1}M,
  \qquad
  \nabla H=I^{-1}M=\Omega,
\]
the equation for any observable \(F\) is
\[
  \dot F=\Poisson{F}{H}.
\]
Taking \(F(M)=M_i\) yields
\[
  \dot M=-\Omega\times M.
\]

\begin{definition}[Casimir]
A Casimir is a function \(C\) satisfying
\[
  \Poisson{C}{F}=0
\]
for every observable \(F\). For the rigid body bracket,
\[
  C(M)=\norm{M}^2
\]
is a Casimir. Its level sets are the coadjoint orbits, which here are spheres.
\end{definition}

\section{Liouville Integrability And Invariant Tori}
\label{sec:liouville-tori}

The rigid-body calculations above show a pattern: conservation laws confine
motion to lower-dimensional sets, and the reduced flow on those sets can be
simple even when the original coordinates look complicated. Liouville
integrability is the precise Hamiltonian version of this pattern. It is not
just ``having many constants of motion.'' The constants must be independent,
they must Poisson commute, and their common level set must be regular enough to
carry the commuting Hamiltonian flows.

The three hypotheses have separate meanings. Independence says the level set
has the expected dimension. Involution says the Hamiltonian flows generated by
the integrals are compatible rather than fighting one another. Compactness says
the flow coordinates eventually repeat, turning straight-line motion on
\(\R^n\) into straight-line motion on a torus. Dropping any one of these
assumptions changes the conclusion.

\begin{definition}[Liouville integrability]
A Hamiltonian system on a \(2n\)-dimensional symplectic manifold is Liouville
integrable if there are \(n\) independent functions
\[
  F_1=H,\ F_2,\ldots,F_n
\]
such that
\[
  \Poisson{F_i}{F_j}=0
\]
for all \(i,j\). The functions are said to be in involution.
\end{definition}

\begin{example}[Basic integrable systems]
Any autonomous Hamiltonian system with one degree of freedom is integrable: the
Hamiltonian itself supplies the one conserved quantity. The fixed-COM rigid body
has three degrees of freedom. The three spatial functions
\[
  H,\qquad J_3,\qquad J^2
\]
are independent and mutually Poisson commuting on regular sets, even though the
individual components \(J_1,J_2,J_3\) do not commute with one another. This is
the rigid-body example used throughout this chapter.
\end{example}

\begin{assumption}[Regular compact level set]
Fix constants \(f=(f_1,\ldots,f_n)\) and define
\[
  M_f=\{x\in P: F_i(x)=f_i,\ i=1,\ldots,n\},
\]
where \(P\) is the \(2n\)-dimensional phase space. We assume that the
differentials \(\dd F_i\) are independent on \(M_f\) and that \(M_f\) is
compact and connected.
\end{assumption}

The regularity assumption says that \(M_f\) is a smooth \(n\)-dimensional
submanifold. Its tangent vectors \(v\in T_xM_f\) are exactly those satisfying
\[
  \dd F_i(v)=0
  \qquad\text{for all }i.
\]
Let \(X_{F_i}\) be the Hamiltonian vector field of \(F_i\), defined by
\[
  \iota_{X_{F_i}}\omega=\dd F_i.
\]
Then
\[
  \dd F_j(X_{F_i})
  =
  \Poisson{F_j}{F_i}
  =
  0.
\]
Thus every \(X_{F_i}\) is tangent to \(M_f\). The commutator of two Hamiltonian
vector fields is
\[
  [X_{F_i},X_{F_j}]
  =
  -X_{\Poisson{F_i}{F_j}},
\]
so involution implies
\[
  [X_{F_i},X_{F_j}]=0.
\]
We therefore get \(n\) independent commuting vector fields on \(M_f\).

\begin{derivation}[Why the regular compact level set is a torus]
The proof has a simple geometric shape. The Hamiltonian vector fields
\(X_{F_i}\) give \(n\) independent directions tangent to \(M_f\), and because
the integrals Poisson commute, flowing in one direction and then another gives
the same point as flowing in the opposite order. Thus \(M_f\) is locally
coordinatized by ``flow times.'' Compactness is what turns these flow-time
coordinates from \(\R^n\) into periodic coordinates.

Commuting vector fields have commuting flows. Starting at a point \(x_0\in
M_f\), flow for times \(s=(s^1,\ldots,s^n)\) along the \(n\) vector fields:
\[
  \Phi_s(x_0)
  =
  \exp(s^1X_{F_1})\cdots\exp(s^nX_{F_n})x_0.
\]
Because the flows commute, this defines an action of \(\R^n\) on \(M_f\).
Independence of the vector fields makes the action locally free, so the
coordinates \(s^i\) are local coordinates on \(M_f\). The only possible global
identifications are translations
\[
  s\sim s+\lambda,\qquad \lambda\in\Lambda,
\]
where \(\Lambda\subset\R^n\) is the period lattice of the orbit. Compactness
forces \(\Lambda\) to have rank \(n\): if the rank were smaller, the quotient
\(\R^n/\Lambda\) would still contain at least one noncompact Euclidean
direction. The action is locally free, so each orbit is open in the
\(n\)-manifold \(M_f\). Since \(M_f\) is connected, that open orbit is the
whole level set. Hence
\[
  M_f\simeq \R^n/\Lambda\simeq \mathbb T^n.
\]
These are the invariant Liouville tori.
\end{derivation}

This argument is often the first place where the word ``torus'' becomes more
than a picture. The torus is not assumed; it is forced by the existence of
enough commuting flows on a compact regular level set. The angle variables
introduced below are coordinates along these flows, and the action variables
are the quantities that label which torus one is on.

\begin{figure}[htbp]
\centering
\begin{tikzpicture}[scale=1.0, >=Latex]
  \draw[thick] (0,0) rectangle (4.2,3.0);
  \draw[->, blue!70!black, thick] (0.45,0.55) -- (3.75,0.55)
    node[midway, below] {\(\gamma_1\)};
  \draw[->, red!75!black, thick] (0.55,0.45) -- (0.55,2.55)
    node[midway, left] {\(\gamma_2\)};
  \draw[purple!75!black, thick, ->]
    (0.25,0.25) -- (1.1,0.9) -- (1.95,1.55) -- (2.8,2.2) -- (3.65,2.85);
  \draw[purple!75!black, thick, ->]
    (3.65,2.85) -- (4.15,3.23);
  \draw[purple!75!black, thick, ->]
    (-0.15,-0.23) -- (0.25,0.25);
  \draw[gray!70, ->] (4.2,1.5) -- (4.8,1.5)
    node[right] {identify};
  \draw[gray!70, ->] (2.1,3.0) -- (2.1,3.55)
    node[above] {identify};
  \node[align=center] at (2.1,-0.75)
    {flat model of a two-dimensional invariant torus\\
     \(\mathbb R^2/\Lambda\) with angle coordinates};
  \node[purple!75!black, align=center] at (5.25,2.45)
    {Hamiltonian flow\\straightens to\\constant slope};
\end{tikzpicture}
\caption{A Liouville torus is locally a space of commuting flow times with
periodic identifications. Action variables label the torus; angle variables
are coordinates along the fundamental cycles.}
\label{fig:liouville-flat-torus}
\end{figure}

The symplectic form vanishes when restricted to an invariant torus. Indeed, on
the tangent basis \(X_{F_i}\),
\[
  \omega(X_{F_i},X_{F_j})
  =
  \dd F_i(X_{F_j})
  =
  \Poisson{F_i}{F_j}
  =
  0.
\]
Since these vectors span \(T_xM_f\), the torus is Lagrangian:
\[
  \omega|_{M_f}=0.
\]

\begin{derivation}[Coordinate proof that the invariant torus is Lagrangian]
A useful coordinate check is the following. Use canonical
coordinates \((q_i,p_i)\), and suppose the level set is written locally as
\[
  F_i(p,q)=f_i,\qquad i=1,\ldots,n.
\]
If the equations can be solved locally for
\[
  p_i=p_i(q,F),
\]
then
\[
\begin{aligned}
  \omega
  &=
  \dd q_i\wedge\dd p_i\\
  &=
  \dd q_i\wedge
  \left(
    \frac{\partial p_i}{\partial q_j}\bigg|_F\dd q_j
    +
    \frac{\partial p_i}{\partial F_j}\bigg|_q\dd F_j
  \right).
\end{aligned}
\]
The key claim is
\[
  \frac{\partial p_i}{\partial q_j}\bigg|_F
  =
  \frac{\partial p_j}{\partial q_i}\bigg|_F.
\]
Then the \(\dd q_i\wedge\dd q_j\) contribution vanishes by antisymmetry, and
\[
  \omega=
  \frac{\partial p_i}{\partial F_j}\bigg|_q\dd q_i\wedge\dd F_j.
\]
On the invariant torus \(F_j=f_j\), \(\dd F_j=0\), so
\[
  \omega|_{M_f}=0.
\]

To prove the claim, differentiate the identities \(F_k(p(F,q),q)=F_k\) at
fixed \(F\):
\[
  0=
  \frac{\partial F_k}{\partial p_j}\bigg|_q
  \frac{\partial p_j}{\partial q_i}\bigg|_F
  +
  \frac{\partial F_k}{\partial q_i}\bigg|_p.
\]
Let
\[
  A_{kj}=\frac{\partial F_k}{\partial p_j}\bigg|_q,\qquad
  M_{ji}=\frac{\partial p_j}{\partial q_i}\bigg|_F,\qquad
  B_{ki}=\frac{\partial F_k}{\partial q_i}\bigg|_p.
\]
Then
\[
  AM+B=0,\qquad M=-A^{-1}B.
\]
The involution condition gives
\[
  0=\Poisson{F_i}{F_j}
  =
  A_{ik}B_{jk}-A_{jk}B_{ik},
\]
or
\[
  AB^T-BA^T=0.
\]
Hence
\[
  A^{-1}B=B^T(A^T)^{-1}=(A^{-1}B)^T,
\]
so \(M=M^T\), proving the symmetry claim.
\end{derivation}

\section{Action Variables And Angle Variables}
\label{sec:action-angle-variables}

Choose a basis of one-cycles \(\gamma_1,\ldots,\gamma_n\) on an invariant torus
\(M_f\). In canonical coordinates the action variables are
\[
  I_j=\frac{1}{2\pi}\oint_{\gamma_j} p_i\,\dd q^i.
\]
Actions are therefore not arbitrary constants of integration. They are
phase-space areas measured around the independent cycles of the torus. This is
why they are insensitive to small deformations of the cycle and why they
survive as adiabatic invariants in Chapter~\ref{chap:hamilton-jacobi}.
They are invariant under continuous deformations of \(\gamma_j\) on the torus.
If \(\gamma_j\) and \(\gamma_j'\) bound a two-chain \(D\subset M_f\), then by
Stokes' theorem
\[
  \oint_{\gamma_j}p_i\,\dd q^i
  -
  \oint_{\gamma_j'}p_i\,\dd q^i
  =
  \int_D \dd(p_i\,\dd q^i)
  =
  \int_D \omega
  =
  0.
\]
Thus \(I_j\) depends on the torus, not on the representative loop. Under
regularity assumptions, the \(I_j\) can replace the original integrals \(F_j\)
as labels for nearby tori.

To find conjugate angles, construct a local generating function
\[
  S(I,q)
  =
  \int^q p_i(I,q')\,\dd q'^i,
\]
where the integration is performed on a local branch. The compatibility
condition for this definition is
\[
  \frac{\partial p_i}{\partial q^j}\bigg|_I
  =
  \frac{\partial p_j}{\partial q^i}\bigg|_I,
\]
which is another form of \(\omega|_{M_I}=0\). Define
\[
  \theta^j=\frac{\partial S}{\partial I_j}\bigg|_q.
\]
When one goes around the cycle \(\gamma_k\), the multivalued function \(S\)
jumps by
\[
  \Delta_{\gamma_k}S=\oint_{\gamma_k}p_i\,\dd q^i=2\pi I_k.
\]
Therefore
\[
  \Delta_{\gamma_k}\theta^j
  =
  \frac{\partial}{\partial I_j}(2\pi I_k)
  =
  2\pi\delta^j{}_k.
\]
The \(\theta^j\) are angles on the torus:
\[
  \theta^j\sim \theta^j+2\pi.
\]
In these action-angle coordinates the Hamiltonian is a function of the actions
alone,
\[
  H=H(I),
\]
and Hamilton's equations reduce to
\[
  \dot I_j=0,\qquad
  \dot\theta^j=\omega^j(I)=\frac{\partial H}{\partial I_j}.
\]
Thus motion on each invariant torus is linear. If the frequencies are rationally
related, the orbit closes; otherwise the orbit is dense on the torus.

\section{Rigid Body Tori}
\label{sec:rigid-body-tori}

The full rigid body with fixed center of mass has three degrees of freedom. On
\(T^*\SO(3)\) the functions
\[
  H,\qquad J_3,\qquad J^2
\]
are mutually Poisson commuting when \(J_3\) is a fixed spatial component of
angular momentum and \(J^2=\norm{J}^2\). For regular values, their common level
set is a three-torus.

The geometry can be seen in two stages. First, the body angular momentum
\(M=R^TJ\) is constrained by
\[
  J^2=M_1^2+M_2^2+M_3^2=\text{constant}
\]
and
\[
  H=\frac12\left(\frac{M_1^2}{I_1}
  +\frac{M_2^2}{I_2}
  +\frac{M_3^2}{I_3}\right)=\text{constant}.
\]
This puts \(M\) on the closed curve obtained by intersecting a sphere with an
ellipsoid. Second, fixing \(J_3\) constrains the orientation \(R\): the third
spatial component of \(J=RM\) is fixed. For a given \(M\) on the curve, this
condition leaves a circle of allowed attitudes. A circle fibered over the
momentum-space circle gives a two-dimensional torus in the reduced picture, and
including the remaining cyclic angle gives the full three-torus.

The reduced free rigid body has one effective degree of freedom on each angular
momentum sphere. The energy level on that sphere gives closed curves except at
separatrices. This is the concrete meaning of integrability here.

\section{From Rigid Body To Deforming Body}
\label{sec:rigid-to-deforming}

The rigid body has one orientation \(R(t)\). A Cosserat body assigns an
orientation to each material point:
\[
  R(t)
  \quad\longrightarrow\quad
  R(X,t).
\]
The rigid-body angular velocity
\[
  R^T\dot R
\]
becomes a time connection, while the spatial derivatives
\[
  R^T\partial_iR
\]
become material connection components. This is the first bridge from rigid body
mechanics to gauge theory of deforming bodies.

\section{What To Take Forward}
\label{sec:rigid-take-forward}

\begin{itemize}
\item \(\SO(3)\) is a curved configuration space; Euler angles are useful local
coordinates but not the system itself.
\item Body angular momentum \(M\) and spatial angular momentum \(J\) encode the
same vector in different frames. Confusing them hides the reduction.
\item The free rigid body is integrable because energy and angular momentum
constraints reduce the motion to closed curves and tori.
\item Liouville tori arise from commuting Hamiltonian flows, not from a lucky
choice of variables. Action-angle coordinates are coordinates adapted to those
tori.
\item Replacing one global attitude \(R(t)\) by a field \(R(X,t)\) is the
conceptual step from rigid bodies to gauge descriptions of deformable bodies.
\end{itemize}

\section{Associated Demos}
\label{sec:rigid-associated-demos}

The script \path{demos/python/rigid_body_euler_top.py} integrates Euler's
equations and plots invariant drift. The Mathematica script
\path{demos/mathematica/RigidBodyEulerTop.wl} gives a symbolic and numerical
companion. A useful check is to compare the numerical drift in \(H\) and \(C\)
for several time steps and initial conditions near the intermediate axis.

\chapter{Perturbation, Chaos, And Celestial Mechanics}
\label{chap:perturbation-chaos}

\section{What This Chapter Studies}
\label{sec:perturbation-chapter-purpose}

The central example is a massless asteroid perturbed by Jupiter while orbiting
the Sun. The unperturbed problem is Kepler motion, one of the most integrable
systems in classical mechanics. The perturbation introduces resonances, broken
invariant tori, chaotic transport, and possible ejection.

\begin{assumption}[Restricted planar model]
Unless otherwise stated, this chapter uses a planar restricted model:
\begin{enumerate}
\item the Sun has mass \(M_\odot\);
\item Jupiter has mass \(M_J\) and follows a prescribed circular orbit;
\item asteroids are massless test particles;
\item asteroid-asteroid interactions are neglected;
\item all motion is confined to one plane.
\end{enumerate}
This model is not a high-precision Solar System ephemeris. It is a controlled
Hamiltonian laboratory for perturbation, resonance, and transport.
\end{assumption}

\begin{table}[htbp]
\centering
\caption{Notation for perturbation theory, KAM, and the restricted
Sun-Jupiter-asteroid model.}
\label{tab:perturbation-chaos-notation}
\small
\begin{tabular}{p{0.18\textwidth}p{0.48\textwidth}p{0.25\textwidth}}
\toprule
Symbol & Meaning & Units in demos\\
\midrule
\(G\) & gravitational constant & \(4\pi^2\) in AU, yr, solar masses\\
\(M_\odot\) & solar mass & \(1\)\\
\(M_J\) & Jupiter mass & \(9.543\times 10^{-4}M_\odot\)\\
\(r,v,p\) & asteroid position, velocity, canonical momentum & AU, AU/yr, mass-scaled\\
\(h\) & specific Kepler Hamiltonian & energy per unit mass\\
\(\mu\) & reduced mass in two-body problem & mass\\
\(\ell\) & angular momentum per unit mass & AU\(^2\)/yr\\
\(a\) & asteroid semimajor axis & AU\\
\(a_J\) & Jupiter semimajor axis & AU\\
\(e,\varpi\) & eccentricity and longitude of periapse & dimensionless, radians\\
\(n\) & asteroid mean motion & rad/yr\\
\(n_J\) & Jupiter mean motion & rad/yr\\
\(\eps\) & perturbation strength & often proportional to \(M_J/M_\odot\)\\
\((I,\theta)\) & action-angle coordinates & \(I\in U,\theta\in\mathbb T^N\)\\
\(H_0,H_1\) & integrable Hamiltonian and perturbation & model dependent\\
\(\omega_0(I)\) & unperturbed frequency vector & rad/time\\
\(k\) & integer Fourier/resonance vector & element of \(\Z^N\)\\
\(S_1\) & first generating-function correction & canonical perturbation theory\\
\(H_{1k}\) & \(k\)-th Fourier coefficient of \(H_1\) & same units as \(H\)\\
\(\alpha,\tau\) & Diophantine constants & lower-bound parameters\\
\(K\) & standard-map kick strength & dimensionless\\
\(T_K,\Psi\) & standard map or other area-preserving map & discrete-time map\\
\(N,T,\Delta t\) & sample count, integration horizon, time step & demo parameters\\
\(\widehat P_{\mathrm{ej}}\) & finite-time ejection-probability estimate & dimensionless\\
\bottomrule
\end{tabular}
\end{table}

\paragraph{Roadmap.}
The chapter begins with the Kepler problem because perturbation theory needs an
integrable reference system. It then introduces resonances, small divisors, and
Diophantine conditions, which explain why some invariant tori are robust while
others are fragile. The KAM theorem gives the positive result: many sufficiently
irrational tori persist under small perturbation. The Poincare-map and
standard-map sections show the complementary phenomenon: resonant tori break
into islands and chaotic layers. The restricted Sun-Jupiter-asteroid model then
turns this geometry into a numerical ejection experiment.

There are three levels of description to keep separate. The analytic level asks
whether a formal canonical transformation can remove angle dependence. The
geometric level asks which invariant tori, islands, and separatrices organize
phase space. The numerical level asks how an ensemble moves through that
geometry over a finite time. The asteroid ejection probability belongs to the
third level, but it is meaningful only because the first two levels explain
where transport channels come from.

\paragraph{How approximations are labelled.}
Perturbation theory is useful only when the bookkeeping is explicit. In this
chapter \(H_0\) denotes the exactly integrable reference problem, \(\eps H_1\)
denotes a controlled perturbation, \(O(\eps^2)\) denotes terms dropped by a
first-order calculation, and resonance conditions such as \(k\cdot\omega=0\)
identify where that first-order calculation is least trustworthy. Numerical
experiments later in the chapter do not replace these estimates; they test the
transport suggested by them over a stated finite time horizon.

\section{The Kepler Problem}
\label{sec:kepler-perturbation}

For two bodies of masses \(m_1,m_2\), let
\[
  r=x_1-x_2
\]
be the relative position and
\[
  \mu=\frac{m_1m_2}{m_1+m_2}
\]
the reduced mass. The relative Hamiltonian is
\[
  H(r,p)=\frac{\norm{p}^2}{2\mu}-\frac{Gm_1m_2}{\norm r}.
\]
For a massless asteroid around the Sun, one usually divides by asteroid mass and
works with the specific Hamiltonian
\[
  h(r,v)=\frac12\norm v^2-\frac{GM_\odot}{r},
  \qquad r=\norm r.
\]

\begin{derivation}[Kepler orbit equation]
The angular momentum per unit mass is
\[
  \ell=r^2\dot\theta.
\]
This \(\ell\) is the specific angular momentum. In
Chapter~\ref{chap:lag-ham}, the two-body Kepler derivation used the massful
angular momentum \(\ell_{\mathrm{mass}}=\mu r^2\dot\theta\); the two symbols
differ by the reduced mass because this chapter has divided the Hamiltonian by
the asteroid mass. Section~\ref{subsec:detailed-binet} gives the same Binet
equation with an arbitrary particle mass \(m\); setting \(m=1\) gives the
specific convention used here.
For a central potential, \(\ell\) is conserved. The radial equation is
\[
  \ddot r-r\dot\theta^2=-\frac{GM_\odot}{r^2}.
\]
Set \(u=1/r\). Using
\[
  \frac{\dd}{\dd t}=\dot\theta\frac{\dd}{\dd\theta}
  =
  \ell u^2\frac{\dd}{\dd\theta},
\]
the radial equation loses its explicit time dependence and becomes Binet's
equation for the orbit shape:
\[
  \frac{\dd^2u}{\dd\theta^2}+u=\frac{GM_\odot}{\ell^2}.
\]
Therefore
\[
  u(\theta)=\frac{GM_\odot}{\ell^2}\left(1+e\cos(\theta-\varpi)\right),
\]
or
\[
  \boxed{
  r(\theta)=\frac{a(1-e^2)}{1+e\cos(\theta-\varpi)}
  }.
\]
Here \(a\) is the semimajor axis, \(e\) is eccentricity, and \(\varpi\) is the
longitude of periapse.
\end{derivation}

Bound Kepler orbits have energy
\[
  h=-\frac{GM_\odot}{2a}.
\]
Kepler's third law follows from the period \(T\):
\[
  n^2a^3=GM_\odot,\qquad n=\frac{2\pi}{T}.
\]

\section{Action-Angle View Of Kepler Motion}
\label{sec:kepler-action-angle}

The Kepler problem is more than solvable; it is superintegrable. For
perturbation theory, the most useful idea is that unperturbed motion can be
written in action-angle form:
\[
  H_0=H_0(I),
  \qquad
  \dot I=0,
  \qquad
  \dot\theta=\omega(I).
\]
For a near-integrable Hamiltonian,
\[
  H(I,\theta)=H_0(I)+\eps H_1(I,\theta),
\]
the actions are no longer exactly constant:
\[
  \dot I=-\eps\frac{\partial H_1}{\partial \theta},
  \qquad
  \dot\theta=\frac{\partial H_0}{\partial I}
  +
  \eps\frac{\partial H_1}{\partial I}.
\]
The small parameter \(\eps\) should be interpreted physically. In the
Sun-Jupiter-asteroid problem, a natural scale is \(M_J/M_\odot\), though close
encounters and resonances can amplify its effect.

\section{Resonance}
\label{sec:resonance}

\begin{definition}[Resonance]
A resonance occurs when there is a nonzero integer vector \(k\in\Z^n\) such
that
\[
  k\cdot\omega(I)=0.
\]
Near resonance, phase combinations that normally average out become slow and
can accumulate over long times.
\end{definition}

For mean-motion resonance with Jupiter, a \(p:q\) resonance means
\[
  \frac{n}{n_J}\approx \frac{p}{q}.
\]
Using Kepler's law,
\[
  \frac{n}{n_J}\approx
  \left(\frac{a_J}{a}\right)^{3/2},
\]
so the nominal resonance location is
\[
  \boxed{
  a_{p:q}=a_J\left(\frac{q}{p}\right)^{2/3}
  }.
\]
This formula is only the center of a resonance zone. The width of the zone
depends on the perturbation strength, eccentricity, and the order of the
resonance. The point of marking the nominal location is to identify where slow
angle combinations can form; the later chaos discussion explains what happens
inside and near those zones.

\begin{figure}[htbp]
\centering
\begin{tikzpicture}[x=0.72cm,y=1cm, >=Latex]
  \draw[->] (0,0) -- (13.8,0) node[right] {\(a\) [AU]};
  \foreach \x/\lab in {0/2.0,2/2.5,4/3.0,6/3.5,8/4.0,10/4.5,12/5.0}
    \draw (\x,0.07) -- (\x,-0.07) node[below] {\lab};
  \draw[ultra thick, blue!35] (0.1,0) -- (7.0,0);
  \node[blue!60!black] at (3.6,0.45) {asteroid-belt-like annulus};
  \foreach \x/\lab in {2.008/3:1,3.302/5:2,3.833/7:3,5.114/2:1}
  {
    \draw[red, thick] (\x,-0.55) -- (\x,0.9);
    \node[red, rotate=90, anchor=south] at (\x,0.95) {\lab};
  }
  \fill[orange!80!black] (12.82,0) circle (2.2pt);
  \node[orange!80!black, above] at (12.82,0.12) {Jupiter};
\end{tikzpicture}
\caption{Nominal mean-motion resonance locations with Jupiter, shown
across an asteroid-belt-like range. The resonance positions use
\(a_{p:q}=a_J(q/p)^{2/3}\); the figure now includes Jupiter's orbital radius to
show that the listed resonances lie interior to Jupiter.}
\label{fig:jupiter-resonance-locations}
\end{figure}

\section{Perturbing An Integrable System}
\label{sec:perturbing-integrable-system}

The central perturbation question is the following. Suppose the unperturbed
Hamiltonian is integrable and has action-angle coordinates:
\[
  H_0=H_0(I),
  \qquad
  (I,\theta)\in U\times\mathbb T^N,
  \qquad
  \dot I=0,\quad \dot\theta=\omega_0(I),
\]
where
\[
  \omega_0(I)=\frac{\partial H_0}{\partial I}.
\]
Now add a small generic perturbation:
\[
  H(I,\theta)=H_0(I)+\eps H_1(I,\theta),
  \qquad
  0<\eps\ll 1.
\]
What happens to the invariant tori \(I=\text{constant}\)?

Historically, two tempting answers were both too simple. One might hope that
the perturbed system is still integrable after a clever change of variables.
One might also expect that a generic perturbation destroys the tori and makes a
typical orbit wander ergodically through the available energy surface. KAM
theory gives the subtler answer: for small analytic perturbations, most
sufficiently irrational tori persist, while resonant regions create islands,
separatrices, and chaotic transport.

The historical landmarks are the Fermi-Pasta-Ulam-Tsingou numerical experiment
of 1953, which contradicted naive ergodicity, and the
Kolmogorov-Arnold-Moser work from roughly 1954--1963, which explained why
small perturbations can leave most invariant tori intact while still allowing
chaotic regions.

\section{Canonical Perturbation Theory And Small Divisors}
\label{sec:canonical-perturbation-small-divisors}

We first ask whether a canonical transformation can remove the angle dependence
order by order. Use a type-II generating function
\[
  S(I',\theta)=I'\cdot\theta+\eps S_1(I',\theta)+\eps^2S_2+\cdots.
\]
The transformation is
\[
  I_i=\frac{\partial S}{\partial \theta^i},
  \qquad
  \theta'^i=\frac{\partial S}{\partial I'_i}.
\]
To first order,
\[
  I=I'+\eps\frac{\partial S_1}{\partial \theta}+O(\eps^2).
\]
If the transformed Hamiltonian is angle-independent to order \(\eps\),
\[
  H'(I')=H(I,\theta),
\]
then expanding gives
\[
  H'(I')
  =
  H_0(I')
  +
  \eps\left[
    \omega_0(I')\cdot\frac{\partial S_1}{\partial \theta}
    +
    H_1(I',\theta)
  \right]
  +
  O(\eps^2).
\]
Thus we need
\[
  \omega_0(I')\cdot\frac{\partial S_1}{\partial \theta}
  =
  f_1(I')-H_1(I',\theta),
\]
where \(f_1(I')\) will become the first correction to the new Hamiltonian.
Since \(S_1\) must be periodic in each angle, the average over the torus of the
left side is zero. Therefore
\[
  f_1(I')=\langle H_1(I',\theta)\rangle_\theta.
\]

Expand the perturbation in Fourier modes:
\[
  H_1(I,\theta)=\sum_{k\in\Z^N} H_{1k}(I)e^{ik\cdot\theta}.
\]
The \(k=0\) mode is the average. For \(k\neq0\), the homological equation gives
\[
  i(k\cdot\omega_0(I'))S_{1k}(I')=-H_{1k}(I'),
\]
or
\[
  \boxed{
  S_{1k}(I')
  =
  \frac{iH_{1k}(I')}{k\cdot\omega_0(I')}
  \qquad (k\neq0).
  }
\]
The denominators \(k\cdot\omega_0\) are the small divisors.

\begin{definition}[Resonant torus]
The unperturbed torus with action \(I\) is resonant if there exists
\[
  k\in\Z^N\setminus\{0\}
\]
such that
\[
  k\cdot\omega_0(I)=0.
\]
On such a torus the phase \(k\cdot\theta\) is constant along the unperturbed
motion, so the corresponding perturbation does not average away.
\end{definition}

Even if no divisor is exactly zero, the Fourier series for \(S_1\) can diverge
if some \(k\cdot\omega_0\) are too small too often. This is the analytic heart
of KAM theory.

\section{Fourier Decay And The Diophantine Condition}
\label{sec:diophantine-condition}

Small divisors can be overcome when the numerators decay rapidly. If
\(H_1(I,\theta)\) is real analytic in \(\theta\), it extends to a complex strip
\[
  |\operatorname{Im}\theta^i|<\sigma.
\]
Shifting the Fourier integral into the complex strip gives the estimate
\[
  |H_{1k}(I)|\leq C(I)e^{-\sigma |k|_1},
  \qquad
  |k|_1=|k_1|+\cdots+|k_N|.
\]
Thus large Fourier modes are exponentially suppressed.
The estimate is the analytic reason KAM is possible: small divisors grow at
worst polynomially under a Diophantine condition, while analytic Fourier
coefficients shrink exponentially. Smooth but nonanalytic perturbations require
more delicate bookkeeping because the numerator may decay only as a power of
\(|k|\).

\begin{example}[Fourier decay distinguishes analytic and nonanalytic functions]
One-angle examples make the difficulty visible. The series
\[
  \sum_{k=1}^{\infty}\frac{\sin k\theta}{k}
\]
has Fourier coefficients that decay only like \(1/k\); its graph has a jump and
is not real analytic. Likewise
\[
  \sum_{k=1}^{\infty}\frac{\cos k\theta}{k^2}
  =
  \frac14(\theta-\pi)^2-\frac{\pi^2}{12},
  \qquad |\theta|<2\pi,
\]
is periodic but not analytic at the endpoints of a period. A genuinely
real-analytic example is
\[
  \sum_{k=-\infty}^{\infty}\cos(k\theta)e^{-k^2/2},
\]
a theta-function-type series with rapidly suppressed Fourier coefficients.
\end{example}

The frequency vector must also avoid being too well approximated by resonances.

\begin{definition}[Diophantine frequency]
A vector \(\omega\in\R^N\) is Diophantine with constants \(\alpha>0\) and
\(\tau>N-1\) if
\[
  |k\cdot\omega|
  \geq
  \frac{\alpha}{|k|_1^\tau}
  \qquad
  \text{for all } k\in\Z^N\setminus\{0\}.
\]
\end{definition}

For Diophantine \(\omega_0(I')\),
\[
  |S_{1k}(I')|
  \leq
  \frac{C(I')}{\alpha}|k|_1^\tau e^{-\sigma |k|_1},
\]
so the first-order generating function is well-defined. The polynomial factor
from the small divisor cannot beat the exponential decay from analyticity.

\begin{example}[How irrational is irrational?]
For \(N=2\), the Diophantine inequality can be recast as a rational
approximation condition for
\[
  x=\frac{\omega_1}{\omega_2}.
\]
It asks, roughly, that
\[
  \left|x-\frac{p}{q}\right|
  \geq
  \frac{A}{q^{\tau+1}}
\]
for all rational approximants \(p/q\). For \(x=\sqrt2\), use
\[
  f(x)=x^2-2=0.
\]
If \(p/q=x-\eta\), then
\[
  f(p/q)=\frac{p^2-2q^2}{q^2}.
\]
The numerator is a nonzero integer unless \(p/q=\sqrt2\), so its absolute value
is at least one. Linearizing
\[
  f(x-\eta)=f(x)-\eta f'(x)+O(\eta^2)
\]
with \(f'(x)=2\sqrt2\) gives
\[
  |\eta|\geq \frac{A}{q^2}
\]
for some \(A>0\). Thus \(\sqrt2\) satisfies a Diophantine-type lower bound.
More generally, algebraic irrational numbers obey polynomial rational-
approximation bounds. By contrast, numbers such as
\[
  x=\sum_{n=0}^{\infty}\frac{1}{2^{n!}}
\]
are too well approximated by rationals and fail every fixed lower bound of this
kind.
\end{example}

The Diophantine condition excludes a dense set of resonant hyperplanes, but the
excluded set has small measure when \(\alpha\) is small. To see the estimate,
fix a bounded frequency domain \(\Omega\subset\R^N\). For a given
\(k\neq0\), the bad set is the slab
\[
  R_k
  =
  \left\{
    \omega\in\Omega:
    |k\cdot\omega|<
    \frac{\alpha}{|k|_1^\tau}
  \right\}.
\]
Its thickness normal to the hyperplane \(k\cdot\omega=0\) is of order
\[
  \frac{\alpha}{|k|_1^{\tau+1}},
\]
because the normal vector has length comparable to \(|k|_1\). Hence
\[
  \operatorname{Vol}(R_k\cap\Omega)
  \leq
  C_\Omega\frac{\alpha}{|k|_1^{\tau+1}}.
\]
Summing over \(k\neq0\) converges when \(\tau+1>N\), i.e. \(\tau>N-1\). Thus
the non-Diophantine set can be made to have arbitrarily small measure by taking
\(\alpha\) small, even though resonances are dense.

\section{KAM Theorem: Statement And Sketch}
\label{sec:kam-theorem-sketch}

At this point the perturbation problem has two competing facts. Analyticity
makes high Fourier modes small, which helps the perturbation series. Small
divisors make near-resonant modes large, which threatens the perturbation
series. KAM theory is the result of balancing these facts: it restricts
attention to frequencies that are irrational in a quantified way, and it asks
the unperturbed Hamiltonian to twist enough that frequency can be adjusted when
the torus is deformed.

The theorem is best read as a statement about labelled invariant tori, not
about individual trajectories. An unperturbed action value \(I\) determines a
frequency \(\omega_0(I)=\partial H_0/\partial I\). If that frequency is
Diophantine and the twist condition below holds, then the perturbed system has
a nearby torus carrying quasiperiodic motion with the corresponding frequency.
Points on resonant or nearly resonant tori are precisely the points for which
this labelling argument breaks down.

\begin{assumption}[Kolmogorov nondegeneracy]
The integrable Hamiltonian \(H_0(I)\) has a locally nondegenerate frequency map:
\[
  \det\left(\frac{\partial^2H_0}{\partial I_i\partial I_j}\right)\neq0.
\]
Equivalently, nearby actions give genuinely different frequency vectors.
\end{assumption}

\begin{proposition}[KAM theorem, lecture-note form]
Let
\[
  H(I,\theta)=H_0(I)+\eps H_1(I,\theta)
\]
be a sufficiently smooth, typically analytic, nearly integrable Hamiltonian on
\(U\times\mathbb T^N\). Assume \(H_0\) is Kolmogorov nondegenerate. For
sufficiently small \(|\eps|\), every unperturbed torus whose frequency is
Diophantine with fixed constants \((\alpha,\tau)\) persists as a nearby
invariant torus of the perturbed system, provided \(\eps\) is small enough
relative to \(\alpha\). The perturbed tori fill phase space up to a set whose
measure tends to zero as \(\eps\to0\).
\end{proposition}

The proof is not ordinary first-order perturbation theory. The first-order
calculation above shows the mechanism, but higher orders produce new small
divisors and require control of the coordinate change itself:
\[
  I=I'+\eps\frac{\partial S_1}{\partial\theta}+\cdots,
  \qquad
  \theta'=\theta+\eps\frac{\partial S_1}{\partial I'}+\cdots.
\]
The important scale warning is:
\[
  \frac{\partial S_1}{\partial I'},\
  \frac{\partial S_1}{\partial \theta}
  =
  O(\alpha^{-1}),
\]
so the coordinate change is small only when \(\eps/\alpha\) is small. A common
lecture-note form is therefore: for \(|\eps|<\delta\alpha^2\), the tori whose
frequencies satisfy the Diophantine inequality with constant \(\alpha\) persist
and fill phase space up to \(O(\alpha)\) measure.

The point of the exponent and smallness conditions is bookkeeping, not
mysticism. If \(H_1\) is analytic in a strip of width \(r\), its Fourier
coefficients satisfy a bound of the form
\[
  |H_{1,k}|\leq C e^{-r|k|}.
\]
Solving the homological equation divides by \(k\cdot\omega\), and the
Diophantine condition gives
\[
  \left|\frac{H_{1,k}}{k\cdot\omega}\right|
  \leq
  \frac{C}{\alpha}|k|^\tau e^{-r|k|}.
\]
The polynomial factor \(|k|^\tau\) is harmless compared with exponential
decay, provided one gives up a little analyticity width. This is the analytic
heart of the proof: small divisors are dangerous, but not fatal, on the
Diophantine set.

KAM uses a Newton-style iteration. At each step, one solves a homological
equation like the first-order equation, removes most of the angle dependence,
and shrinks the analytic domain to pay for derivative losses. The error is made
quadratically smaller at each step. Diophantine estimates keep the small
divisors under control; nondegeneracy lets one adjust the action label so the
desired frequency remains fixed.

In proof-outline form, the iteration has four repeating moves:
\[
\begin{array}{rcl}
\text{linearize} &:& \text{separate resonant size from nonresonant modes},\\
\text{solve} &:& \text{invert } \omega\cdot\partial_\theta
    \text{ using the Diophantine bound},\\
\text{renormalize} &:& \text{change the action label using nondegeneracy},\\
\text{shrink} &:& \text{reduce the analytic domain so derivatives stay controlled}.
\end{array}
\]
The theorem is difficult because all four moves must be made quantitatively at
once. The conceptual message is simpler: sufficiently irrational frequencies
avoid the worst resonances, while twist lets the torus move to the nearby
frequency that the perturbed system actually supports.

One useful way to read the theorem is as a filtered survival statement. Start
with all invariant tori of the integrable system. Remove narrow neighborhoods
of the low-order resonance surfaces \(k\cdot\omega=0\). The remaining set is
not an interval of actions but a Cantor-like family: it is full of holes at all
orders, yet the total measure of the holes can be made small. The KAM iteration
then proves that each torus in this filtered family has a nearby perturbed
torus with conjugate quasiperiodic motion.

\begin{figure}[htbp]
\centering
\begin{tikzpicture}[scale=1.0, >=Latex]
  \draw[->] (-0.2,0) -- (7.2,0) node[right] {action label \(I\)};
  \draw[thick, blue!70!black] (0.4,0) -- (1.25,0);
  \draw[thick, blue!70!black] (1.55,0) -- (2.35,0);
  \draw[thick, blue!70!black] (2.62,0) -- (3.55,0);
  \draw[thick, blue!70!black] (3.85,0) -- (4.18,0);
  \draw[thick, blue!70!black] (4.52,0) -- (5.42,0);
  \draw[thick, blue!70!black] (5.78,0) -- (6.65,0);
  \foreach \x/\w in {1.4/0.18,2.48/0.10,3.7/0.18,4.34/0.16,5.60/0.18}
  {
    \draw[red!70!black, very thick] (\x-\w,0.18) -- (\x+\w,-0.18);
    \draw[red!70!black, very thick] (\x-\w,-0.18) -- (\x+\w,0.18);
  }
  \node[blue!70!black, align=center] at (2.0,0.8)
    {Diophantine\\surviving labels};
  \node[red!70!black, align=center] at (5.25,0.9)
    {resonant gaps\\and chaotic layers};
  \draw[->, gray!70] (1.4,-0.75) -- (1.4,-0.18)
    node[below=0.25cm, align=center] {low-order\\resonance};
  \draw[decorate, decoration={brace, amplitude=4pt}, thick]
    (0.4,-0.55) -- (6.65,-0.55)
    node[midway, below=5pt] {large measure remains when \(\eps\) is small};
\end{tikzpicture}
\caption{KAM survival is Cantor-like. Resonant action labels are removed in
thin layers around \(k\cdot\omega=0\); the surviving Diophantine labels are
nowhere all of the line, but their total measure approaches the integrable
measure as \(\eps\to0\).}
\label{fig:kam-cantor-survival}
\end{figure}

The theorem does not say that every torus survives. Resonant tori are destroyed
or reorganized into lower-dimensional structures. Nor does it say that all
surviving tori are easy to see numerically. In systems with \(N=2\), surviving
one-dimensional invariant curves on a Poincare section can block transport. In
higher-dimensional systems, the geometry leaves more room for slow drift
through resonance channels.

\section{From Resonance To Chaos: Poincare Map Picture}
\label{sec:poincare-resonance-chaos}

A Poincare section converts continuous Hamiltonian motion into an
area-preserving map. This loses the time parametrization but keeps the geometry
needed to see resonance. Near a resonant torus, some iterate of the unperturbed
map is the identity; after perturbation, that near-identity map unfolds into
fixed points, island chains, and separatrices.

This is why maps are so useful in a chapter about flows. A continuous orbit can
hide its organization in time; a section records each return to a transverse
surface and turns invariant tori into curves, resonances into island chains,
and chaotic layers into visible scatter. The standard map below is not a toy
because it is simple; it is useful because it isolates this return-map
geometry.

For two degrees of freedom, begin with an integrable Hamiltonian
\[
  H_0=H_0(I_1,I_2),
  \qquad
  (I_1,I_2,\theta_1,\theta_2)\in U\times\mathbb T^2,
\]
and fix an energy surface \(H_0=E\). Choose a Poincare section \(S\), for
example \(\theta_2=0\). Starting from \(x\in S\), follow the orbit until it
next intersects \(S\). This defines the return map
\[
  T:S\to S.
\]
The unperturbed flow has
\[
  \dot\theta_1=\omega_1(I),
  \qquad
  \dot\theta_2=\omega_2(I).
\]
Returning to \(\theta_2=0\) takes time \(2\pi/\omega_2\), so the return map
twists the remaining angle by
\[
  \Delta\theta_1=2\pi\frac{\omega_1}{\omega_2}.
\]
On a resonant torus with \(\omega_1/\omega_2=p/q\),
\[
  T^q=\id
\]
for the unperturbed return map \(T\). A concrete period-three resonance is
\[
  \frac{\omega_1}{\omega_2}=\frac13,
  \qquad
  T:\theta_1\mapsto\theta_1+\frac{2\pi}{3},
  \qquad
  T^3=\id.
\]
Let \(C\subset S\) be the invariant curve obtained by intersecting this
resonant torus with the section. Nearby invariant curves \(C_-\) and \(C_+\)
have slightly different rotation numbers. With the usual twist assumption,
\(T^3\) rotates points one way on \(C_-\) and the opposite way on \(C_+\).

Now perturb the Hamiltonian:
\[
  H=H_0+\eps H_1.
\]
The perturbed Poincare map \(T_\eps\) is still defined by first return to
\(S\), even though the invariant tori may no longer exist. Near
\(C\cap S\), \(T_\eps^3\) is close to the identity. At fixed energy it has the
local form
\[
  T_\eps^3(I_1,\theta_1)
  =
  \left(
    I_1+f_\eps(I_1,\theta_1),
    \theta_1+g_\eps(I_1,\theta_1)
  \right).
\]
For small \(\eps\), the opposite twists on \(C_-\) and \(C_+\) persist.
Therefore an intermediate curve \(C'\) has
\[
  g_\eps(I_1,\theta_1)=0;
\]
on \(C'\), the third-return map produces no angular displacement.

The missing radial displacement is controlled by area preservation. Hamiltonian
flow preserves \(\omega\), and hence the Poincare map preserves the induced
area form on the section. Equivalently, for a closed curve \(\gamma\) bounding
a disk \(D_\gamma\) in the section,
\[
  \oint_\gamma p\,\dd q
  =
  -\int_{D_\gamma}\omega,
\]
so the signed symplectic area enclosed by \(\gamma\) is invariant. If
\(f_\eps\) were nonzero with a fixed sign everywhere on \(C'\), the curve
would be displaced monotonically inward or outward under \(T_\eps^3\), changing
the enclosed area. Thus \(f_\eps\) must have at least a pair of zeros on
\(C'\). At such points
\[
  f_\eps=0,
  \qquad
  g_\eps=0,
\]
so \(T_\eps^3\) has fixed points. In the period-three example, if \(x\) is a
fixed point of \(T_\eps^3\), then \(T_\eps(x)\) is also a fixed point of
\(T_\eps^3\). Thus the resonant curve generically produces six fixed points:
three elliptic and three hyperbolic, alternating around the resonant layer.

Elliptic points are surrounded by island chains and new local invariant curves.
Hyperbolic points have stable and unstable manifolds. Once those manifolds
intersect transversely, the separatrix is replaced by a chaotic layer. Farther
from the resonance, sufficiently Diophantine KAM curves persist and confine
motion.

This is the local mechanism behind the pictures one sees in the standard map,
the driven pendulum, and the asteroid-belt resonances: KAM tori survive away
from strong resonances, while resonant tori break into islands and chaotic
separatrix webs.

\section{The Standard Map As A Minimal Model}
\label{sec:standard-map}

The standard map is
\[
\begin{aligned}
  p_{n+1} &= p_n+K\sin q_n,\\
  q_{n+1} &= q_n+p_{n+1}\pmod{2\pi}.
\end{aligned}
\]
This is naturally a map of the cylinder, with \(q\) taken modulo \(2\pi\) and
\(p\in\R\). For phase portraits on a compact square, the Python demo reduces
both \(q\) and \(p\) modulo \(2\pi\), giving the corresponding map on the
two-torus.
It is area-preserving because the Jacobian matrix has determinant one:
\[
  D\Psi=
  \begin{pmatrix}
  1+K\cos q_n & 1\\
  K\cos q_n & 1
  \end{pmatrix},
  \qquad
  \det D\Psi=1.
\]
The parameter \(K\) controls resonance overlap. For small \(K\), invariant
curves dominate. For larger \(K\), chaotic regions grow.

\begin{example}[Period-three resonance in the standard map]
A common drift-kick convention for the standard map is
\[
  T_K(\theta,p)=\left(\theta+p,\ p+K\sin(\theta+p)\right),
\]
which is symplectically equivalent to the convention above after shifting where
the kick is placed. At \(K=0\), the invariant curve
\[
  C:\quad p=\frac{2\pi}{3}
\]
has
\[
  T_0^3=\id
\]
because the angle advances by \(2\pi\) after three iterates.

Let
\[
  p=\frac{2\pi}{3}+\delta(\theta),
  \qquad
  \delta=O(K).
\]
To first order in \(K\) and \(\delta\),
\[
\begin{aligned}
  \theta_3-\theta-2\pi
  &=
  3\delta
  +
  2K\sin\left(\theta+\frac{2\pi}{3}\right)
  +
  K\sin\left(\theta+\frac{4\pi}{3}\right)
  +O(K^2).
\end{aligned}
\]
The nearby curve \(C'\) on which \(T_K^3\) fixes the angle therefore satisfies
\[
  \boxed{
  \delta(\theta)
  =
  \frac{K}{2}\sin\theta
  -
  \frac{K\sqrt3}{6}\cos\theta
  +
  O(K^2).
  }
\]
For example, at \(\theta=\pi/3\),
\[
  p=\frac{2\pi}{3}+\frac{K}{2\sqrt3}+O(K^2),
\]
which is the initial condition singled out in the corresponding numerical
exercise. At the next order, area preservation forces the radial displacement
on \(C'\) to have alternating zeros. These become the six period-three points:
three elliptic island centers and three hyperbolic points with separatrix
structure.
\end{example}

\begin{figure}[htbp]
\centering
\begin{tikzpicture}[scale=0.92]
  \draw[->] (-0.2,0) -- (6.5,0) node[right] {\(q\)};
  \draw[->] (0,-0.2) -- (0,4.8) node[above] {\(p\)};
  \draw (0,0.05) -- (0,-0.05) node[below] {\(0\)};
  \draw (1.5708,0.05) -- (1.5708,-0.05) node[below] {\(\pi/2\)};
  \draw (3.1416,0.05) -- (3.1416,-0.05) node[below] {\(\pi\)};
  \draw (4.7124,0.05) -- (4.7124,-0.05) node[below] {\(3\pi/2\)};
  \draw (6.2832,0.05) -- (6.2832,-0.05) node[below] {\(2\pi\)};
  \draw (0.05,2.094) -- (-0.05,2.094) node[left] {\(2\pi/3\)};
  \draw (0.05,4.189) -- (-0.05,4.189) node[left] {\(4\pi/3\)};
  \foreach \y in {0.65,1.15,3.25,3.85}
    \draw[blue, thick, domain=0.4:6.0, samples=100]
      plot (\x,{\y+0.15*sin(100*\x+\y*20)});
  \draw[dashed, gray] (0.25,2.1) -- (6.2,2.1)
    node[right] {\(p=2\pi/3\)};
  \foreach \x in {1.05,3.15,5.25}
  {
    \draw[red!75!black, thick] (\x,2.1) ellipse (0.48 and 0.27);
    \fill[red!75!black] (\x,2.1) circle (1.1pt);
  }
  \draw[red!75!black, densely dashed]
    (1.53,2.1) .. controls (2.0,2.55) and (2.6,2.55) .. (3.15,2.37)
    .. controls (3.7,2.55) and (4.6,2.55) .. (4.77,2.1);
  \draw[red!75!black, densely dashed]
    (1.53,2.1) .. controls (2.0,1.65) and (2.6,1.65) .. (3.15,1.83)
    .. controls (3.7,1.65) and (4.6,1.65) .. (4.77,2.1);
  \foreach \x in {2.10,4.20,6.05}
  {
    \draw[purple!75!black, thick] (\x-0.12,1.98) -- (\x+0.12,2.22);
    \draw[purple!75!black, thick] (\x-0.12,2.22) -- (\x+0.12,1.98);
  }
  \foreach \x/\y in {0.8/0.95,1.6/3.55,2.9/1.35,4.7/3.75,5.8/0.85}
    \fill[black!55] (\x,\y) circle (1.0pt);
  \node[blue] at (1.45,4.45) {surviving KAM curves};
  \node[red!75!black] at (3.15,2.85) {period-three island chain};
  \node[purple!75!black] at (5.1,1.45) {hyperbolic points};
\end{tikzpicture}
\caption{Schematic phase portrait of an area-preserving map near a
period-three resonance: invariant curves, a chain of three islands, hyperbolic
points, and scattered chaotic points coexist.}
\label{fig:standard-map-island-chain}
\end{figure}

\section{Restricted Sun-Jupiter-Asteroid Equations}
\label{sec:restricted-sun-jupiter-asteroid}

Let \(r\in\R^2\) be the asteroid position relative to the Sun. Jupiter is
prescribed by
\[
  r_J(t)=a_J(\cos n_Jt,\sin n_Jt),
  \qquad
  n_J^2=\frac{G(M_\odot+M_J)}{a_J^3}.
\]
In a heliocentric non-inertial frame, the asteroid acceleration is
\[
  \boxed{
  \ddot r
  =
  -\frac{GM_\odot r}{\norm r^3}
  -
  GM_J
  \left(
    \frac{r-r_J}{\norm{r-r_J}^3}
    +
    \frac{r_J}{\norm{r_J}^3}
  \right)
  }.
\]
The first term is solar gravity. The first part of the Jupiter term is direct
gravitational attraction to Jupiter. The final term is the indirect term caused
by the acceleration of the heliocentric frame.

\begin{derivation}[Indirect term]
Let \(x\) be the asteroid position in an inertial frame and \(x_\odot\) the Sun
position. The heliocentric vector is \(r=x-x_\odot\). Then
\[
  \ddot r=\ddot x-\ddot x_\odot.
\]
The direct Jupiter acceleration of the asteroid is
\[
  -GM_J\frac{x-x_J}{\norm{x-x_J}^3}.
\]
The acceleration of the Sun due to Jupiter is
\[
  \ddot x_\odot=
  GM_J\frac{x_J-x_\odot}{\norm{x_J-x_\odot}^3}
  =
  GM_J\frac{r_J}{\norm{r_J}^3}.
\]
Subtracting \(\ddot x_\odot\) produces the indirect term
\[
  -GM_J\frac{r_J}{\norm{r_J}^3}.
\]
\end{derivation}

\section{Numerical Ejection Probability}
\label{sec:numerical-ejection-probability}

The highlighted simulation samples \(N\) test particles with initial semimajor
axes \(a_0\in[a_{\min},a_{\max}]\), small eccentricities, random orbital phases,
and random periapse directions. Let \(T\) be the integration horizon.

\begin{definition}[Finite-time ejection probability]
For a specified model, sampling distribution, time step, and ejection criterion,
the finite-time ejection probability estimate is
\[
  \widehat P_{\mathrm{ej}}(T)
  =
  \frac{1}{N}
  \sum_{i=1}^{N}\mathbf 1\{\text{particle \(i\) is ejected before time \(T\)}\}.
\]
\end{definition}

\begin{assumption}[Ejection criterion in the demo]
The Python demo counts a particle as ejected if it reaches a large heliocentric
radius or if its heliocentric two-body energy is positive after it has moved
outside Jupiter's orbital scale. It separately records close approaches to
Jupiter and collisions with the inner solar cutoff.
\end{assumption}

This probability is not an intrinsic number attached to nature. It depends on:
\[
  \text{model}+\text{initial ensemble}+\text{time horizon}
  +\text{numerical method}+\text{loss criterion}.
\]
That dependence should be reported, not hidden.

\section{Integrator Used In The Demo}
\label{sec:asteroid-integrator}

The demo uses a kick-drift-kick style update. If \(a(r,t)\) is the acceleration,
then
\[
\begin{aligned}
  v_{n+1/2} &= v_n+\frac{\Delta t}{2}a(r_n,t_n),\\
  r_{n+1} &= r_n+\Delta t\,v_{n+1/2},\\
  v_{n+1} &= v_{n+1/2}+\frac{\Delta t}{2}a(r_{n+1},t_{n+1}).
\end{aligned}
\]
Because Jupiter's position is explicitly time-dependent in the prescribed
circular model, the method is not exactly the same as a symplectic map for an
autonomous finite-dimensional Hamiltonian. It is still much better pedagogically
than an opaque black-box integrator because every force evaluation is visible.

\section{What To Take Forward}
\label{sec:perturbation-chaos-take-forward}

\begin{itemize}
\item Integrable systems are organized by invariant tori, not merely by closed
form solutions.
\item Resonance is a geometric condition \(k\cdot\omega=0\). It marks where
averaging fails and where perturbations can accumulate coherently.
\item Small divisors are not a technical nuisance; they are the arithmetic
trace of dense resonances.
\item KAM theory says that sufficiently irrational tori persist, but it also
highlights what is lost: resonant tori reorganize into island chains,
separatrices, and chaotic transport channels.
\item Numerical ejection probabilities are finite-time, model-dependent
observables. They become meaningful only when the sampling rule, integration
time, timestep, and loss criterion are stated.
\end{itemize}

\section{Questions For The Reader}
\label{sec:asteroid-reader-questions}

\begin{enumerate}
\item How does \(\widehat P_{\mathrm{ej}}(T)\) change with \(T\)?
\item Which semimajor-axis bins have the largest ejection fractions?
\item What changes when \(M_J\) is scaled upward?
\item Which losses are close encounters rather than clean ejections?
\item How sensitive are the results to \(\Delta t\)?
\item Are apparent resonance structures stable under a different random seed?
\end{enumerate}

\section{Associated Demos}
\label{sec:chaos-associated-demos}

\begin{itemize}
\item \path{demos/python/standard_map.py}
\item \path{demos/python/hamiltonian_pendulum.py}
\item \path{demos/python/asteroid_ejection_probability.py}
\item \path{demos/mathematica/AsteroidEjectionProbability.wl}
\end{itemize}

\chapter[Fluid Mechanics]{Fluid Mechanics: Balance Laws, Navier-Stokes, And Stability}
\label{chap:fluid}

\section{What Is A Fluid?}
\label{sec:what-is-fluid}

Fluid mechanics is continuum mechanics with a special constitutive idea: the
material does not resist a static shear deformation. A solid can sit in a
sheared shape and store elastic energy; a simple fluid responds to shear through
stress proportional to the rate of deformation. This chapter therefore treats
fluids on their own terms. Elastic and deforming bodies are separated into the
next chapter.

The basic fields are functions on physical space:
\[
  \rho(x,t),\qquad u(x,t),\qquad p(x,t),
\]
where \(\rho\) is mass density, \(u\) is velocity, and \(p\) is pressure. A
fluid particle can still be labelled by a material coordinate \(X\), with
\[
  x=\varphi(X,t),\qquad u(\varphi(X,t),t)=\partial_t\varphi(X,t).
\]
The Eulerian view keeps \(x\) fixed and watches fluid pass through; the
Lagrangian view follows the particle labelled by \(X\).

\begin{table}[htbp]
\centering
\caption{Fluid-mechanics notation. The symbol \(p\) denotes pressure in this
chapter, not canonical momentum.}
\label{tab:fluid-notation}
\small
\begin{tabular}{p{0.18\textwidth}p{0.48\textwidth}p{0.25\textwidth}}
\toprule
Symbol & Meaning & Type\\
\midrule
\(X\) & material label of a fluid particle & coordinate in reference region\\
\(\varphi(X,t)\) & flow map & material-to-spatial map\\
\(F^i{}_A\) & \(\partial\varphi^i/\partial X^A\) & deformation gradient of flow map\\
\(J\) & \(\det F\) & local volume ratio\\
\(x\) & Eulerian spatial position & point in physical space\\
\(\rho(x,t)\) & mass density & scalar field\\
\(u(x,t)\) & velocity & vector field\\
\(p(x,t)\) & pressure & scalar field\\
\(V(t)\) & material volume & moving region in space\\
\(n\) & outward unit normal to a surface & vector field on boundary\\
\(b(x,t)\) & body force per unit mass & vector field\\
\(t(n)\) & surface traction on normal \(n\) & vector field\\
\(\sigma\) & Cauchy stress tensor & \(3\times3\) tensor field\\
\(D u/Dt\) & material acceleration & \(\partial_tu+u\cdot\nabla u\)\\
\(D\) & rate-of-strain tensor & \(\frac12(\nabla u+\nabla u^T)\)\\
\(W\) & antisymmetric spin tensor & \(\frac12(\nabla u-\nabla u^T)\)\\
\(\omega\) & vorticity & \(\nabla\times u\)\\
\(\mu\) & dynamic viscosity & \(ML^{-1}T^{-1}\)\\
\(\nu\) & kinematic viscosity & \(\mu/\rho\)\\
\(\lambda\) & second viscosity coefficient & stress coefficient\\
\(U,L\) & characteristic velocity and length & nondimensionalization scales\\
\(\mathrm{Re}\) & Reynolds number & \(UL/\nu\)\\
\(\Omega\) & fixed spatial domain in energy estimates & subset of \(\R^3\)\\
\(\gamma(t)\) & closed material loop & curve transported by the flow\\
\(h,\Phi\) & specific enthalpy and conservative body-force potential & Kelvin theorem\\
\(V(y),G\) & laminar profile and pressure-gradient magnitude & shear-flow sections\\
\(\alpha,\sigma,\phi(y)\) & wave number, growth rate, normal-mode amplitude & Orr-Sommerfeld section\\
\(\Gamma_i\) & point-vortex circulation & vortex strength\\
\bottomrule
\end{tabular}
\end{table}

\begin{assumption}[Continuum fluid]
The molecular scale is much smaller than the length scales being modelled, so
density, velocity, pressure, and stress are treated as smooth fields. Across
shocks, contact discontinuities, or interfaces, the differential equations must
be replaced by weak forms and jump conditions.
\end{assumption}

\paragraph{Roadmap.}
The derivation is deliberately staged. First we learn how to differentiate
quantities attached to moving fluid parcels. Then we convert integral balance
laws into local partial differential equations for mass and momentum. Only
after those kinematic and balance statements are in place do we introduce the
Newtonian constitutive law that turns Cauchy's momentum equation into the
Navier-Stokes equation. The later sections use that equation in three ways:
exact laminar solutions, stability analysis, and Hamiltonian point-vortex
reduction.

The staging matters because the Navier-Stokes equation is not a single
assumption. Its inertial term is kinematics, its pressure and stress terms come
from momentum balance, and its viscous term comes from the Newtonian
constitutive law. Separating these ingredients makes it clear which part of the
model changes when one studies inviscid flow, non-Newtonian fluids, elasticity,
or low-Reynolds-number swimming.

\paragraph{Derivation audit trail.}
Every continuum equation below should be read with its conservation statement
attached. The material derivative comes from following one labelled parcel.
The Reynolds transport theorem converts derivatives of moving integrals into
Eulerian fields. Mass conservation comes from conserving mass in every material
volume. Cauchy's equation comes from momentum balance and the existence of a
traction stress tensor. Navier-Stokes appears only after the Newtonian stress
law is imposed. This audit trail is the best protection against treating
\(\rho Du/Dt=-\nabla p+\mu\Delta u\) as an isolated formula rather than the
endpoint of several modelling choices.

\section{Material Derivative}
\label{sec:material-derivative}

If \(f(x,t)\) is a scalar or vector field, its value along a particle path
\(x(t)=\varphi(X,t)\) changes as
\[
  \frac{\dd}{\dd t}f(x(t),t)
  =
  \partial_t f(x(t),t)+\dot x(t)\cdot\nabla f(x(t),t).
\]
Since \(\dot x=u(x(t),t)\), the material derivative is
\[
  \boxed{
  \frac{D f}{D t}
  =
  \partial_t f+u\cdot\nabla f.
  }
\]
For the velocity field itself,
\[
  \frac{D u}{D t}
  =
  \partial_tu+(u\cdot\nabla)u
\]
is the particle acceleration. The nonlinear term is kinematic: it appears even
before any force law is chosen, because the velocity field is sampled along a
moving trajectory.

\section{Reynolds Transport Theorem}
\label{sec:reynolds-transport}

Let \(V(t)\) be a material volume: it moves with the fluid, so its boundary
velocity is \(u\). For a scalar field \(f\), we want the time derivative of
\[
  I(t)=\int_{V(t)} f(x,t)\,\dd^3x.
\]
There are two sources of change. The value of \(f\) can change along each
particle path, and the volume element occupied by neighboring particles can
expand or contract. The Reynolds transport theorem is exactly the formula that
keeps both effects.

\begin{figure}[htbp]
\centering
\begin{tikzpicture}[scale=1.0, >=Latex]
  \draw[->] (-0.4,0) -- (6.8,0) node[right] {\(x^1\)};
  \draw[->] (0,-0.4) -- (0,4.0) node[above] {\(x^2\)};
  \draw[blue!70!black, thick, fill=blue!5]
    (1.0,0.9) .. controls (1.8,0.35) and (3.2,0.55) .. (3.85,1.15)
    .. controls (4.6,1.85) and (4.15,3.0) .. (3.0,3.2)
    .. controls (1.75,3.45) and (0.75,2.45) .. (1.0,0.9);
  \draw[red!75!black, thick, ->] (1.3,1.0) -- (1.75,0.55);
  \draw[red!75!black, thick, ->] (3.9,1.45) -- (4.5,1.55);
  \draw[red!75!black, thick, ->] (2.95,3.15) -- (3.15,3.75);
  \draw[red!75!black, thick, ->] (0.9,2.4) -- (0.45,2.75);
  \node[blue!70!black] at (2.45,2.0) {\(V(t)\)};
  \node[red!75!black] at (5.25,2.85) {boundary velocity \(u\)};
  \draw[red!75!black, ->] (4.75,2.65) -- (3.25,3.45);
  \node[align=center] at (5.15,0.75)
    {material points stay\\inside the moving volume};
\end{tikzpicture}
\caption{A material volume is carried by the fluid velocity. Its boundary
velocity is \(u\), so the time derivative of an integral over \(V(t)\) must
include both the material change of the integrand and the volume change
measured by \(\nabla\cdot u\).}
\label{fig:reynolds-material-volume}
\end{figure}

\begin{derivation}[Transport through a moving volume]
Write \(V(t)=\varphi(U,t)\) for a fixed material region \(U\). Let
\[
  F^i{}_A=\frac{\partial\varphi^i}{\partial X^A},
  \qquad
  J=\det F.
\]
Changing variables gives
\[
  I(t)=\int_U f(\varphi(X,t),t)J(X,t)\,\dd^3X.
\]
Differentiate:
\[
  \dot I
  =
  \int_U
  \left(
    \frac{D f}{D t}J+f\dot J
  \right)\dd^3X.
\]
The Jacobian identity is
\[
  \dot J=J\nabla\cdot u.
\]
One way to see it is Jacobi's formula:
\[
  \dot J=J\,\tr(F^{-1}\dot F),
\]
and \(\dot F^i{}_A=\partial_Au^i(\varphi,t)\), so
\(\tr(F^{-1}\dot F)=\partial_i u^i\). Returning to spatial variables,
\[
  \boxed{
  \frac{\dd}{\dd t}\int_{V(t)} f\,\dd^3x
  =
  \int_{V(t)}
  \left(
    \frac{D f}{D t}+f\nabla\cdot u
  \right)\dd^3x
  =
  \int_{V(t)}
  \left(
    \partial_t f+\nabla\cdot(fu)
  \right)\dd^3x.
  }
\]
\end{derivation}

The theorem separates two effects: the field \(f\) changes along particles, and
the material volume expands or contracts.

\section{Conservation Of Mass}
\label{sec:mass-conservation}

Mass in a material volume is constant:
\[
  \frac{\dd}{\dd t}\int_{V(t)}\rho\,\dd^3x=0.
\]
Using the transport theorem,
\[
  \int_{V(t)}
  \left(
    \partial_t\rho+\nabla\cdot(\rho u)
  \right)\dd^3x=0.
\]
Because this holds for every material volume,
\[
  \boxed{
  \partial_t\rho+\nabla\cdot(\rho u)=0.
  }
\]
Equivalently,
\[
  \frac{D\rho}{Dt}+\rho\nabla\cdot u=0.
\]
For constant-density incompressible flow,
\[
  \boxed{\nabla\cdot u=0.}
\]

\section{Cauchy Momentum Balance}
\label{sec:cauchy-momentum}

Newton's second law for a material volume \(V(t)\) is
\[
  \frac{\dd}{\dd t}\int_{V(t)}\rho u\,\dd^3x
  =
  \int_{V(t)}\rho b\,\dd^3x
  +
  \int_{\partial V(t)} t(n)\,\dd S,
\]
where \(b\) is body force per unit mass and \(t(n)\) is surface traction on a
surface whose outward unit normal is \(n\).

\begin{assumption}[Cauchy stress hypothesis]
The traction is local and linear in the normal:
\[
  t(n)=\sigma n,
\]
where \(\sigma_{ij}(x,t)\) is the Cauchy stress tensor. Balance of angular
momentum, in the absence of couple stresses, implies
\[
  \sigma=\sigma^T.
\]
\end{assumption}

\begin{derivation}[Local momentum equation]
Apply the transport theorem to \(f=\rho u\). Using mass conservation, the
left-hand side becomes
\[
\begin{aligned}
  \frac{\dd}{\dd t}\int_{V(t)}\rho u\,\dd^3x
  &=
  \int_{V(t)}
  \left(
    \partial_t(\rho u)+\nabla\cdot(\rho u\otimes u)
  \right)\dd^3x\\
  &=
  \int_{V(t)}
  \rho\left(\partial_tu+u\cdot\nabla u\right)\dd^3x.
\end{aligned}
\]
The traction term is converted by the divergence theorem:
\[
  \int_{\partial V(t)}\sigma n\,\dd S
  =
  \int_{V(t)}\nabla\cdot\sigma\,\dd^3x.
\]
Therefore
\[
  \int_{V(t)}
  \left[
    \rho\frac{D u}{D t}
    -
    \nabla\cdot\sigma
    -
    \rho b
  \right]\dd^3x=0.
\]
Since \(V(t)\) is arbitrary,
\[
  \boxed{
  \rho\frac{D u}{D t}=\nabla\cdot\sigma+\rho b.
  }
\]
\end{derivation}

This equation is not yet a closed fluid equation. It is a balance law. Closure
comes from a constitutive relation for \(\sigma\).

\section{Newtonian Stress And The Navier-Stokes Equation}
\label{sec:navier-stokes-derivation}

The balance laws so far are universal for continua, but they are not yet a
closed fluid model: the stress tensor \(\sigma\) is still unknown. A
constitutive law supplies the missing relation between stress and motion. For a
simple Newtonian fluid, the stress may depend on the instantaneous rate of
deformation but not on the accumulated shear history. This is the point at
which fluids separate from elastic solids.

The stress of a simple Newtonian fluid is assumed to depend locally and
linearly on the instantaneous rate of deformation, not on the deformation
itself. The velocity gradient splits into symmetric and antisymmetric parts:
\[
  \nabla u=D+W,
  \qquad
  D_{ij}=\frac12(\partial_i u_j+\partial_j u_i),
  \qquad
  W_{ij}=\frac12(\partial_i u_j-\partial_j u_i).
\]
The antisymmetric part \(W\) represents local rigid rotation. It should not
produce viscous stress in an isotropic simple fluid. This statement is a
constitutive assumption with three ingredients:
\[
\begin{array}{ll}
\text{locality:} & \sigma(x,t) \text{ depends on } u \text{ through }
  \nabla u(x,t),\\
\text{objectivity:} & \text{superposed rigid rotations do not create stress},\\
\text{isotropy:} & \text{there is no preferred material direction}.
\end{array}
\]
Objectivity removes \(W\). Isotropy says that a linear map from the symmetric
tensor \(D\) to a symmetric stress can only use the deviatoric part of \(D\)
and its trace. The most general isotropic linear relation is therefore
\[
  \boxed{
  \sigma=-pI+2\mu D+\lambda(\nabla\cdot u)I.
  }
\]
Here \(p\) is pressure, \(\mu\geq0\) is shear viscosity, and \(\lambda\) is the
second viscosity coefficient. For constant \(\mu,\lambda\),
\[
\begin{aligned}
  (\nabla\cdot\sigma)_i
  &=
  -\partial_i p
  +
  \mu\partial_j(\partial_i u_j+\partial_j u_i)
  +
  \lambda\partial_i(\partial_k u_k)\\
  &=
  -\partial_i p
  +
  \mu\Delta u_i
  +
  (\lambda+\mu)\partial_i(\nabla\cdot u).
\end{aligned}
\]
Thus the compressible Navier-Stokes momentum equation is
\[
  \boxed{
  \rho
  \left(
    \partial_tu+u\cdot\nabla u
  \right)
  =
  -\nabla p+\mu\Delta u+(\lambda+\mu)\nabla(\nabla\cdot u)+\rho b.
  }
\]
Together with mass conservation, an equation of state, and an energy or
temperature equation when needed, this describes a Newtonian compressible fluid.

For incompressible flow with constant \(\rho\),
\[
  \nabla\cdot u=0,
\]
and the pressure is best understood as the Lagrange multiplier enforcing this
constraint. The equations reduce to
\[
  \boxed{
  \begin{aligned}
  \rho(\partial_tu+u\cdot\nabla u)&=-\nabla p+\mu\Delta u+\rho b,\\
  \nabla\cdot u&=0.
  \end{aligned}
  }
\]

Read the incompressible equations term by term. The material acceleration
\(\partial_tu+u\cdot\nabla u\) is the acceleration of a moving fluid particle.
The pressure gradient enforces the volume-preserving constraint and transmits
normal stress. The Laplacian \(\nu\Delta u\) diffuses momentum. The body force
adds external acceleration. This interpretation prevents a common mistake:
pressure is not specified independently in incompressible flow; it is solved
simultaneously with \(u\) so that \(\nabla\cdot u=0\) remains true.
Taking the divergence of the incompressible momentum equation makes this
constraint role explicit:
\[
  -\Delta p
  =
  \rho\,\nabla\cdot(u\cdot\nabla u)-\rho\,\nabla\cdot b,
\]
with boundary conditions determined by the normal component of the momentum
equation at the wall. The pressure is therefore a nonlocal field: it adjusts
instantaneously, within the incompressible model, to project the acceleration
back onto divergence-free velocity fields.

\section{Nondimensionalization And Reynolds Number}
\label{sec:reynolds-number}

Let \(U\) be a typical speed, \(L\) a typical length, and define
\[
  x=Lx',\qquad
  t=\frac{L}{U}t',\qquad
  u=Uu',\qquad
  p=\rho U^2p'.
\]
Dropping primes, the incompressible Navier-Stokes equation without body force
becomes
\[
  \partial_tu+u\cdot\nabla u
  =
  -\nabla p+\frac1{\mathrm{Re}}\Delta u,
  \qquad
  \nabla\cdot u=0,
\]
where
\[
  \boxed{\mathrm{Re}=\frac{UL}{\nu}},\qquad \nu=\frac{\mu}{\rho}.
\]
The Reynolds number compares inertial transport to viscous diffusion. Small
\(\mathrm{Re}\) flow is viscously dominated and often reversible in time after
the velocity is reversed. Large \(\mathrm{Re}\) flow can carry disturbances
long distances before viscosity damps them.

\section{Energy Identity And Viscous Dissipation}
\label{sec:fluid-energy-identity}

For a smooth incompressible solution in a fixed domain \(\Omega\), assume either
periodic boundary conditions or no-slip walls \(u=0\) on \(\partial\Omega\).
Dot the Navier-Stokes equation with \(u\) and integrate:
\[
  \rho\int_\Omega u\cdot\partial_tu\,\dd^3x
  +
  \rho\int_\Omega u\cdot(u\cdot\nabla u)\,\dd^3x
  =
  -\int_\Omega u\cdot\nabla p\,\dd^3x
  +
  \mu\int_\Omega u\cdot\Delta u\,\dd^3x.
\]
The nonlinear term is a divergence:
\[
  u\cdot(u\cdot\nabla u)=u\cdot\nabla\left(\frac12\norm u^2\right)
  =
  \nabla\cdot\left(\frac12\norm u^2u\right),
\]
because \(\nabla\cdot u=0\). The pressure term is also a divergence:
\[
  u\cdot\nabla p=\nabla\cdot(pu).
\]
The boundary conditions remove both boundary fluxes. Integration by parts gives
\[
  \int_\Omega u\cdot\Delta u\,\dd^3x
  =
  -\int_\Omega |\nabla u|^2\,\dd^3x.
\]
Hence
\[
  \boxed{
  \frac{\dd}{\dd t}
  \int_\Omega \frac12\rho\norm u^2\,\dd^3x
  =
  -\mu\int_\Omega |\nabla u|^2\,\dd^3x\leq0.
  }
\]
Viscosity irreversibly converts organized kinetic energy into heat. This
identity is also the first stability estimate for laminar flows: disturbances
must draw energy from mean shear fast enough to beat viscous dissipation.

\section{Vorticity And Circulation}
\label{sec:vorticity-circulation}

The vorticity is
\[
  \omega=\nabla\times u.
\]
Taking the curl of the incompressible Navier-Stokes equation gives
\[
  \boxed{
  \partial_t\omega+u\cdot\nabla\omega
  =
  \omega\cdot\nabla u+\nu\Delta\omega.
  }
\]
The term \(\omega\cdot\nabla u\) stretches and tilts vortex lines in three
dimensions. In two dimensions, vorticity is perpendicular to the plane and the
stretching term vanishes:
\[
  \partial_t\omega+u\cdot\nabla\omega=\nu\Delta\omega.
\]
This contrast is central. Two-dimensional inviscid vorticity is transported
like a scalar dye, which creates many conserved quantities. Three-dimensional
vorticity can be stretched by the flow itself, allowing vortex intensification
and making the analysis of Navier-Stokes much harder.
For \(\nu=0\), vorticity is transported by the flow.

\begin{proposition}[Kelvin circulation theorem]
Let \(\gamma(t)\) be a closed material loop moving with an inviscid barotropic
fluid under conservative body forces. Then
\[
  \frac{\dd}{\dd t}\oint_{\gamma(t)}u\cdot \dd x=0.
\]
\end{proposition}

\begin{proof}
Parametrize the loop by \(s\mapsto x(s,t)\), with
\[
  \partial_t x(s,t)=u(x(s,t),t).
\]
Then
\[
  \frac{\dd}{\dd t}\oint u\cdot \partial_sx\,\dd s
  =
  \oint
  \left(
    \frac{D u}{D t}\cdot \partial_sx
    +
    u\cdot\partial_su
  \right)\dd s.
\]
For inviscid barotropic flow with conservative potential \(\Phi\),
\[
  \frac{D u}{D t}=-\nabla h-\nabla\Phi,
\]
where \(h\) is specific enthalpy. The integrand becomes
\[
  \partial_s\left(
    -h-\Phi+\frac12\norm u^2
  \right).
\]
The integral of a single-valued total derivative around a closed loop is zero.
\end{proof}

\section{Laminar Flow Solutions}
\label{sec:laminar-solutions}

Laminar solutions are exact, organized solutions of Navier-Stokes. They are
important not because all flows are laminar, but because stability theory asks
which organized solutions persist under disturbance.
Each solution follows from three choices: geometry, boundary conditions, and
pressure gradient or wall motion. Changing any one of these choices changes the
profile and the instability problem.

\subsection{Plane Couette Flow}
\label{subsec:plane-couette}

Two infinite plates are separated by distance \(h\). The lower plate is fixed;
the upper plate moves with velocity \(Ue_x\). Seek a steady parallel flow
\[
  u=(V(y),0,0),\qquad p=\text{constant}.
\]
The incompressibility condition is automatic. The \(x\)-momentum equation is
\[
  0=\mu V''(y),
\]
with boundary conditions \(V(0)=0\), \(V(h)=U\). Thus
\[
  \boxed{V(y)=U\frac{y}{h}.}
\]
The shear stress is constant:
\[
  \sigma_{xy}=\mu V'(y)=\mu\frac{U}{h}.
\]

\subsection{Plane Poiseuille Flow}
\label{subsec:plane-poiseuille}

Now both plates are fixed and a pressure gradient drives the flow. Let
\[
  u=(V(y),0,0),
  \qquad
  -\frac{\dd p}{\dd x}=G>0,
  \qquad
  0<y<h.
\]
The steady \(x\)-momentum equation is
\[
  0=-\frac{\dd p}{\dd x}+\mu V''(y)=G+\mu V''(y),
\]
so
\[
  V''(y)=-\frac{G}{\mu}.
\]
With no-slip conditions \(V(0)=V(h)=0\),
\[
  \boxed{
  V(y)=\frac{G}{2\mu}y(h-y).
  }
\]
The maximum speed occurs at \(y=h/2\):
\[
  V_{\max}=\frac{Gh^2}{8\mu}.
\]
The volume flux per unit span is
\[
  Q=\int_0^h V(y)\,\dd y
  =
  \frac{Gh^3}{12\mu}.
\]

\subsection{Hagen-Poiseuille Pipe Flow}
\label{subsec:hagen-poiseuille}

For a circular pipe of radius \(a\), use cylindrical radius \(r\) and seek
\[
  u=V(r)e_z,\qquad -\frac{\dd p}{\dd z}=G>0.
\]
The axial equation is
\[
  0=G+\mu\frac1r\frac{\dd}{\dd r}
  \left(
    r\frac{\dd V}{\dd r}
  \right).
\]
Regularity at \(r=0\) and no slip at \(r=a\) give
\[
  \boxed{
  V(r)=\frac{G}{4\mu}(a^2-r^2).
  }
\]
The flux is
\[
  Q=2\pi\int_0^a V(r)r\,\dd r
  =
  \boxed{\frac{\pi a^4G}{8\mu}}.
\]
The fourth power in \(a\) is one of the most important practical consequences
of viscous laminar flow.

\begin{figure}[htbp]
\centering
\begin{tikzpicture}[scale=1.0, >=Latex]
  \draw[->] (-0.25,0) -- (-0.25,3.25) node[above] {\(y\)};
  \draw[->] (-0.25,0) -- (1.2,0) node[right] {\(x\)};
  \draw[thick] (0,0) -- (3.0,0);
  \draw[thick] (0,3.0) -- (3.0,3.0);
  \draw[->, thick, black!70] (2.1,3.18) -- (2.8,3.18) node[right] {wall speed \(U\)};
  \foreach \y/\len in {0.45/0.25,0.95/0.55,1.45/0.85,1.95/1.15,2.45/1.45}
    \draw[blue, thick, ->] (0.35,\y) -- ++(\len,0);
  \node[blue] at (1.5,3.45) {Couette};
  \node at (1.5,-0.35) {\(V(0)=0,\ V(h)=U\)};

  \begin{scope}[xshift=4.4cm]
    \draw[thick] (0,0) -- (3.0,0);
    \draw[thick] (0,3.0) -- (3.0,3.0);
    \draw[->, thick, black!70] (0.25,3.25) -- (1.35,3.25)
      node[right] {\(-\partial_xp=G\)};
    \foreach \y/\len in {0.45/0.45,0.95/1.05,1.50/1.45,2.05/1.05,2.55/0.45}
      \draw[red!75!black, thick, ->] (0.35,\y) -- ++(\len,0);
    \node[red!75!black] at (1.5,3.45) {Poiseuille};
    \node at (1.5,-0.35) {\(V(0)=V(h)=0\)};
  \end{scope}
\end{tikzpicture}
\caption{Couette flow is linear because it is driven by wall motion. Poiseuille
flow is parabolic because it balances pressure forcing against viscous
diffusion of momentum. The arrows are spatial velocity vectors, not a graph
with velocity on the vertical axis.}
\label{fig:couette-poiseuille-profiles}
\end{figure}

\section{Linear Stability Of Laminar Shear Flow}
\label{sec:linear-stability-shear}

Let a steady base flow be
\[
  U(y)e_x.
\]
Write a small perturbation as
\[
  u=U(y)e_x+\tilde u,\qquad
  p=P(x)+\tilde p,
\]
and retain only terms linear in \(\tilde u\). In nondimensional variables,
\[
  \boxed{
  \partial_t\tilde u
  +
  U\partial_x\tilde u
  +
  \tilde v\,U'(y)e_x
  =
  -\nabla\tilde p
  +
  \frac1{\mathrm{Re}}\Delta\tilde u,
  \qquad
  \nabla\cdot\tilde u=0.
  }
\]
The term \(\tilde v\,U'e_x\) is the mechanism by which cross-stream velocity
tilts fluid into faster or slower parts of the mean shear and extracts energy
from the base flow.

\begin{figure}[htbp]
\centering
\begin{tikzpicture}[scale=1.0, >=Latex]
  \draw[->] (-0.25,0) -- (-0.25,3.4) node[above] {\(y\)};
  \draw[->] (-0.25,0) -- (5.9,0) node[right] {\(x\)};
  \draw[thick] (0,0.35) -- (5.3,0.35);
  \draw[thick] (0,3.05) -- (5.3,3.05);
  \foreach \y/\len in {0.65/0.35,1.10/0.62,1.55/0.92,2.00/1.25,2.45/1.55,2.85/1.8}
    \draw[blue!70!black, thick, ->] (0.45,\y) -- ++(\len,0);
  \draw[red!75!black, thick, ->] (2.75,1.15) -- (2.75,2.05)
    node[midway,right] {\(\tilde v\)};
  \draw[red!75!black, thick, ->] (2.75,2.05) -- (3.55,2.05)
    node[right] {\(\tilde v\,U'(y)e_x\)};
  \draw[purple!75!black, thick, domain=0.3:5.0, samples=120]
    plot (\x,{1.75+0.22*sin(115*\x)});
  \node[blue!70!black] at (4.25,2.75) {base shear \(U(y)e_x\)};
  \node[purple!75!black] at (4.0,1.25) {normal-mode disturbance};
\end{tikzpicture}
\caption{Linear shear-flow stability. A wall-normal perturbation velocity
\(\tilde v\) moves fluid across the base shear; the term
\(\tilde v\,U'(y)e_x\) is the linear energy-exchange channel between the base
flow and the disturbance.}
\label{fig:shear-stability-mechanism}
\end{figure}

\subsection{Orr-Sommerfeld Equation}
\label{subsec:orr-sommerfeld}

For two-dimensional perturbations, introduce a streamfunction
\[
  \tilde u=\partial_y\psi,\qquad
  \tilde v=-\partial_x\psi.
\]
Seek normal modes
\[
  \psi(x,y,t)=\phi(y)e^{i\alpha x+\sigma t}.
\]
Taking the curl of the linearized equation removes pressure and gives the
linear vorticity equation. Since the perturbation vorticity is
\[
  \tilde\omega_z
  =
  \partial_x\tilde v-\partial_y\tilde u
  =
  -(\partial_y^2-\alpha^2)\phi\,e^{i\alpha x+\sigma t},
\]
substitution gives a fourth-order eigenvalue equation for the wall-normal
shape \(\phi(y)\):
\[
  \boxed{
  (\sigma+i\alpha U)(D^2-\alpha^2)\phi
  -
  i\alpha U''\phi
  =
  \frac1{\mathrm{Re}}(D^2-\alpha^2)^2\phi,
  }
\]
where \(D=\dd/\dd y\). No-slip and no-penetration at walls give
\[
  \phi=0,\qquad D\phi=0
\]
at the walls. This fourth-order eigenvalue problem is the
Orr-Sommerfeld equation. A base flow is linearly unstable if some eigenvalue has
\[
  \operatorname{Re}\sigma>0.
\]
Many fluid-mechanics texts write the same normal mode as
\(e^{i\alpha(x-ct)}\). With the convention used here, the complex phase speed
\(c\) and growth rate \(\sigma\) are related by
\[
  \sigma=-i\alpha c,
\]
so \(\operatorname{Re}\sigma=\alpha\,\operatorname{Im}c\) for real
\(\alpha\).

\begin{remark}[Classical facts]
Plane Poiseuille flow has a linear instability above a critical Reynolds number
about \(5772\), using centerline speed and half-channel width. Plane Couette
flow is linearly stable for all Reynolds numbers in the classical
Orr-Sommerfeld analysis, yet it can transition through finite-amplitude
disturbances. Stability is therefore not a single concept: linear eigenvalue
stability, energy stability, transient growth, and nonlinear transition answer
different questions.
\end{remark}

\subsection{Energy Stability Estimate}
\label{subsec:energy-stability}

For perturbations of a parallel flow, the kinetic energy of the disturbance
satisfies
\[
  \frac12\frac{\dd}{\dd t}\int_\Omega |\tilde u|^2\,\dd^3x
  =
  -\int_\Omega \tilde u\,\tilde v\,U'(y)\,\dd^3x
  -
  \frac1{\mathrm{Re}}
  \int_\Omega |\nabla\tilde u|^2\,\dd^3x.
\]
The second term is viscous dissipation. The first is production by mean shear.
If dissipation dominates production for all admissible perturbations, the base
flow is energy stable. If not, the estimate does not prove instability, but it
reveals the only possible energy source for perturbation growth.

\section{Point Vortices As A Hamiltonian Fluid Model}
\label{sec:point-vortices}

In two dimensions, an idealized vorticity distribution can be written as
\[
  \omega(x,y)=\sum_{i=1}^{N}\Gamma_i\delta(x-x_i,y-y_i),
\]
where \(\Gamma_i\) is circulation. The point vortex equations are
\[
\begin{aligned}
  \Gamma_i\dot x_i &= \frac{\partial H}{\partial y_i},\\
  \Gamma_i\dot y_i &= -\frac{\partial H}{\partial x_i},
\end{aligned}
\]
with Hamiltonian
\[
  H
  =
  -\frac{1}{4\pi}
  \sum_{i<j}\Gamma_i\Gamma_j
  \log\left((x_i-x_j)^2+(y_i-y_j)^2\right).
\]
This is Hamiltonian mechanics, but the symplectic form is weighted:
\[
  \omega_{\mathrm{vortex}}
  =
  \sum_i \Gamma_i\,\dd x_i\wedge\dd y_i.
\]
The weight \(\Gamma_i\) is the circulation carried by the \(i\)-th vortex. One
can derive this finite-dimensional symplectic form by reducing the ideal-fluid
Hamiltonian bracket to delta-function vorticity; here we take the reduced
point-vortex bracket as the model.

\begin{figure}[htbp]
\centering
\begin{tikzpicture}[scale=1.05, >=Latex]
  \draw[->] (-2.5,0) -- (2.7,0) node[right] {\(x\)};
  \draw[->] (0,-2.2) -- (0,2.4) node[above] {\(y\)};
  \foreach \x/\y/\g/\lab in {-0.65/-0.35/+/\Gamma,0.65/0.35/-/-\Gamma}
  {
    \draw[thick] (\x,\y) circle (0.16);
    \node at (\x,\y) {\(\g\)};
    \node[above right] at (\x,\y) {\(\lab\)};
  }
  \draw[blue, thick, ->] (-1.7,-0.35) -- (1.45,-0.35);
  \draw[blue, thick, ->] (-0.4,0.35) -- (2.75,0.35);
  \draw[dashed] (-0.65,-0.35) -- (0.65,0.35);
  \draw[<->] (-0.65,-0.62) -- (0.65,0.08) node[midway, below right] {separation};
  \node at (0,-1.75) {a vortex dipole translates by mutual induction};
\end{tikzpicture}
\caption{A vortex dipole is the simplest point-vortex motion: the two vortices
translate together with fixed separation. More complicated point-vortex
systems give finite-dimensional Hamiltonian models of ideal-fluid motion.}
\label{fig:vortex-dipole}
\end{figure}

\section{What To Take Forward}
\label{sec:fluid-take-forward}

\begin{itemize}
\item The material derivative is the bridge between following particles and
describing fields at fixed spatial points.
\item Conservation of mass and momentum are balance laws before they are
differential equations; the local PDEs come from applying the divergence theorem
to arbitrary moving volumes.
\item Navier-Stokes is Cauchy momentum balance plus the Newtonian stress law.
Changing the constitutive law changes the continuum model.
\item The Reynolds number compares inertial transport with viscous diffusion of
momentum. It is the organizing scale behind laminar solutions, instabilities,
and transition.
\item Point vortices show that Hamiltonian structure survives inside fluid
mechanics after a strong but explicit reduction.
\end{itemize}

\section{Associated Demos}
\label{sec:fluid-associated-demos}

\begin{itemize}
\item \path{demos/python/fluids_vorticity.py}, a point-vortex Hamiltonian model.
\item \path{demos/python/hamiltonian_pendulum.py}, useful for comparing
separatrix geometry with shear-flow stability.
\item Suggested extension: a numerical Orr-Sommerfeld spectrum for Couette and
Poiseuille flow.
\end{itemize}

\chapter{Elastic Bodies And Gauge Kinematics Of Deformation}
\label{chap:elastic-gauge}

\section{Why This Chapter Is Separate From Fluids}
\label{sec:elastic-vs-fluid}

Fluids and elastic bodies share continuum balance laws, but they organize
material response differently. A fluid has no preferred static shear shape: its
deviatoric stress is tied to the rate of deformation. An elastic body stores
energy in deformation itself. This chapter therefore returns to the material
map
\[
  \varphi:B\times I\to\R^3,
  \qquad
  x=\varphi(X,t),
\]
and treats \(X\in B\) as a persistent label of a material point.

The second purpose of the chapter is to introduce the gauge-theoretic
description of deforming bodies in the spirit of Shapere and Wilczek.\footnote{
A. Shapere and F. Wilczek, ``Gauge kinematics of deformable bodies,''
American Journal of Physics 57, 514--518 (1989),
\href{https://doi.org/10.1119/1.15986}{doi:10.1119/1.15986}. Their closely
related low-Reynolds-number swimming formulation appears in
\href{https://doi.org/10.1017/S002211208900025X}{doi:10.1017/S002211208900025X}.}
The central idea is simple and powerful: a deformable body has a space of
shapes, and choosing axes for each shape is a gauge choice. Cyclic changes of
shape can produce a net rotation or translation, not because of a force impulse,
but because a connection over shape space has nonzero curvature.

\begin{table}[htbp]
\centering
\caption{Elasticity, Cosserat, and Shapere-Wilczek gauge notation.}
\label{tab:elastic-gauge-notation}
\small
\begin{tabular}{p{0.18\textwidth}p{0.48\textwidth}p{0.25\textwidth}}
\toprule
Symbol & Meaning & Type\\
\midrule
\(B\) & reference body & material domain\\
\(\partial B\) & boundary of the reference body & material surface\\
\(X^A\) & material coordinate & coordinate on \(B\)\\
\(x^i\) & spatial coordinate & coordinate in \(\R^3\)\\
\(\varphi(X,t)\) & deformation map & \(B\to\R^3\)\\
\(F^i{}_A\) & deformation gradient & \(\partial_A\varphi^i\)\\
\(J\) & local volume ratio & \(\det F\)\\
\(C_{AB}\) & right Cauchy-Green tensor & \(F^i{}_AF^i{}_B\)\\
\(E_{AB}\) & Green strain & \(\frac12(C_{AB}-\delta_{AB})\)\\
\(\xi(X,t)\) & small displacement field & \(\varphi=X+\xi\)\\
\(e_{AB}\) & infinitesimal strain & symmetric gradient of \(\xi\)\\
\(W(F,X)\) & stored energy density & energy per reference volume\\
\(\rho_0(X)\) & reference mass density & mass per reference volume\\
\(P^A{}_i\) & first Piola-Kirchhoff stress & \(\partial W/\partial F^i{}_A\)\\
\(\sigma_{ij}\) & Cauchy stress & spatial stress tensor\\
\(\lambda,\mu\) & Lame moduli & elastic constants\\
\(c_P,c_S\) & longitudinal and shear wave speeds & speed\\
\(g,s,\mathcal S\) & rigid placement, shape, and shape space & gauge kinematics\\
\(\xi=g^{-1}\dot g\) & rigid velocity in moving frame & Lie-algebra element\\
\(M(s),m(s)\) & rigid and shape blocks of kinetic metric & matrices/operators\\
\(\Pi\) & momentum conjugate to rigid motion & element of \(\mathfrak g^*\)\\
\(\mathcal A\) & mechanical/gauge connection & Lie-algebra-valued one-form\\
\(\mathcal F\) & curvature of \(\mathcal A\) & Lie-algebra-valued two-form\\
\(\Gamma\) & closed loop in shape space & stroke/cycle\\
\(\mathcal P\) & path-ordering symbol & ordered exponential\\
\(R(X,t)\) & local material frame & \(\SO(3)\)-valued field\\
\(A_A,A_t\) & spatial and time connection components & \(\mathfrak{so}(3)\)-valued\\
\(\Gamma_A\) & translational strain in local frame & \(R^{-1}\partial_A\varphi\)\\
\(T^a\) & torsion two-form & dislocation density model\\
\(\kappa,B,\kappa_0\) & rod curvature, bending stiffness, intrinsic curvature & rod variables\\
\bottomrule
\end{tabular}
\end{table}

\paragraph{Roadmap.}
The chapter has two intertwined narratives. The first is classical elasticity:
deformation gradients measure local stretching, compatibility distinguishes
true displacement fields from incompatible strain data, and a hyperelastic
action gives elastodynamic equations and stress. The second is gauge
kinematics: once a body can change shape, choosing body axes becomes a
shape-dependent convention, and cyclic shape change can produce rigid motion
through holonomy. The Cosserat and defect sections show how these two
narratives meet in local frame fields.

The guiding distinction is between intrinsic deformation and frame description.
Quantities such as \(C=F^TF\) and the Green strain do not change under a rigid
rotation of the observer. Quantities such as local frame components and
connection coefficients do change, but in a controlled gauge-covariant way. The
goal is not to avoid frames; it is to know exactly which statements depend on
them.

\paragraph{Vocabulary bridge.}
In the elasticity sections, ``configuration'' means a deformation map
\(\varphi:B\to\R^3\). In the gauge sections, the same physical configuration is
split into a rigid placement \(g\) and a shape \(s\). This split is useful, but
not unique: changing the convention for body axes changes \(g\) and changes the
connection coefficients \(\mathcal A_\alpha\). The observable statement is the
reconstructed rigid motion around a closed shape cycle, especially its
curvature-controlled small-loop limit. Thus the gauge language is not an extra
force law; it is a precise bookkeeping system for frame-dependent descriptions
of deformation.

\section{Deformation Gradient And Strain}
\label{sec:deformation-gradient-strain}

Elasticity begins by comparing nearby material points before and after
deformation. The object that does this comparison is not the displacement
itself but its derivative: two deformations that differ by a rigid translation
have the same local strain. This is why the deformation gradient \(F\), and the
metric-like tensor \(C=F^TF\), are the natural first objects.

The deformation gradient is the local linearization of the deformation map:
\[
  F^i{}_A=\frac{\partial \varphi^i}{\partial X^A}.
\]
If \(\delta X\) is a small material vector, then its spatial image is
\[
  \delta x^i=F^i{}_A\delta X^A.
\]
The right Cauchy-Green tensor measures squared lengths after deformation:
\[
  |\delta x|^2
  =
  C_{AB}\delta X^A\delta X^B,
  \qquad
  C_{AB}=F^i{}_AF^i{}_B.
\]
The Green strain is
\[
  E_{AB}=\frac12(C_{AB}-\delta_{AB}).
\]
It vanishes for every rigid motion
\[
  \varphi(X)=a+QX,\qquad Q\in\SO(3),
\]
because \(F=Q\) and \(Q^TQ=I\). This invariance is why \(E\), not \(F-I\), is a
geometric strain measure for finite deformations.

For small displacement
\[
  \varphi(X,t)=X+\xi(X,t),
\]
one has
\[
  F^i{}_A=\delta^i_A+\partial_A\xi^i.
\]
Keeping only first-order displacement gradients,
\[
  E_{AB}
  =
  \frac12(\partial_A\xi_B+\partial_B\xi_A)
  +
  O(|\nabla\xi|^2).
\]
The infinitesimal strain is therefore
\[
  \boxed{
  e_{AB}=\frac12(\partial_A\xi_B+\partial_B\xi_A).
  }
\]

\section{Compatibility}
\label{sec:elastic-compatibility}

Not every matrix field is a deformation gradient. If
\[
  F^i{}_A=\partial_A\varphi^i
\]
for a smooth deformation, then mixed partial derivatives commute:
\[
  \boxed{
  \partial_AF^i{}_B-\partial_BF^i{}_A=0.
  }
\]
On a simply connected body, this curl-free condition is also enough to recover
\(\varphi\) from \(F\), up to a rigid translation.

For small strain, the compatibility question is subtler because \(e\) is only
the symmetric part of \(\nabla\xi\). Given a symmetric tensor field \(e_{ij}\),
we ask whether there exists a displacement \(\xi\) such that
\[
  e_{ij}=\frac12(\partial_i\xi_j+\partial_j\xi_i).
\]
The answer is controlled by the Saint-Venant tensor:
\[
  \boxed{
  \epsilon_{ik\ell}\epsilon_{jmn}
  \partial_k\partial_m e_{\ell n}=0.
  }
\]
This is often written
\[
  \operatorname{curl}\operatorname{curl} e=0.
\]
When the condition fails, the local strain field cannot be assembled into a
single-valued displacement without cuts, defects, residual stress, or a
non-Euclidean reference geometry.

\begin{remark}[Geometric meaning]
The compatibility equations are the linearized version of flatness. A strain
field defines a metric \(g_{ij}=\delta_{ij}+2e_{ij}\). If this metric is the
pullback of the Euclidean metric by an embedding, its Riemann curvature must
vanish. Linearizing that curvature gives Saint-Venant compatibility.
\end{remark}

\section{Hyperelastic Action And Piola Stress}
\label{sec:hyperelastic-action}

A hyperelastic material has stored energy density \(W(F,X)\) per reference
volume. Let \(\rho_0(X)\) be reference mass density. The action is
\[
  S[\varphi]
  =
  \int_{t_0}^{t_1}\int_B
  \left(
    \frac12\rho_0|\partial_t\varphi|^2-W(F,X)
  \right)\dd^3X\,\dd t.
\]

\begin{assumption}[Hyperelasticity]
The stress is derived from stored energy:
\[
  P^A{}_i=\frac{\partial W}{\partial F^i{}_A}.
\]
This excludes plasticity, viscosity, fracture, and history dependence. Those
effects require additional internal variables or dissipation laws.
\end{assumption}

\begin{derivation}[Elastodynamic equation]
Let \(\delta\varphi=\eta\), with \(\eta=0\) at the time endpoints. Since
\[
  \delta F^i{}_A=\partial_A\eta^i,
\]
the first variation is
\[
\begin{aligned}
  \delta S
  &=
  \int_{t_0}^{t_1}\int_B
  \left(
    \rho_0\dot\varphi_i\dot\eta_i
    -
    P^A{}_i\partial_A\eta^i
  \right)\dd^3X\,\dd t\\
  &=
  \int_{t_0}^{t_1}\int_B
  \left(
    -\rho_0\ddot\varphi_i
    +
    \partial_A P^A{}_i
  \right)\eta^i\dd^3X\,\dd t
  -
  \int_{t_0}^{t_1}\int_{\partial B}
  P^A{}_iN_A\eta^i\,\dd S\,\dd t.
\end{aligned}
\]
For arbitrary interior variations,
\[
  \boxed{
  \rho_0\ddot\varphi_i=\partial_A P^A{}_i.
  }
\]
If a boundary traction \(t_i\) per reference area is prescribed, the natural
boundary condition is
\[
  P^A{}_iN_A=t_i.
\]
\end{derivation}

\section{Spatial Stress And Frame Indifference}
\label{sec:spatial-stress-frame-indifference}

The Cauchy stress \(\sigma\) is the physical stress acting across spatial
surfaces. It is related to Piola stress by
\[
  \boxed{
  \sigma_{ij}
  =
  \frac1J P^A{}_iF^j{}_A.
  }
\]
This relation is the Piola transform: it converts force per reference area to
force per current spatial area.

A stored energy is frame indifferent if a superposed rigid rotation does not
change it:
\[
  W(QF,X)=W(F,X)\qquad\text{for all }Q\in\SO(3).
\]
Let \(Q(\eps)=I+\eps A+O(\eps^2)\), with \(A^T=-A\). Differentiating
frame indifference at \(\eps=0\) gives
\[
  0=
  P^A{}_iA^i{}_jF^j{}_A.
\]
Since \(A\) is any skew matrix, the tensor \(PF^T\) must be symmetric:
\[
  P^A{}_iF^j{}_A=P^A{}_jF^i{}_A.
\]
Therefore \(\sigma=\sigma^T\). Thus symmetry of stress follows from rotational
invariance of stored energy, matching angular-momentum balance.

\section{Linear Isotropic Elasticity}
\label{sec:linear-isotropic-elasticity}

For small strain in an isotropic homogeneous solid, the quadratic energy density
is
\[
  W(e)=\frac{\lambda}{2}(\tr e)^2+\mu e_{ij}e_{ij},
\]
where \(\lambda,\mu\) are Lame moduli. The stress is
\[
  \sigma_{ij}
  =
  \frac{\partial W}{\partial e_{ij}}
  =
  \lambda(\tr e)\delta_{ij}+2\mu e_{ij}.
\]
The linearized equation of motion is
\[
  \rho_0\partial_t^2\xi_i=\partial_j\sigma_{ij}+\rho_0 b_i.
\]
Substituting the stress-strain relation,
\[
\begin{aligned}
  \partial_j\sigma_{ij}
  &=
  \lambda\partial_i(\partial_k\xi_k)
  +
  \mu\partial_j(\partial_i\xi_j+\partial_j\xi_i)\\
  &=
  (\lambda+\mu)\partial_i(\nabla\cdot\xi)
  +
  \mu\Delta\xi_i.
\end{aligned}
\]
Hence
\[
  \boxed{
  \rho_0\partial_t^2\xi
  =
  (\lambda+\mu)\nabla(\nabla\cdot\xi)
  +
  \mu\Delta\xi
  +
  \rho_0 b.
  }
\]
Using
\[
  \Delta\xi=\nabla(\nabla\cdot\xi)-\nabla\times\nabla\times\xi,
\]
one sees two wave families:
\[
  c_P=\sqrt{\frac{\lambda+2\mu}{\rho_0}},
  \qquad
  c_S=\sqrt{\frac{\mu}{\rho_0}}.
\]
Longitudinal waves change volume; shear waves change shape at nearly fixed
volume.

\section{Gauge Kinematics Of Deformable Bodies}
\label{sec:gauge-kinematics-deformable-bodies}

The preceding sections describe how an elastic body deforms relative to a
reference body. Gauge kinematics asks a different but compatible question: if
the body changes shape internally, how much of its observed motion is a genuine
rigid displacement, and how much is a convention for choosing body axes at each
shape? This is the same distinction that appears in gauge theory: local
descriptions depend on a frame choice, while holonomy around a closed loop is
observable.

The Wilczek-Shapere point of view begins with a finite deformable body moving in
space. Separate its configuration into:
\[
  g\in G,\qquad s\in\mathcal S,
\]
where \(G\) is the rigid-motion group, usually \(\SE(3)\) or \(\SO(3)\), and
\(\mathcal S\) is shape space. A point of the full configuration space is a
rigid placement \(g\) together with an internal shape \(s\).

The subtlety is that the split into \(g\) and \(s\) is not unique. To say what
the body's orientation is at a given shape, one must choose body axes for that
shape. A shape-dependent change of axes
\[
  h:\mathcal S\to G
\]
changes the representative \(g\). This is exactly a gauge choice.

\begin{figure}[htbp]
\centering
\begin{tikzpicture}[scale=1.0, >=Latex]
  \draw[->] (-0.2,0) -- (6.5,0) node[right] {\(s^1\)};
  \draw[->] (0,-0.2) -- (0,4.1) node[above] {\(s^2\)};
  \draw[thick, blue, ->] (1.2,1.2) .. controls (2.8,3.4) and (5.0,3.1) .. (5.0,1.4);
  \draw[thick, blue, ->] (5.0,1.4) .. controls (3.9,0.4) and (2.2,0.3) .. (1.2,1.2);
  \node[blue] at (3.3,3.55) {shape cycle};
  \begin{scope}[shift={(1.2,1.2)}, scale=0.28]
    \draw[fill=gray!20, thick] (-0.8,0) ellipse (0.45 and 0.25);
    \draw[fill=gray!20, thick] (0.8,0) ellipse (0.45 and 0.25);
    \draw[thick] (-0.35,0) -- (0.35,0);
  \end{scope}
  \begin{scope}[shift={(5.0,1.4)}, scale=0.28, rotate=-20]
    \draw[fill=gray!20, thick] (-1.0,0) ellipse (0.55 and 0.18);
    \draw[fill=gray!20, thick] (1.0,0) ellipse (0.55 and 0.18);
    \draw[thick] (-0.45,0) -- (0.45,0);
  \end{scope}
  \begin{scope}[shift={(5.35,3.05)}, scale=0.28, rotate=45]
    \draw[fill=gray!20, thick] (-0.8,0) ellipse (0.45 and 0.25);
    \draw[fill=gray!20, thick] (0.8,0) ellipse (0.45 and 0.25);
    \draw[thick] (-0.35,0) -- (0.35,0);
    \draw[red!80!black, thick, ->] (0.0,0.65) arc (90:300:0.45);
  \end{scope}
  \draw[red, thick, ->] (4.7,2.55) -- (5.15,2.85);
  \node[align=center] at (3.2,0.65) {closed shape loop\\can have nonzero holonomy};
  \node[red!80!black, align=center] at (5.6,2.45) {net rigid\\rotation};
\end{tikzpicture}
\caption{Gauge kinematics: a closed loop in shape space can produce a net rigid
motion. The start and end shape can be the same while the rigid placement has
changed; the effect is the holonomy of a connection.}
\label{fig:gauge-shape-cycle-holonomy}
\end{figure}

\subsection{Mechanical Connection From Zero Momentum}
\label{subsec:mechanical-connection-zero-momentum}

The mechanical connection is obtained by solving a constraint, not by adding a
new force law. If the body has no total rigid momentum, then any internal shape
velocity that would otherwise create angular or linear momentum must be
cancelled by an accompanying rigid velocity. The connection is the linear rule
that performs this cancellation at each shape.

Let the body velocity in the moving frame be
\[
  \xi=g^{-1}\dot g\in\mathfrak g,
\]
where \(\mathfrak g\) is the Lie algebra of \(G\). Suppose the kinetic energy
has the block form
\[
  T
  =
  \frac12
  \begin{pmatrix}
  \xi\\ \dot s
  \end{pmatrix}^T
  \begin{pmatrix}
  M(s) & A(s)\\
  A(s)^T & m(s)
  \end{pmatrix}
  \begin{pmatrix}
  \xi\\ \dot s
  \end{pmatrix}.
\]
The momentum conjugate to rigid motion is
\[
  \Pi=\frac{\partial T}{\partial \xi}=M(s)\xi+A(s)\dot s.
\]
The symbol \(\Pi\) is the total rigid-motion momentum written in body
variables: for \(G=\SE(3)\) it packages total linear and angular momentum.
If the total angular and linear momentum vanish,
\[
  \Pi=0,
\]
as they do for an isolated body prepared with zero net external impulse, then
\(\Pi\) remains zero by conservation and
\[
  \xi=-M(s)^{-1}A(s)\dot s.
\]
Writing \(\dot s=\dot s^\alpha\partial_\alpha\), define the connection
\[
  \boxed{
  \mathcal A_\alpha(s)=M(s)^{-1}A_\alpha(s).
  }
\]
The rigid velocity is then
\[
  \boxed{
  g^{-1}\dot g=-\mathcal A_\alpha(s)\dot s^\alpha.
  }
\]
This is the cleanest derivation of the gauge field: it is the rule that tells a
zero-momentum deforming body how much rigid motion is forced by an infinitesimal
shape change.

The minus sign is part of the reconstruction convention. A positive connection
component means that the chosen shape change would create rigid momentum in the
body frame; the actual zero-momentum motion supplies the opposite rigid
velocity to cancel it. This is the mechanical meaning behind the gauge-theory
formula.

\subsection{Explicit Inertia-Operator Form}
\label{subsec:explicit-gauge-equation}

A concrete \(\SO(3)\) version of the same construction is useful because it
turns the abstract mechanical connection into a linear algebra problem. Let
\[
  A=R^{-1}\dot R\in\mathfrak{so}(3)
\]
be the angular velocity matrix of the body frame. Let \(Q=Q^T\) encode the
instantaneous shape through the second moment tensor
\[
  Q_{ij}=\int_B\rho(X)X_iX_j\,\dd^3X
\]
in the chosen internal frame. The corresponding inertia tensor is
\[
  \widetilde I_{ij}
  =
  \delta_{ij}\sum_{k=1}^3Q_{kk}-Q_{ij}.
\]
Let \(Y(X,s(t))\) be the position of a material point in the chosen internal
frame, with the center of mass at the origin. The spatial position is
\(x=RY\). In the body frame,
\[
  R^{-1}\dot x=A Y+\dot Y.
\]
The first term is the rigid velocity required by the frame motion; the second
term is internal deformation. The body angular momentum is therefore
\[
  M_{\mathrm{tot}}
  =
  \int_B\rho\,Y\times(AY+\dot Y)\,\dd^3X.
\]
With the hat convention used for the rigid body,
\(\widehat v\,w=v\times w\), or
\[
  (\widehat v)_{ij}=-\epsilon_{ijk}v_k,
\]
the rigid part obeys the matrix identity
\[
  \widehat{\widetilde I\Omega}=AQ+QA,
  \qquad A=\widehat\Omega.
\]
The internal part is
\[
  J_{\mathrm{int}}
  =
  \int_B\rho\,Y\times\dot Y\,\dd^3X.
\]
For zero total angular momentum,
\[
  \widetilde I\Omega+J_{\mathrm{int}}=0.
\]
It is convenient to define the deformation-generated momentum that the rigid
motion must cancel by
\[
  \widetilde J=-J_{\mathrm{int}}.
\]
Then the zero-momentum condition takes the linear matrix form
\[
  \boxed{
  AQ+QA=\widehat{\widetilde J}.
  }
\]
Solving this equation is the explicit finite-dimensional gauge connection:
the angular velocity \(A\) is determined by the shape velocity through
\(\widetilde J\).

\begin{derivation}[Solving the finite-dimensional gauge equation]
Write
\[
  A=\widehat\Omega,\qquad A_{ij}=-\epsilon_{ijk}\Omega_k.
\]
To verify the identity \(AQ+QA=\widehat{\widetilde I\Omega}\), diagonalize the
symmetric tensor \(Q\) at the instant under consideration:
\[
  Q=\diag(q_1,q_2,q_3).
\]
Then
\[
  (AQ+QA)_{12}=-(q_1+q_2)\Omega_3,\quad
  (AQ+QA)_{23}=-(q_2+q_3)\Omega_1,\quad
  (AQ+QA)_{31}=-(q_3+q_1)\Omega_2.
\]
But the inertia tensor has principal moments
\[
  \widetilde I_1=q_2+q_3,\qquad
  \widetilde I_2=q_3+q_1,\qquad
  \widetilde I_3=q_1+q_2,
\]
so these entries are exactly those of
\(\widehat{\widetilde I\Omega}\). Since both sides transform tensorially under
orthogonal changes of the internal frame, the identity holds in any frame.
At zero total angular momentum, the rigid angular momentum must equal
\(\widetilde J\), the negative of the internal contribution. The linear map
from angular velocity \(\Omega\) to rigid angular momentum is precisely the
inertia tensor:
\[
  \widetilde J_i=\widetilde I_{ij}\Omega_j.
\]
Equivalently,
\[
  \Omega_i=(\widetilde I^{-1})_{ij}\widetilde J_j.
\]
Therefore
\[
  \boxed{
  A_{ij}
  =
  -\epsilon_{ijk}
  (\widetilde I^{-1})_{k\ell}\widetilde J_\ell.
  }
\]
This is the concrete \(\SO(3)\) deforming-body gauge formula. In the abstract
notation above, \(M(s)\) is the inertia operator
\(\widetilde I(s)\), while the cross term \(A(s)\dot s\) is the internally
generated angular momentum \(\widetilde J(s,\dot s)\). The helicopter rotor
laboratory in Section~\ref{sec:lab-helicopter-gauge} is the finite-dimensional
special case where the shape variables are two rotor angles and this abstract
inertia-operator equation reduces to a pair of scalar connection components.
\end{derivation}

\subsection{Gauge Transformation And Holonomy}
\label{subsec:gauge-holonomy}

Under a shape-dependent change of body frame \(h(s)\), the connection transforms
as
\[
  \boxed{
  \mathcal A'
  =
  \Ad_{h^{-1}}\mathcal A-h^{-1}\dd h.
  }
\]
The sign of the inhomogeneous term follows from the reconstruction convention
\(g^{-1}\dot g=-\mathcal A(\dot s)\). With the opposite convention, the same
geometric connection is written with the more familiar plus sign. The base space
is shape space and the structure group is the group of rigid motions.

For a path \(s(t)\), the rigid motion is obtained by integrating
\[
  g^{-1}\dot g=-\mathcal A_\alpha(s)\dot s^\alpha.
\]
For a closed shape cycle \(\Gamma\), the net displacement is
\[
  g(T)
  =
  g(0)\,
  \mathcal P\exp\left(
    -\oint_\Gamma \mathcal A
  \right),
\]
where \(\mathcal P\) denotes path ordering. The result can be nontrivial even
though the shape returns to its starting point.

The curvature of the connection is
\[
  \boxed{
  \mathcal F
  =
  \dd\mathcal A+\mathcal A\wedge\mathcal A,
  }
\]
or in components,
\[
  \mathcal F_{\alpha\beta}
  =
  \partial_\alpha\mathcal A_\beta
  -
  \partial_\beta\mathcal A_\alpha
  +
  [\mathcal A_\alpha,\mathcal A_\beta].
\]
For a small loop enclosing area element
\(\Delta S^{\alpha\beta}\),
\[
  \Delta g
  \approx
  \exp\left(
    -\mathcal F_{\alpha\beta}\Delta S^{\alpha\beta}
  \right).
\]

\begin{figure}[htbp]
\centering
\begin{tikzpicture}[scale=1.0, >=Latex]
  \draw[->] (-0.2,0) -- (6.3,0) node[right] {\(s^1\)};
  \draw[->] (0,-0.2) -- (0,3.7) node[above] {\(s^2\)};
  \coordinate (a) at (1.6,1.0);
  \coordinate (b) at (4.4,1.0);
  \coordinate (c) at (4.4,2.7);
  \coordinate (d) at (1.6,2.7);
  \fill[blue!10] (a) -- (b) -- (c) -- (d) -- cycle;
  \draw[blue!70!black, thick, ->] (a) -- node[below] {\(\delta s^1\)} (b);
  \draw[blue!70!black, thick, ->] (b) -- node[right] {\(\delta s^2\)} (c);
  \draw[blue!70!black, thick, ->] (c) -- (d);
  \draw[blue!70!black, thick, ->] (d) -- (a);
  \draw[red!75!black, thick, ->] (3.0,1.85) arc (210:-80:0.45);
  \node[red!75!black, align=center] at (5.15,2.35)
    {net rigid motion\\\(-\mathcal F_{12}\delta s^1\delta s^2\)};
  \node[align=center] at (3.0,0.45)
    {same infinitesimal shape loop,\\different final placement};
\end{tikzpicture}
\caption{Curvature is the local density of holonomy. Traversing a small
rectangle in shape space returns to the same shape but can leave a residual
rigid displacement proportional to the curvature through the enclosed area.}
\label{fig:gauge-curvature-small-loop}
\end{figure}

Thus curvature is the local density of geometric phase. This is the same
mathematical idea as Hannay angle, Berry phase, and the rigid-body falling-cat
effect, now written for deforming bodies.

Flat connections can still produce frame changes along open paths, but a small
closed loop detects curvature. This is why the curvature, rather than the
connection coefficients themselves, is the local observable of the gauge
description. Different choices of body frame change \(\mathcal A\), while the
holonomy around a loop changes by conjugation and therefore represents the same
physical displacement.

\section{Swimming At Low Reynolds Number}
\label{sec:low-reynolds-swimming}

Shapere and Wilczek's low-Reynolds-number swimming problem gives a particularly
physical gauge theory. At very small Reynolds number, inertia drops out of
Navier-Stokes and the surrounding fluid satisfies Stokes equations:
\[
  -\nabla p+\mu\Delta u=0,\qquad \nabla\cdot u=0.
\]
Boundary velocity is prescribed by the swimmer's rigid motion plus deformation
rate. Because the Stokes problem is linear, the force and torque on the body
depend linearly on:
\[
  \xi=g^{-1}\dot g
  \qquad\text{and}\qquad
  \dot s.
\]
A freely swimming body has zero net force and torque, so the same algebra as
above gives
\[
  \xi=-\mathcal A_\alpha(s)\dot s^\alpha.
\]
The connection is now computed from a hydrodynamic boundary-value problem
rather than from an inertia tensor. A reciprocal one-parameter stroke encloses
no area in shape space and produces no net motion: this is Purcell's scallop
theorem in gauge language. A two-parameter stroke can enclose area and swim by
curvature.

\section{Local Frames In Elastic Continua}
\label{sec:local-frames-elastic-continua}

The finite-dimensional gauge picture becomes a field theory when every material
point carries a local frame. A Cosserat body has, in addition to
\(\varphi(X,t)\), a rotation field
\[
  R(X,t)\in\SO(3).
\]
The directors are
\[
  d_a(X,t)=R(X,t)E_a,
\]
where \(E_a\) is a fixed reference basis. The rotational connection is
\[
  A_A=R^{-1}\partial_A R\in\mathfrak{so}(3),
  \qquad
  A_t=R^{-1}\partial_t R.
\]
The translational strain measured in the local frame is
\[
  \Gamma_A=R^{-1}\partial_A\varphi.
\]
For a rigid body, \(R\) depends only on time. For a Cosserat continuum, \(R\)
depends on \(X\), so neighboring material points can rotate relative to one
another.

This is the field-theoretic version of the rigid-body frame issue. In the
rigid body there is one moving frame; in a Cosserat body there is a moving frame
at every material point. The connection \(A_A\) compares nearby frames, while
\(A_t\) records how a given local frame evolves in time.

\begin{definition}[Local frame gauge transformation]
Let \(G(X,t)\in\SO(3)\) be a change of local material frame:
\[
  R'(X,t)=R(X,t)G(X,t).
\]
Then
\[
  A'_A=G^{-1}A_AG+G^{-1}\partial_A G,
  \qquad
  \Gamma'_A=G^{-1}\Gamma_A.
\]
Thus \(A_A\) is a gauge connection and \(\Gamma_A\) is a vector-valued
matter field in the local frame.
\end{definition}

\section{Curvature, Torsion, And Defects}
\label{sec:curvature-torsion-defects}

If a connection is globally of the form
\[
  A_A=R^{-1}\partial_A R,
\]
then it is pure gauge and satisfies the Maurer-Cartan identity:
\[
  F_{AB}
  =
  \partial_AA_B-\partial_BA_A+[A_A,A_B]
  =
  0.
\]
Nonzero curvature means that local frames cannot be integrated into a single
smooth orientation field without singularities or defects. In defect theory,
curvature is associated with disclinations, or rotational defects.

To describe translational defects, introduce a local coframe \(\theta^a\), which
measures material line elements in a local frame. The connection one-forms
\(\omega^a{}_b\) are the components of the same \(\mathfrak{so}(3)\)-valued
connection written above: if the matrix entries of \(A_A\) are
\((A_A)^a{}_b\), then
\[
  \omega^a{}_b=(A_A)^a{}_b\,\dd X^A.
\]
With this notation, the torsion two-form is
\[
  T^a=\dd\theta^a+\omega^a{}_b\wedge\theta^b.
\]
Torsion measures closure failure of infinitesimal parallelograms and is
associated with dislocation density. In this language:
\[
\begin{aligned}
  \text{curvature} &\longleftrightarrow \text{rotational mismatch},\\
  \text{torsion} &\longleftrightarrow \text{translational mismatch},\\
  \text{incompatible strain} &\longleftrightarrow \text{nonintegrable local data}.
\end{aligned}
\]
The gauge theory is not decorative terminology. It records which descriptions
depend on arbitrary local choices of frame and which quantities are invariant.

\section{Cosserat Rod As A One-Dimensional Gauge Theory}
\label{sec:cosserat-rod}

A planar rod has material coordinate \(s\), centerline \(r(s)\in\R^2\), and
orientation angle \(\theta(s)\). Its director is
\[
  d_1=(\cos\theta,\sin\theta).
\]
For an inextensible unshearable rod,
\[
  r'(s)=d_1(s).
\]
The connection is the scalar curvature
\[
  \kappa(s)=\theta'(s).
\]
Given \(\kappa\),
\[
  \theta(s)=\theta(0)+\int_0^s\kappa(\sigma)\,\dd\sigma,
\]
and
\[
  r(s)=r(0)+\int_0^s
  (\cos\theta(\sigma),\sin\theta(\sigma))\,\dd\sigma.
\]
This is the reconstruction problem: a gauge field \(\kappa\,\dd s\) determines
shape after choosing initial position and frame.

The bending energy
\[
  E_{\mathrm{bend}}
  =
  \frac12\int_0^L B(\kappa-\kappa_0)^2\,\dd s
\]
compares actual curvature with preferred curvature. Here \(B\) is bending
stiffness and \(\kappa_0\) is intrinsic curvature. In gauge language, the
energy penalizes mismatch between a connection and a preferred connection.

\section{Conceptual Summary}
\label{sec:elastic-conceptual-summary}

The progression from rigid body to elastic gauge theory is:
\[
\begin{aligned}
  \text{rigid body:}&\quad R(t),\quad R^{-1}\dot R,\\
  \text{deformable body:}&\quad g(t),\ s(t),\quad g^{-1}\dot g=-\mathcal A(s)\dot s,\\
  \text{Cosserat continuum:}&\quad R(X,t),\quad R^{-1}\partial_A R,\\
  \text{defect theory:}&\quad A,\quad F=\dd A+A\wedge A,\quad T=\dd\theta+\omega\wedge\theta.
\end{aligned}
\]
The first line is ordinary rigid-body mechanics. The second is
Shapere-Wilczek gauge kinematics over shape space. The third attaches a local
rigid body to every material point. The fourth allows the local frame and
translation data to be incompatible, producing curvature and torsion as physical
defect densities.

\section{Associated Demos}
\label{sec:elastic-associated-demos}

\begin{itemize}
\item \path{demos/python/cosserat_rod_demo.py}, a planar frame reconstruction
from curvature.
\item \path{demos/mathematica/GaugeDeformingBody.wl}, a finite-dimensional
gauge-connection toy model.
\item Suggested extension: a two-shape-parameter swimmer showing geometric phase
as the area integral of connection curvature.
\end{itemize}

\chapter{Hamilton-Jacobi Theory And Adiabatic Invariants}
\label{chap:hamilton-jacobi}

\section{Why Hamilton-Jacobi Belongs Here}
\label{sec:why-hamilton-jacobi}

Hamilton-Jacobi theory belongs naturally between canonical transformations and
integrable systems. The Hamilton-Jacobi equation is not merely another way to
solve special examples. It explains why a canonical transformation generated by
one scalar function can turn a Hamiltonian flow into trivial motion, and it
gives the conceptual bridge from Hamiltonian mechanics to action-angle
variables.

\begin{table}[htbp]
\centering
\caption{Notation for Hamilton-Jacobi theory and adiabatic invariants.}
\label{tab:hj-notation}
\small
\begin{tabular}{p{0.18\textwidth}p{0.48\textwidth}p{0.25\textwidth}}
\toprule
Symbol & Meaning & Type\\
\midrule
\(q^i,p_i\) & old canonical coordinates and momenta & local coordinates on \(T^*Q\)\\
\(Q^i,P_i\) & new canonical coordinates and momenta & transformed phase-space coordinates\\
\(S(q,\alpha,t)\) & Hamilton's principal function or complete integral & generating function\\
\(\alpha_i\) & constants used as new momenta \(P_i\) & integration constants\\
\(\beta^i\) & constants \(\partial S/\partial\alpha_i\) & new coordinates\\
\(W(q,\alpha)\) & Hamilton's characteristic function & time-independent part of \(S\)\\
\(E\) & conserved energy in autonomous systems & scalar\\
\(I\) & one-dimensional action variable & \((2\pi)^{-1}\oint p\,\dd q\)\\
\(\theta\) & angle variable conjugate to \(I\) & \(2\pi\)-periodic\\
\(L(t)\) & slowly varying wall separation in the box example & length\\
\(\epsilon_{\mathrm{ad}}\) & adiabatic small parameter & period / driving time scale\\
\bottomrule
\end{tabular}
\end{table}

\paragraph{Roadmap.}
The chapter starts with generating functions because the Hamilton-Jacobi
equation is best understood as a canonical transformation problem. A complete
integral \(S(q,\alpha,t)\) is not just a solution of a PDE; it is a device that
turns the original Hamiltonian flow into constant new momenta and linearly
recoverable old coordinates. The autonomous case introduces the characteristic
function \(W\), the action integral connects this to invariant tori, and the
moving-wall example shows what remains true when parameters vary slowly rather
than exactly stay fixed.

\paragraph{Two meanings of ``action''.}
The word action appears in two related but distinct senses. Hamilton's
principle uses the time integral \(\int L\,\dd t\) as a functional on paths.
Action-angle theory uses \(I=(2\pi)^{-1}\oint p\,\dd q\) as a coordinate
labelling a closed orbit or invariant torus. Hamilton-Jacobi theory is the
bridge: its generating function accumulates the canonical one-form
\(p_i\,\dd q^i-H\,\dd t\), and its period integrals become the action
coordinates used in integrability and adiabatic invariance.

\section{Generating Functions And The Hamilton-Jacobi Equation}
\label{sec:hj-generating-function}

The key idea is to choose new canonical coordinates in which the motion is as
simple as possible. If the new Hamiltonian vanishes, then the new coordinates
and momenta are constants. The price is that the generating function must solve
a nonlinear first-order PDE. That PDE is the Hamilton-Jacobi equation.

Let \((q,p)\) be old canonical coordinates and \((Q,P)\) new canonical
coordinates. A time-dependent type-II generating function
\[
  S=S(q,P,t)
\]
defines a canonical transformation by
\[
  p_i=\frac{\partial S}{\partial q^i},
  \qquad
  Q^i=\frac{\partial S}{\partial P_i}.
\]
These formulas are not arbitrary. They are chosen so that
\[
  \dd S=p_i\,\dd q^i+Q^i\,\dd P_i+\frac{\partial S}{\partial t}\,\dd t.
\]
Since \(Q^i\dd P_i=\dd(P_iQ^i)-P_i\dd Q^i\), this can be rearranged as
\[
  p_i\,\dd q^i
  =
  P_i\,\dd Q^i
  +
  \dd(S-P_iQ^i)
  -
  \frac{\partial S}{\partial t}\,\dd t.
\]
The phase-space action transforms as
\[
  p_i\,\dd q^i-H\,\dd t
  =
  P_i\,\dd Q^i-K\,\dd t+\dd(S-P_iQ^i),
\]
where \(K(Q,P,t)\) is the new Hamiltonian. Comparing coefficients gives
\[
  K(Q,P,t)
  =
  H\!\left(q,\frac{\partial S}{\partial q},t\right)
  +
  \frac{\partial S}{\partial t}.
\]
The Hamilton-Jacobi strategy is to choose \(S\) so that the new Hamiltonian is
as simple as possible. The strongest choice is
\[
  K=0.
\]
Then \(P_i\) and \(Q^i\) are all constants along the new motion. Writing
\(P_i=\alpha_i\), the required equation is
\[
  \boxed{
  \frac{\partial S}{\partial t}
  +
  H\!\left(q,\frac{\partial S}{\partial q},t\right)=0.
  }
\]
This is the Hamilton-Jacobi equation.

\begin{figure}[htbp]
\centering
\begin{tikzpicture}[scale=1.0, >=Latex]
  \draw[->] (-0.2,0) -- (6.4,0) node[right] {\(q^1\)};
  \draw[->] (0,-0.2) -- (0,3.6) node[above] {\(q^2\)};
  \draw[blue!70!black, thick]
    (0.6,0.6) .. controls (1.8,1.1) and (3.1,1.0) .. (5.4,0.65);
  \draw[blue!70!black, thick]
    (0.8,1.35) .. controls (2.0,1.85) and (3.35,1.7) .. (5.65,1.35);
  \draw[blue!70!black, thick]
    (1.0,2.1) .. controls (2.25,2.65) and (3.65,2.45) .. (5.9,2.1);
  \draw[red!75!black, thick, ->]
    (1.45,0.78) .. controls (1.65,1.28) and (1.9,1.68) .. (2.2,2.24);
  \draw[red!75!black, thick, ->]
    (3.0,0.9) .. controls (3.15,1.35) and (3.35,1.75) .. (3.72,2.42);
  \draw[red!75!black, thick, ->]
    (4.65,0.76) .. controls (4.85,1.20) and (5.05,1.55) .. (5.38,2.15);
  \draw[->, black!70] (3.35,1.7) -- (3.85,2.05)
    node[right] {\(p=\nabla_qS\)};
  \node[blue!70!black, align=center] at (1.05,3.15)
    {level sets\\\(S(q,t)=\text{constant}\)};
  \node[red!75!black, align=center] at (5.15,3.05)
    {characteristics\\classical paths};
\end{tikzpicture}
\caption{Hamilton-Jacobi geometry. At a fixed time the level sets of
\(S(q,t)\) are wavefronts in configuration space; the momentum covector
\(p=\nabla_qS\) is normal to those wavefronts, and Hamilton's equations carry
these normals along the characteristic curves.}
\label{fig:hj-wavefront-characteristics}
\end{figure}

\begin{definition}[Complete integral]
A function \(S(q,\alpha,t)\) is a complete integral of the Hamilton-Jacobi
equation for an \(n\)-degree-of-freedom system if it depends on \(n\) independent
parameters \(\alpha_1,\ldots,\alpha_n\) and satisfies
\[
  \det\left(
    \frac{\partial^2S}{\partial q^i\partial\alpha_j}
  \right)\neq0
\]
on the region under consideration. The nondegeneracy condition is exactly what
lets \((q,\alpha)\mapsto(q,p)\) define a local canonical transformation.
\end{definition}

Once a complete integral is known, the solution is obtained algebraically:
\[
  p_i(t)=\frac{\partial S}{\partial q^i}(q(t),\alpha,t),
  \qquad
  \beta^i=\frac{\partial S}{\partial\alpha_i}(q(t),\alpha,t),
\]
where \(\alpha_i\) and \(\beta^i\) are constants fixed by initial data. This is
why Hamilton-Jacobi theory is sometimes called integration by quadrature: the
hard part is constructing \(S\); after that the motion is recovered by
differentiation and inversion.

The nondegeneracy condition in the definition is essential. Without it,
\(\alpha\) would not provide a valid set of new momenta and the equations
\(\partial S/\partial\alpha=\beta\) could fail to determine \(q(t)\). Thus a
Hamilton-Jacobi solution is useful not merely because it solves a PDE, but
because it generates a canonical coordinate system adapted to the flow.

\section{Autonomous Systems And The Characteristic Function}
\label{sec:hj-autonomous}

If \(H(q,p)\) is independent of time, seek a separated form
\[
  S(q,\alpha,t)=W(q,\alpha)-Et,
\]
where one of the constants \(\alpha_i\) is the energy \(E\). The
Hamilton-Jacobi equation becomes
\[
  \boxed{
  H\!\left(q,\frac{\partial W}{\partial q}\right)=E.
  }
\]
This is the time-independent Hamilton-Jacobi equation. The function \(W\) is
Hamilton's characteristic function.

For one degree of freedom with Hamiltonian
\[
  H(q,p)=\frac{p^2}{2m}+V(q),
\]
the equation reads
\[
  \frac1{2m}\left(\frac{\dd W}{\dd q}\right)^2+V(q)=E.
\]
Thus
\[
  \frac{\dd W}{\dd q}=p(q,E)=\pm\sqrt{2m(E-V(q))},
\]
and
\[
  W(q,E)=\int^q p(q',E)\,\dd q'.
\]
This formula is simple but profound: the phase of the Hamilton-Jacobi function
is the accumulated canonical one-form \(p\,\dd q\).

\section{Action As A Period Integral}
\label{sec:hj-action-period}

Suppose a one-degree-of-freedom autonomous system has a bounded orbit between
turning points \(q_-(E)\) and \(q_+(E)\). The action variable is
\[
  \boxed{
  I(E)=\frac{1}{2\pi}\oint p\,\dd q
  =
  \frac1{\pi}\int_{q_-(E)}^{q_+(E)}
  \sqrt{2m(E-V(q))}\,\dd q.
  }
\]
The factor \(1/\pi\) appears because the closed orbit has a positive-momentum
branch and a negative-momentum branch.

\begin{figure}[htbp]
\centering
\begin{tikzpicture}[scale=1.0, >=Latex]
  \draw[->] (-0.4,0) -- (6.6,0) node[right] {\(q\)};
  \draw[->] (0,-2.0) -- (0,2.2) node[above] {\(p\)};
  \draw[densely dotted] (1.0,-1.7) -- (1.0,1.8);
  \draw[densely dotted] (5.3,-1.7) -- (5.3,1.8);
  \node[below] at (1.0,0) {\(q_-\)};
  \node[below] at (5.3,0) {\(q_+\)};
  \draw[blue!70!black, thick, ->]
    (1.0,0) .. controls (2.0,1.55) and (4.2,1.55) .. (5.3,0);
  \draw[blue!70!black, thick, ->]
    (5.3,0) .. controls (4.2,-1.55) and (2.0,-1.55) .. (1.0,0);
  \node[blue!70!black] at (3.2,1.65) {\(+p(q,E)\)};
  \node[blue!70!black] at (3.2,-1.65) {\(-p(q,E)\)};
  \node[align=center] at (3.2,0.55)
    {enclosed area\\\(\displaystyle \oint p\,\dd q=2\pi I\)};
\end{tikzpicture}
\caption{For a one-dimensional bound orbit the action is the signed
phase-space area enclosed by the energy curve divided by \(2\pi\). The turning
points contribute no endpoint term when differentiating with respect to
energy because \(p(q_\pm,E)=0\).}
\label{fig:hj-action-loop}
\end{figure}

The period follows from differentiating \(I(E)\). The moving endpoint terms
vanish because the momentum is zero at the turning points:
\[
  p(q_\pm(E),E)=0.
\]
Thus
\[
  \frac{\dd I}{\dd E}
  =
  \frac1{\pi}\int_{q_-}^{q_+}
  \frac{m\,\dd q}{\sqrt{2m(E-V(q))}}
  =
  \frac1{2\pi}T(E).
\]
Therefore
\[
  \boxed{
  \omega(E)=\frac{\dd H}{\dd I}
  =
  \left(\frac{\dd I}{\dd E}\right)^{-1},
  \qquad
  T(E)=2\pi\frac{\dd I}{\dd E}.
  }
\]
This formula is the one-dimensional seed of action-angle theory.

\begin{example}[Harmonic oscillator]
For
\[
  H=\frac{p^2}{2m}+\frac12m\omega_0^2q^2,
\]
the phase curve is an ellipse:
\[
  \frac{p^2}{2mE}+\frac{q^2}{2E/(m\omega_0^2)}=1.
\]
The area enclosed by the ellipse is
\[
  \pi\sqrt{2mE}\sqrt{\frac{2E}{m\omega_0^2}}
  =
  \frac{2\pi E}{\omega_0}.
\]
Thus
\[
  I=\frac{1}{2\pi}\operatorname{Area}=\frac{E}{\omega_0},
  \qquad
  H(I)=\omega_0 I.
\]
The conjugate angle advances uniformly:
\[
  \dot\theta=\frac{\partial H}{\partial I}=\omega_0.
\]
\end{example}

\section{Hamilton-Jacobi And Wave Optics}
\label{sec:hj-wave-analogy}

Hamilton-Jacobi theory is also the classical side of a broader ray-wave
analogy. The analogy begins with the eikonal approximation. If a wave has
rapidly oscillating phase
\[
  \psi(q,t)\approx A(q,t)e^{\frac{i}{\hbar}S(q,t)},
\]
then substituting into the Schrodinger equation and keeping the leading power
of \(\hbar\) gives
\[
  \partial_tS+H\left(q,\partial_qS\right)=0.
\]
Thus classical mechanics is the short-wavelength or stationary-phase limit of
wave mechanics. In optics, rays are normals to wavefronts. In mechanics,
trajectories are characteristics of the Hamilton-Jacobi equation, and the
momentum is the phase gradient:
\[
  p_i=\partial_iS.
\]

\section{Adiabatic Invariant: Particle Between Slowly Moving Walls}
\label{sec:adiabatic-box}

A clean adiabatic-invariant model is a particle that bounces elastically
between a fixed wall and a slowly moving wall. Let \(L(t)\) be the wall
separation and \(p(t)=m v(t)>0\) the magnitude of the particle momentum between
collisions.
The example is intentionally elementary: no differential geometry is needed to
see the invariant. The point is that the same quantity obtained from collision
kinematics is also the action integral \((2\pi)^{-1}\oint p\,\dd q\). This is
why adiabatic invariance belongs with action-angle variables.

\begin{assumption}[Adiabatic wall motion]
The wall speed is small compared with the particle speed:
\[
  \epsilon_{\mathrm{ad}}
  =
  \frac{|\dot L|}{v}\ll1.
\]
Over one round trip, the wall speed can be treated as approximately constant.
\end{assumption}

First derive the invariant by elementary collision kinematics. If the moving
wall approaches the fixed wall with speed \(u=-\dot L>0\), an elastic collision
with that wall increases the particle speed by approximately \(2u\). During one
round trip,
\[
  \Delta t=\frac{2L}{v},
  \qquad
  \Delta L=-u\Delta t=-\frac{2Lu}{v}.
\]
The momentum change is
\[
  \Delta p=2mu=\frac{2pu}{v}.
\]
Eliminating \(u/v\) gives
\[
  \Delta p=-\frac{p}{L}\Delta L.
\]
Therefore
\[
  \Delta(pL)=L\Delta p+p\Delta L=0
\]
to leading adiabatic order. Hence
\[
  \boxed{p(t)L(t)\approx \text{constant}.}
\]

The action interpretation explains why this result is general. For a particle
in a one-dimensional box, the phase-space loop consists of a branch with
momentum \(+p\) from \(0\) to \(L\) and a branch with momentum \(-p\) from
\(L\) back to \(0\). Hence
\[
  \oint p\,\dd q=2pL,
  \qquad
  I=\frac{pL}{\pi}.
\]
Slow parameter changes preserve \(I\) to adiabatic accuracy, so \(pL\) is the
adiabatic invariant. The energy scales as
\[
  E(t)=\frac{p(t)^2}{2m}
  =
  E(0)\frac{L(0)^2}{L(t)^2}.
\]

\begin{figure}[htbp]
\centering
\begin{tikzpicture}[scale=1.0, >=Latex]
  \draw[->] (-0.4,0) -- (6.4,0) node[right] {\(q\)};
  \draw[->] (0,-1.9) -- (0,2.1) node[above] {\(p\)};
  \draw[thick] (0.8,-1.25) rectangle (5.4,1.25);
  \draw[blue!75!black, thick, ->] (0.8,1.25) -- (5.4,1.25);
  \draw[blue!75!black, thick, ->] (5.4,-1.25) -- (0.8,-1.25);
  \draw[dashed] (5.4,-1.55) -- (5.4,1.55);
  \draw[->, red!75!black, thick] (5.4,1.65) -- (4.8,1.65)
    node[left] {slow wall motion};
  \node[below] at (0.8,0) {\(0\)};
  \node[below] at (5.4,0) {\(L(t)\)};
  \node[right] at (5.4,1.25) {\(+p\)};
  \node[right] at (5.4,-1.25) {\(-p\)};
  \node[align=center] at (3.1,0)
    {area \(=2pL\)\\\(I=pL/\pi\)};
\end{tikzpicture}
\caption{The adiabatic box in phase space. Slow wall motion changes the
rectangle's width \(L\) and height \(p\), but preserves the enclosed area
\(2pL\) to leading adiabatic order.}
\label{fig:adiabatic-box-action}
\end{figure}

This example also explains the limitation of the word ``invariant.'' The
quantity \(I\) is not exactly conserved for arbitrary wall motion. It is
preserved to high accuracy when the parameter changes slowly compared with the
orbital period, and the accuracy improves as the ratio of time scales becomes
smaller. Adiabatic invariance is therefore a controlled approximation tied to a
separation of time scales.

\section{What Must Be Remembered}
\label{sec:hj-summary}

\begin{itemize}
\item The Hamilton-Jacobi equation is the condition that a canonical
transformation make the new Hamiltonian trivial.
\item A complete integral \(S(q,\alpha,t)\) contains the full solution, but
only when the mixed Hessian in \((q,\alpha)\) is nondegenerate.
\item In autonomous systems, \(S=W-Et\), and \(W\) solves
\(H(q,\partial W/\partial q)=E\).
\item Action variables are period integrals of the canonical one-form.
\item Adiabatic invariance is the slow-parameter version of action
conservation.
\end{itemize}

The conceptual bridge to the earlier chapters is this: Liouville integrability
produces invariant tori, action variables measure cycles on those tori, and
Hamilton-Jacobi theory constructs the canonical transformation that reveals
those variables when the system is separable enough.

\chapter{Worked Problem Laboratories}
\label{chap:worked-problem-labs}

\section{Purpose And Scope}
\label{sec:worked-problem-labs-purpose}

This chapter collects self-contained lecture-note laboratories. Each laboratory
has a precise model, a defined notation set, and a derivation that can be read
from the definitions and assumptions given here. The examples are intentionally
varied: they show how the abstract machinery in the earlier chapters behaves
when used on concrete mechanics problems.

\begin{table}[htbp]
\centering
\caption{Notation used in the worked problem laboratories. Local definitions in
each section override this table when stated.}
\label{tab:worked-lab-notation}
\small
\begin{tabular}{p{0.18\textwidth}p{0.48\textwidth}p{0.25\textwidth}}
\toprule
Symbol & Meaning & Type\\
\midrule
\(s^a\) & new generalized coordinates under a point transformation & coordinates on \(Q\)\\
\(q^i\) & old generalized coordinates & coordinates on \(Q\)\\
\(V(x)\) & one-dimensional potential energy & scalar function\\
\(\gamma\) & damping coefficient in the oscillator example & inverse time\\
\(\omega\) & imposed angular speed in rotating-hoop example & inverse time\\
\(a\) & hoop radius or semimajor axis depending on section & locally defined\\
\(R,r\) & cylinder radii or Schwarzschild length/radius depending on section & locally defined\\
\(N\) & normal force in rolling-contact examples & force\\
\(A(r)\) & \(1-R/r\) in the Schwarzschild-type example & dimensionless\\
\(E,L_z\) & conserved energy and angular momentum in central examples & constants of motion\\
\(H_{\mathrm{HH}}\) & Henon-Heiles Hamiltonian & Hamiltonian function\\
\(\lambda_{\mathrm{HH}}\) & Henon-Heiles perturbation strength & small parameter\\
\(I_i,\theta_i\) & oscillator actions and angles & canonical coordinates\\
\(r_i,v_i\) & positions and velocities of three equal masses & vectors in \(\R^2\)\\
\(m,k\) & mass and Newtonian coupling in three-body lab & positive constants\\
\(\alpha,\beta\) & internal rotor angles in the helicopter lab & shape coordinates\\
\(\mathcal A_\alpha,\mathcal A_\beta\) & gauge connection components over rotor shape space & \(\mathfrak{so}(3)\)-valued\\
\bottomrule
\end{tabular}
\end{table}

\paragraph{How to use this chapter.}
Each laboratory is deliberately small, but each one isolates a recurring
mechanics move: change coordinates, test regularity, encode dissipation with
time dependence, reduce a constrained rolling problem, test circular-orbit
stability, expose resonance, or reconstruct motion from a gauge connection.
The value of the chapter is cumulative. A reader should be able to take the
workflow of one section and recognize the same workflow in a more complicated
system later.

Each laboratory should be read with three questions in mind. What is the
minimal model? What calculation turns the model into a diagnostic quantity?
What would fail if one of the assumptions were changed? These questions are
what turn a worked example into a reusable method rather than a solved problem
to memorize.

\paragraph{How figures function in this chapter.}
The figures are not decorative summaries after the calculation. Each one is
attached to a modelling decision: a coordinate choice, a constraint force, an
effective potential, a phase-space section, an initial-data geometry, or a
gauge loop. When using these laboratories as exercise problems, the figure
should be read before the algebra and then checked again after the final boxed
result. If a student cannot say which equation the figure represents, the
problem has not yet been understood geometrically.

\section{Form-Invariance Of Lagrange Equations}
\label{sec:lab-point-transformations}

We begin with the coordinate-independence of the Euler-Lagrange equations. Let
\[
  q^i=q^i(s^1,\ldots,s^n,t)
\]
be a regular time-dependent point transformation. Regular means that the
Jacobian \(\partial q^i/\partial s^a\) has rank \(n\). The transformed
Lagrangian is
\[
  L'(s,\dot s,t)=L(q(s,t),\dot q(s,\dot s,t),t),
\]
where
\[
  \dot q^i
  =
  \frac{\partial q^i}{\partial s^a}\dot s^a
  +
  \frac{\partial q^i}{\partial t}.
\]
Two elementary identities do the work:
\[
  \frac{\partial \dot q^i}{\partial \dot s^a}
  =
  \frac{\partial q^i}{\partial s^a},
  \qquad
  \frac{\dd}{\dd t}
  \left(
    \frac{\partial q^i}{\partial s^a}
  \right)
  =
  \frac{\partial \dot q^i}{\partial s^a}.
\]
Using the chain rule,
\[
  \frac{\partial L'}{\partial \dot s^a}
  =
  \frac{\partial L}{\partial \dot q^i}
  \frac{\partial q^i}{\partial s^a},
\]
and
\[
  \frac{\partial L'}{\partial s^a}
  =
  \frac{\partial L}{\partial q^i}
  \frac{\partial q^i}{\partial s^a}
  +
  \frac{\partial L}{\partial \dot q^i}
  \frac{\partial \dot q^i}{\partial s^a}.
\]
Therefore
\[
\begin{aligned}
  \frac{\dd}{\dd t}\frac{\partial L'}{\partial \dot s^a}
  -
  \frac{\partial L'}{\partial s^a}
  &=
  \left(
    \frac{\dd}{\dd t}\frac{\partial L}{\partial\dot q^i}
    -
    \frac{\partial L}{\partial q^i}
  \right)
  \frac{\partial q^i}{\partial s^a}.
\end{aligned}
\]
If \(q(t)\) satisfies the Euler-Lagrange equations, the right side vanishes;
if the transformation is regular, the converse holds as well. Thus the
Euler-Lagrange equation is not tied to a preferred coordinate system.

\section{A Nonstandard Lagrangian And Regularity}
\label{sec:lab-nonstandard-lagrangian}

This laboratory separates two ideas that are often conflated. A Lagrangian may
produce a familiar-looking equation of motion and still fail to be globally
regular. Regularity is a statement about whether velocities can be solved from
momenta. Without it, the Legendre transform can fold phase space and the
Hamiltonian description must be handled with care.

Consider the one-dimensional Lagrangian
\[
  L(x,\dot x)
  =
  \frac{m^2\dot x^4}{12}
  +
  m\dot x^2V(x)
  -
  V(x)^2.
\]
The momentum is
\[
  p=\frac{\partial L}{\partial\dot x}
  =
  \frac{m^2}{3}\dot x^3+2mV\dot x.
\]
The Euler-Lagrange equation gives
\[
\begin{aligned}
  0
  &=
  \frac{\dd}{\dd t}
  \left(
    \frac{m^2}{3}\dot x^3+2mV\dot x
  \right)
  -
  \left(
    m\dot x^2V'-2VV'
  \right)\\
  &=
  (m\dot x^2+2V)(m\ddot x+V').
\end{aligned}
\]
There are two branches. The Newtonian branch is
\[
  m\ddot x=-V'(x).
\]
The second branch,
\[
  m\dot x^2+2V(x)=0,
\]
is a singular energy-surface condition. The velocity Hessian
\[
  \frac{\partial^2L}{\partial\dot x^2}
  =
  m(m\dot x^2+2V)
\]
vanishes on this second branch. This example is a useful warning: a Lagrangian
can produce familiar equations on a regular branch while being singular on
another branch. Regularity is not decorative; it is what lets the Legendre
transform define a Hamiltonian system.

\begin{figure}[htbp]
\centering
\begin{tikzpicture}[scale=1.0, >=Latex]
  \draw[->] (-3.0,0) -- (3.15,0) node[right] {\(\dot x\)};
  \draw[->] (0,-2.2) -- (0,2.25) node[above] {\(p\)};
  \draw[blue!70!black, thick, domain=-2.55:2.55, samples=180]
    plot (\x,{0.33*\x*\x*\x-0.9*\x});
  \draw[dashed, red!75!black] (-0.95,-1.9) -- (-0.95,1.9);
  \draw[dashed, red!75!black] (0.95,-1.9) -- (0.95,1.9);
  \fill[red!75!black] (-0.95,0.57) circle (1.3pt)
    node[above left] {\(W=0\)};
  \fill[red!75!black] (0.95,-0.57) circle (1.3pt)
    node[below right] {\(W=0\)};
  \draw[decorate, decoration={brace, amplitude=5pt}]
    (-0.95,-1.15) -- (0.95,-1.15)
    node[midway, below=7pt] {multi-valued velocity branch};
  \node[blue!70!black, align=center] at (2.05,1.45)
    {folded\\Legendre map};
\end{tikzpicture}
\caption{Nonstandard Lagrangian laboratory. For a representative region with
\(V<0\), the momentum-velocity map can fold where the velocity Hessian
\(W=\partial^2L/\partial\dot x^2\) vanishes. At such points the Legendre
transform is not locally invertible, even though the Euler-Lagrange equation has
a familiar Newtonian branch.}
\label{fig:lab-nonstandard-legendre-fold}
\end{figure}
\FloatBarrier

\section{Damped Oscillator From A Time-Dependent Lagrangian}
\label{sec:lab-damped-oscillator}

Let
\[
  L(q,\dot q,t)
  =
  e^{\gamma t}
  \left(
    \frac12m\dot q^2-\frac12kq^2
  \right).
\]
The Euler-Lagrange equation is
\[
  \frac{\dd}{\dd t}(e^{\gamma t}m\dot q)+e^{\gamma t}kq=0,
\]
or
\[
  \boxed{
  m\ddot q+m\gamma\dot q+kq=0.
  }
\]
Although this is a dissipative equation, it comes from a time-dependent
Lagrangian. The canonical energy is not conserved because \(L\) depends
explicitly on time.

Set
\[
  s=e^{\gamma t/2}q,
  \qquad
  q=e^{-\gamma t/2}s.
\]
Then
\[
  \dot q=e^{-\gamma t/2}\left(\dot s-\frac{\gamma}{2}s\right).
\]
Substituting into the action gives
\[
  L
  =
  \frac12m\left(\dot s-\frac{\gamma}{2}s\right)^2
  -
  \frac12ks^2.
\]
The cross term is a total derivative:
\[
  -\frac12m\gamma s\dot s
  =
  -\frac{\dd}{\dd t}\left(\frac14m\gamma s^2\right).
\]
Discarding that total derivative, the equivalent Lagrangian is
\[
  L_{\mathrm{eff}}
  =
  \frac12m\dot s^2
  -
  \frac12m
  \left(
    \frac{k}{m}-\frac{\gamma^2}{4}
  \right)s^2.
\]
Thus
\[
  \ddot s+\left(\frac{k}{m}-\frac{\gamma^2}{4}\right)s=0.
\]
The transformation has moved damping into a shifted oscillator frequency.

\begin{figure}[htbp]
\centering
\begin{tikzpicture}[scale=1.0, >=Latex]
  \draw[->] (-0.2,0) -- (6.45,0) node[right] {\(t\)};
  \draw[->] (0,-1.45) -- (0,1.6) node[above] {amplitude};
  \draw[blue!70!black, dashed, domain=0:6.0, samples=120]
    plot (\x,{1.10*exp(-0.22*\x)});
  \draw[blue!70!black, dashed, domain=0:6.0, samples=120]
    plot (\x,{-1.10*exp(-0.22*\x)});
  \draw[blue!70!black, thick, domain=0:6.0, samples=180]
    plot (\x,{1.10*exp(-0.22*\x)*cos(2.7*\x r)});
  \draw[red!75!black, thick, domain=0:6.0, samples=180]
    plot (\x,{1.10*cos(2.7*\x r)});
  \node[blue!70!black] at (4.7,0.43) {\(q(t)\)};
  \node[red!75!black] at (4.5,1.25) {\(s=e^{\gamma t/2}q\)};
  \node[align=center] at (3.2,-1.85)
    {rescaling absorbs the exponential damping envelope};
\end{tikzpicture}
\caption{Damped-oscillator laboratory. In the original coordinate \(q\), the
motion has a decaying envelope. The rescaled coordinate \(s=e^{\gamma t/2}q\)
obeys an undamped oscillator equation with shifted frequency.}
\label{fig:lab-damped-oscillator-rescaling}
\end{figure}
\FloatBarrier

\section{Rotating Hoop: Effective Potential And Bifurcation}
\label{sec:lab-rotating-hoop}

A bead of mass \(m\) moves without friction on a hoop of radius \(a\). The hoop
is vertical and rotates about its vertical diameter with imposed angular speed
\(\omega\). Let \(\theta\) be the angle of the bead measured from the downward
vertical. The bead position has speed
\[
  v^2=a^2\dot\theta^2+a^2\omega^2\sin^2\theta.
\]
The Lagrangian is
\[
  L
  =
  \frac12ma^2\dot\theta^2
  +
  \frac12ma^2\omega^2\sin^2\theta
  +
  mga\cos\theta.
\]

\begin{figure}[htbp]
\centering
\begin{tikzpicture}[scale=1.0, >=Latex]
  \begin{scope}
    \draw[->] (0,-2.05) -- (0,2.25) node[above] {rotation axis};
    \draw[dashed] (0,-1.85) -- (0,1.85);
    \draw[thick] (0,0) circle (1.65);
    \draw[fill=red!70!black] (1.17,-1.17) circle (0.09);
    \draw[thick] (0,0) -- (1.17,-1.17);
    \draw[->] (0,-0.75) arc (-90:-45:0.75);
    \node at (0.55,-0.95) {\(\theta\)};
    \draw[->, blue!70!black] (1.17,-1.17) -- (1.75,-1.17)
      node[right] {azimuthal speed \(a\omega\sin\theta\)};
    \node[below] at (0,-1.72) {downward vertical};
    \node at (0,-2.45) {bead on rotating vertical hoop};
  \end{scope}
  \begin{scope}[xshift=5.0cm]
    \draw[->] (-2.0,0) -- (2.0,0) node[right] {\(\theta\)};
    \draw[->] (0,-1.5) -- (0,1.55) node[above] {\(V_{\mathrm{eff}}\)};
    \draw[gray!60] (-1.55,-0.05) -- (-1.55,0.05) node[below=3pt] {\(-\pi\)};
    \draw[gray!60] (1.55,-0.05) -- (1.55,0.05) node[below=3pt] {\(\pi\)};
    \draw[blue!70!black, thick, domain=-1.65:1.65, samples=140]
      plot (\x,{-0.35*sin(2*\x r)*sin(2*\x r)-0.80*cos(2*\x r)});
    \draw[red!75!black, thick, domain=-1.65:1.65, samples=140]
      plot (\x,{-0.95*sin(2*\x r)*sin(2*\x r)-0.45*cos(2*\x r)+0.25});
    \node[blue!70!black, align=center] at (-1.25,1.15)
      {\(\omega<\omega_c\)\\one minimum};
    \node[red!75!black, align=center] at (1.2,1.15)
      {\(\omega>\omega_c\)\\two minima};
  \end{scope}
\end{tikzpicture}
\caption{Rotating-hoop laboratory. The geometry fixes the kinetic term
\(a^2\omega^2\sin^2\theta\), and the effective potential changes from one
bottom minimum to two off-bottom minima when \(\omega\) crosses
\(\omega_c=\sqrt{g/a}\). The potential curves are schematic and emphasize the
bifurcation, not numerical scale.}
\label{fig:lab-rotating-hoop}
\end{figure}
\FloatBarrier

This is a one-dimensional system with effective potential
\[
  V_{\mathrm{eff}}(\theta)
  =
  -\frac12ma^2\omega^2\sin^2\theta
  -
  mga\cos\theta.
\]
Equilibria satisfy
\[
  \frac{\dd V_{\mathrm{eff}}}{\dd\theta}
  =
  ma\sin\theta(g-a\omega^2\cos\theta)=0.
\]
Thus either
\[
  \theta=0,\pi,
\]
or
\[
  \cos\theta=\frac{g}{a\omega^2}.
\]
The off-bottom equilibria exist only when
\[
  \omega\geq\omega_c=\sqrt{\frac{g}{a}}.
\]
At \(\omega=\omega_c\), the bottom equilibrium changes stability and two
symmetric equilibria appear. This is a concrete mechanics realization of a
pitchfork bifurcation in an effective potential.

\section{Rolling Contact: Contact Forces As Multipliers}
\label{sec:lab-rolling-contact}

Rolling examples force one to distinguish coordinates, constraints, and
reaction forces.

\subsection{Hoop Rolling On A Fixed Cylinder}

A hoop of mass \(m\) and radius \(r\) rolls without slipping on the outside of a
fixed cylinder of radius \(R\). Let \(\theta\) be the angular position of the
hoop center measured from the top. The center of mass moves on a circle of
radius \(R+r\). The rolling condition gives
\[
  r\dot\phi=(R+r)\dot\theta,
\]
where \(\phi\) is the spin angle of the hoop. Since the hoop has
\[
  I_{\mathrm{hoop}}=mr^2,
\]
the kinetic energy is
\[
  T=
  \frac12m(R+r)^2\dot\theta^2
  +
  \frac12mr^2\dot\phi^2
  =
  m(R+r)^2\dot\theta^2.
\]
Starting from rest at the top, energy gives
\[
  m(R+r)^2\dot\theta^2
  =
  mg(R+r)(1-\cos\theta).
\]
The radial equation is
\[
  mg\cos\theta-N=m(R+r)\dot\theta^2,
\]
where \(N\) is the normal force. Loss of contact occurs at \(N=0\). Using the
energy relation,
\[
  \cos\theta=1-\cos\theta,
\]
so
\[
  \boxed{\cos\theta=\frac12.}
\]
The hoop leaves the cylinder at \(\theta=\pi/3\) from the top.

\begin{figure}[htbp]
\centering
\begin{tikzpicture}[scale=1.0, >=Latex]
  \draw[thick, fill=gray!8] (0,0) circle (1.4);
  \draw[thick, fill=blue!7] (1.65,1.65) circle (0.45);
  \draw[dashed] (0,0) -- (1.65,1.65);
  \draw[dashed] (0,0) -- (0,2.35);
  \draw[->] (0,1.9) arc (90:45:1.9);
  \node at (0.85,1.95) {\(\theta\)};
  \draw[->, red!75!black, thick] (1.65,1.65) -- (2.15,2.15)
    node[right] {\(N\)};
  \draw[->, black!70, thick] (1.65,1.65) -- (1.65,0.75)
    node[below] {\(mg\)};
  \draw[->, blue!70!black, thick] (2.1,1.65) arc (0:65:0.45)
    node[right] {\(\dot\phi\)};
  \node[align=center] at (0,-1.95)
    {loss of contact occurs when the\\normal force reaches \(N=0\)};
  \node at (-1.1,-1.1) {fixed cylinder \(R\)};
  \node at (2.7,2.25) {rolling hoop \(r\)};
\end{tikzpicture}
\caption{Rolling-contact laboratory. The radial force balance must be kept
until the end because the contact force \(N\) is the diagnostic for leaving the
surface. Eliminating the rolling constraint gives the motion, but \(N=0\)
decides the loss-of-contact angle.}
\label{fig:lab-rolling-contact}
\end{figure}
\FloatBarrier

\subsection{Solid Cylinder Inside A Hollow Cylinder}

Now consider a solid cylinder of radius \(R\) rolling inside a fixed hollow
cylinder of radius \(2R\). The center of mass moves on a circle of radius
\(R\). With rolling imposed, the kinetic energy is
\[
  T
  =
  \frac12MR^2\dot\theta^2
  +
  \frac14MR^2\dot\theta^2,
\]
because a solid cylinder has moment of inertia \(I=(1/2)MR^2\). Taking the
potential as \(V=-MgR\cos\theta\), the Lagrangian is
\[
  L=\frac34MR^2\dot\theta^2+MgR\cos\theta.
\]
The equation of motion is
\[
  \frac32MR^2\ddot\theta+MgR\sin\theta=0,
\]
or
\[
  \frac32\ddot\theta+\frac{g}{R}\sin\theta=0.
\]
For small oscillations,
\[
  \boxed{
  \omega_{\mathrm{small}}=\sqrt{\frac{2g}{3R}}.
  }
\]

\section{Schwarzschild-Type Central-Force Stability}
\label{sec:lab-schwarzschild-stability}

Consider the central Lagrangian
\[
  L=-m\sqrt{
    A(r)-\frac{\dot r^2}{A(r)}-r^2\dot\theta^2
  },
  \qquad
  A(r)=1-\frac{R}{r},
  \qquad r>R.
\]
This model is a relativistic-looking laboratory for circular-orbit stability.
Define
\[
  B=A-\frac{\dot r^2}{A}-r^2\dot\theta^2.
\]
The conserved energy and angular momentum are
\[
  E=\frac{mA}{\sqrt B},
  \qquad
  L_z=\frac{mr^2\dot\theta}{\sqrt B}.
\]
Eliminating \(\dot\theta\) and \(B\) gives the radial equation in effective
potential form:
\[
  \dot r^2
  =
  \frac{A^2}{E^2}
  \left[
    E^2-A\left(m^2+\frac{L_z^2}{r^2}\right)
  \right].
\]
Thus circular orbits are critical points of
\[
  U_{\mathrm{eff}}(r)
  =
  A(r)\left(m^2+\frac{L_z^2}{r^2}\right)
\]
at fixed \(L_z\). The circular-orbit condition \(U_{\mathrm{eff}}'(r_0)=0\)
gives
\[
  L_z^2
  =
  \frac{Rm^2r_0^2}{2r_0-3R}.
\]
Hence circular orbits require \(r_0>3R/2\). Stability is determined by
\[
  U_{\mathrm{eff}}''(r_0)
  =
  \frac{2Rm^2(r_0-3R)}
  {r_0^3(2r_0-3R)}.
\]
Therefore
\[
  \boxed{
  r_0>3R \text{ is stable},\qquad
  \frac32R<r_0<3R \text{ is unstable}.
  }
\]
This is the same structural calculation used in more realistic relativistic
orbit problems: circular-orbit stability is a second-derivative question for
an effective potential at fixed angular momentum.

\begin{figure}[htbp]
\centering
\begin{tikzpicture}[scale=1.0, >=Latex]
  \draw[->] (0,0) -- (6.4,0) node[right] {\(r/R\)};
  \draw[->] (0,-0.35) -- (0,3.0) node[above] {\(U_{\mathrm{eff}}\)};
  \draw[dashed] (1.0,-0.2) -- (1.0,2.8) node[above] {\(R\)};
  \draw[dashed] (3.0,-0.2) -- (3.0,2.8) node[above] {\(3R\)};
  \draw[red!75!black, thick, domain=1.18:5.9, samples=180]
    plot (\x,{1.15+1.7*(1-1/\x)*(1+2.9/(\x*\x))-0.33*\x});
  \fill[red!75!black] (2.02,1.83) circle (1.3pt)
    node[above right] {unstable circular orbit};
  \fill[blue!70!black] (4.35,1.54) circle (1.3pt)
    node[below right] {stable circular orbit};
  \draw[blue!70!black, thick, ->] (4.05,1.25) -- (4.30,1.48);
  \draw[red!75!black, thick, ->] (1.8,1.55) -- (2.0,1.78);
  \node[align=center] at (3.25,-0.75)
    {stability is the sign of \(U_{\mathrm{eff}}''(r_0)\)\\
     at fixed angular momentum};
\end{tikzpicture}
\caption{Effective-potential picture for the Schwarzschild-type laboratory.
Circular orbits are critical points; the outer critical point is a minimum
\((r_0>3R)\), while the inner one is a maximum
\((3R/2<r_0<3R)\). The curve is schematic.}
\label{fig:lab-schwarzschild-effective-potential}
\end{figure}
\FloatBarrier

\section{Henon-Heiles As The Finite-Dimensional Chaos Bridge}
\label{sec:lab-henon-heiles}

The Henon-Heiles Hamiltonian is
\[
  H_{\mathrm{HH}}
  =
  H_0+\lambda_{\mathrm{HH}}H_1,
\]
with
\[
  H_0=
  \frac12(p_x^2+p_y^2+x^2+y^2),
  \qquad
  H_1=x^2y-\frac13y^3.
\]
For the isotropic oscillator, introduce action-angle variables
\[
  x=\sqrt{2I_1}\sin\theta_1,
  \qquad
  p_x=\sqrt{2I_1}\cos\theta_1,
\]
\[
  y=\sqrt{2I_2}\sin\theta_2,
  \qquad
  p_y=\sqrt{2I_2}\cos\theta_2.
\]
Then
\[
  H_0=I_1+I_2,
  \qquad
  \omega_0=(1,1).
\]
The first-order homological equation for a generating function \(S_1\) is
\[
  \omega_0\cdot\partial_\theta S_1
  =
  \langle H_1\rangle-H_1.
\]
In Fourier components,
\[
  i(k_1+k_2)S_{1k}=-H_{1k}.
\]
The denominator vanishes whenever
\[
  k_1+k_2=0.
\]
At first order the cubic perturbation has no average and the dangerous modes
can be handled by symmetry, but at second order the canonical transformation
generates resonant Fourier modes with \(k_1+k_2=0\). Those modes cannot be
removed by a periodic \(S_2\). This is the small-divisor obstruction in its
most concrete classroom form: the unperturbed oscillator has equal
frequencies, so resonant combinations of angles do not average away.

The Henon-Heiles system is therefore the right finite-dimensional bridge from
integrability to Hamiltonian chaos. At low energy, Poincare sections show
deformed invariant curves. At higher energy, resonance islands and chaotic
regions appear. Unlike the asteroid problem, all variables and forces are
visible in a four-dimensional phase space.

\begin{figure}[htbp]
\centering
\begin{tikzpicture}[scale=1.0, >=Latex]
  \begin{scope}
    \draw[->] (-2.1,0) -- (2.1,0) node[right] {\(x\)};
    \draw[->] (0,-1.9) -- (0,2.2) node[above] {\(y\)};
    \foreach \s in {0.45,0.75,1.05,1.35}
    {
      \draw[blue!70!black, thick, domain=0:360, samples=120, variable=\t]
        plot ({\s*cos(\t)*(1+0.12*sin(3*\t))},
              {\s*sin(\t)*(1-0.10*cos(3*\t))});
    }
    \draw[red!75!black, thick]
      (0,1.75) -- (-1.55,-0.95) -- (1.55,-0.95) -- cycle;
    \node[red!75!black, align=center] at (0,-1.45)
      {three escape\\valleys};
    \node at (0,2.55) {configuration contours};
  \end{scope}
  \begin{scope}[xshift=5.2cm]
    \draw[->] (-1.8,0) -- (1.9,0) node[right] {\(x\)};
    \draw[->] (0,-1.8) -- (0,1.9) node[above] {\(p_x\)};
    \draw[blue!70!black, thick] (0,0) ellipse (1.2 and 0.65);
    \draw[blue!70!black, thick] (0,0) ellipse (0.72 and 0.38);
    \foreach \a in {0,120,240}
      \draw[red!75!black, thick] ({0.82*cos(\a)},{0.82*sin(\a)})
        ellipse (0.28 and 0.13);
    \foreach \x/\y in {-1.3/0.9,-1.1/-0.8,-0.4/1.15,0.25/-1.2,1.0/0.75,1.25/-0.55}
      \fill[black!55] (\x,\y) circle (0.9pt);
    \node at (0,2.35) {Poincare-section cartoon};
  \end{scope}
\end{tikzpicture}
\caption{Henon-Heiles laboratory. Configuration-space contours show the
threefold cubic deformation of the oscillator potential; a Poincare section
then records the phase-space transition from invariant curves to islands and
chaotic scatter. Both panels are schematic diagnostics for the derivation.}
\label{fig:lab-henon-heiles}
\end{figure}
\FloatBarrier

\section{Equal-Mass Three-Body Laboratory}
\label{sec:lab-three-body}

For an equal-mass three-body laboratory, take three masses in the plane with
pair potential
\[
  V=-k\left(\frac1{r_{12}}+\frac1{r_{23}}+\frac1{r_{13}}\right),
\]
where
\[
  r_{ij}=\norm{r_i-r_j}.
\]
The Lagrangian is
\[
  L=\sum_{i=1}^3\frac12m\norm{\dot r_i}^2
  +
  k\left(\frac1{r_{12}}+\frac1{r_{23}}+\frac1{r_{13}}\right).
\]
The equations are
\[
  m\ddot r_i
  =
  -k\sum_{j\neq i}\frac{r_i-r_j}{\norm{r_i-r_j}^3}.
\]
The symmetric initial positions are
\[
  r_1(0)=(1,0),
  \qquad
  r_2(0)=\left(-\frac12,\frac{\sqrt3}{2}\right),
  \qquad
  r_3(0)=\left(-\frac12,-\frac{\sqrt3}{2}\right).
\]

\begin{figure}[htbp]
\centering
\begin{tikzpicture}[scale=1.15, >=Latex]
  \draw[->] (-2.0,0) -- (2.2,0) node[right] {\(x\)};
  \draw[->] (0,-1.7) -- (0,1.9) node[above] {\(y\)};
  \coordinate (r1) at (1.25,0);
  \coordinate (r2) at (-0.625,1.083);
  \coordinate (r3) at (-0.625,-1.083);
  \draw[gray!70, thick] (r1) -- (r2) -- (r3) -- cycle;
  \fill[blue!70!black] (r1) circle (2pt) node[right] {\(r_1\)};
  \fill[blue!70!black] (r2) circle (2pt) node[above left] {\(r_2\)};
  \fill[blue!70!black] (r3) circle (2pt) node[below left] {\(r_3\)};
  \draw[->, red!75!black, thick] (r1) -- ++(0,0.65);
  \draw[->, red!75!black, thick] (r2) -- ++(-0.56,-0.32);
  \draw[->, red!75!black, thick] (r3) -- ++(0.56,-0.32);
  \draw[dashed] (0,0) circle (1.25);
  \node[red!75!black, align=center] at (1.45,1.2)
    {tangent velocities\\for rotating triangle};
  \node[align=center] at (0,-1.95)
    {the perturbation lab changes the initial velocities\\and monitors
     energy, angular momentum, and separation};
\end{tikzpicture}
\caption{Equal-mass three-body laboratory. The unperturbed comparison solution
starts from an equilateral triangle with tangential velocities; perturbing this
initial data is a controlled way to watch the loss of symmetric motion.}
\label{fig:lab-three-body-initial-data}
\end{figure}
\FloatBarrier

For the rotating equilateral solution, velocities are tangent to the
circumcircle. A representative perturbed initial condition uses a speed
parameter \(v=0.8\) and perturbation \(\epsilon=0.05\), with initial velocities
such as
\[
  \dot x_1(0)=\epsilon,\qquad
  \dot x_2(0)=-\frac{\sqrt3}{2}v-\epsilon,\qquad
  \dot x_3(0)=\frac{\sqrt3}{2}v.
\]
The important lesson is not the particular number \(0.8\). It is the workflow:
derive Newton's equations from a Lagrangian, integrate the exact pairwise
forces, monitor energy and angular momentum, and compare unperturbed symmetric
motion with perturbed motion. This is the finite-mass cousin of the
Sun-Jupiter-asteroid ejection simulation.

\section{Helicopter Rotor Model As A Gauge Connection}
\label{sec:lab-helicopter-gauge}

A toy helicopter model makes the Shapere-Wilczek connection finite-dimensional.
It is the most compact laboratory version of the inertia-operator connection
derived in Section~\ref{subsec:explicit-gauge-equation}: the general shape
coordinate \(s\) is replaced by two rotor angles \((\alpha,\beta)\), and the
zero-momentum equation can be solved by inspection.
Let the body have isotropic base inertia
\[
  I_{\mathrm{body}}=I_0\mathbf 1.
\]
Let one internal rotor spin about the body \(z\)-axis with angle \(\alpha\) and
inertia \(I_1\), and let a second rotor spin about the body \(y\)-axis with
angle \(\beta\) and inertia \(I_2\). The internal shape space is the torus
\[
  \mathcal S=S^1_\alpha\times S^1_\beta.
\]
Let \(\Omega\in\R^3\) be the body angular velocity. In the body frame, the
rotor angular velocities are
\[
  \Omega+\dot\alpha e_z,
  \qquad
  \Omega+\dot\beta e_y.
\]
Keeping only the rotational kinetic energy relevant for reconstruction,
\[
  T=
  \frac12I_0\norm{\Omega}^2
  +
  \frac12I_1(\Omega_z+\dot\alpha)^2
  +
  \frac12I_2(\Omega_y+\dot\beta)^2.
\]
The body angular momentum conjugate to \(\Omega\) is
\[
  \Pi
  =
  I_0\Omega
  +
  I_1(\Omega_z+\dot\alpha)e_z
  +
  I_2(\Omega_y+\dot\beta)e_y.
\]
Set total angular momentum to zero. Componentwise,
\[
  (I_0+I_1)\Omega_z+I_1\dot\alpha=0,
\]
\[
  (I_0+I_2)\Omega_y+I_2\dot\beta=0,
\]
\[
  I_0\Omega_x=0.
\]
Therefore
\[
  \Omega
  =
  -\frac{I_1}{I_0+I_1}\dot\alpha\,e_z
  -
  \frac{I_2}{I_0+I_2}\dot\beta\,e_y.
\]
In connection notation,
\[
  g^{-1}\dot g
  =
  -\mathcal A_\alpha\dot\alpha-\mathcal A_\beta\dot\beta,
\]
with
\[
  \mathcal A_\alpha=
  \frac{I_1}{I_0+I_1}e_z,
  \qquad
  \mathcal A_\beta=
  \frac{I_2}{I_0+I_2}e_y.
\]
The components are constant, but the structure group is nonabelian:
\[
  [e_z,e_y]\neq0.
\]
Thus the curvature contains the commutator term
\[
  \mathcal F_{\alpha\beta}
  =
  [\mathcal A_\alpha,\mathcal A_\beta].
\]
A small rectangular cycle in \((\alpha,\beta)\) can therefore produce a net
body rotation even though each rotor angle returns to its original value. This
is the finite-dimensional version of falling-cat holonomy.

\begin{figure}[htbp]
\centering
\begin{tikzpicture}[scale=1.0, >=Latex]
  \begin{scope}
    \draw[->] (-0.2,0) -- (4.4,0) node[right] {\(\alpha\)};
    \draw[->] (0,-0.2) -- (0,3.1) node[above] {\(\beta\)};
    \fill[blue!8] (0.85,0.65) rectangle (3.25,2.35);
    \draw[blue!70!black, thick, ->] (0.85,0.65) -- (3.25,0.65)
      node[midway, below] {\(\Delta\alpha\)};
    \draw[blue!70!black, thick, ->] (3.25,0.65) -- (3.25,2.35)
      node[midway, right] {\(\Delta\beta\)};
    \draw[blue!70!black, thick, ->] (3.25,2.35) -- (0.85,2.35);
    \draw[blue!70!black, thick, ->] (0.85,2.35) -- (0.85,0.65);
    \node[align=center] at (2.05,3.35) {closed shape loop\\on the rotor torus};
  \end{scope}
  \begin{scope}[xshift=5.4cm]
    \draw[->, thick] (0,0) -- (1.6,0) node[right] {\(e_x\)};
    \draw[->, thick] (0,0) -- (0,1.6) node[above] {\(e_z\)};
    \draw[->, thick] (0,0) -- (-1.1,-0.8) node[left] {\(e_y\)};
    \draw[red!75!black, thick, ->] (0.7,0.55) arc (30:290:0.55);
    \node[red!75!black, align=center] at (0.85,-1.2)
      {net body rotation\\from \([\mathcal A_\alpha,\mathcal A_\beta]\)};
  \end{scope}
\end{tikzpicture}
\caption{Helicopter-rotor gauge laboratory. The connection components are
constant, but they point in noncommuting body directions. A closed rectangle in
shape space therefore has curvature from the commutator term and can produce a
net body rotation.}
\label{fig:lab-helicopter-gauge-loop}
\end{figure}
\FloatBarrier

\section{What These Laboratories Practice}
\label{sec:worked-lab-summary}

\begin{itemize}
\item Coordinate changes preserve the Euler-Lagrange form when the
transformation is regular; the same physical path can have many coordinate
descriptions.
\item Regularity of the velocity Hessian is what permits a Lagrangian system to
be rewritten as a Hamiltonian system without losing information.
\item Constraints and contact forces should be solved together. Eliminating a
constraint too early can hide the force that decides loss of contact.
\item Effective potentials turn circular-orbit and bifurcation questions into
ordinary one-variable stability tests.
\item The Henon-Heiles, three-body, and rotor examples show how quickly
integrable intuition must be supplemented by phase-space geometry and
numerical diagnostics.
\end{itemize}

\section{Problem-To-Demo Conversion Rules}
\label{sec:lab-demo-rules}

Every computational laboratory in this repository should record the following
metadata:
\[
  \text{model}+\text{variables}+\text{parameters}
  +\text{units}+\text{invariants}+\text{numerical method}.
\]
For Hamiltonian demos, the minimum checks are:
\[
  \Delta H(t),\qquad
  \Delta J(t) \text{ if angular momentum is conserved},\qquad
  \text{step-size sensitivity}.
\]
For probabilistic demos such as asteroid ejection, the output must also state:
\[
  \text{sampling distribution},\quad
  \text{random seed},\quad
  \text{time horizon},\quad
  \text{ejection criterion}.
\]
Without these, a numerical picture is not a reproducible mechanics experiment.

The same rule applies in reverse when reading a demo. The code should reveal
the equations it integrates, the units in which parameters are measured, the
events it declares, and the invariants it monitors. A plot is the last line of
the argument, not the first.

\chapter{Detailed Worked Derivations And Examples}
\label{chap:detailed-worked-derivations}

\section{Purpose Of This Chapter}
\label{sec:detailed-purpose}

The earlier chapters develop the structure of mechanics. This chapter slows
down and works through representative calculations in full. The goal is not to
collect tricks. The goal is to show how assumptions, notation, conserved
quantities, reductions, and stability criteria fit together in examples that
students can use as templates.

Each example follows the same pattern:
\[
\begin{gathered}
  \text{model assumptions}
  \longrightarrow
  \text{coordinates and Lagrangian/Hamiltonian}\\
  \longrightarrow
  \text{conserved quantities or constraints}
  \longrightarrow
  \text{reduced equation}
  \longrightarrow
  \text{interpretation}.
\end{gathered}
\]
This pattern is more important than any single example. The most common source
of confusion in advanced mechanics is mixing levels: treating a coordinate
artifact as a physical effect, treating a conserved quantity as an independent
coordinate before checking regularity, or treating a numerical picture as a
theorem. The derivations in this chapter slow down exactly at those points.

For that reason, some calculations repeat ideas that appeared earlier in a more
compressed form. The repetition is intentional. A fast derivation is good for
orientation; a slow derivation is where one sees which assumptions are doing
work, which signs come from conventions, and which quantities are actually
invariant.

\paragraph{How to compare these examples.}
Read each example as a template with replaceable parts. The central-force
example replaces a rotational degree of freedom by angular momentum. The
pendulum replaces an orbit by a phase portrait with a separatrix. The rigid top
replaces a configuration-space picture by motion on invariant tori. The
resonance and Henon-Heiles examples show how small denominators enter. The
fluid and gauge examples show how constraints produce pressure or a mechanical
connection. The calculations are different, but the logic is the same: identify
the reduced variables, derive the diagnostic equation, and interpret the
geometry of the result.

\section{Central Forces From First Principles}
\label{sec:detailed-central-forces}

The central-force problem is the prototype of reduction by symmetry. Rotational
invariance gives conservation of angular momentum; angular momentum confines
the motion to a plane; polar coordinates in that plane reduce the dynamics to a
one-dimensional radial equation with an effective potential. The later Kepler
and precession calculations are refinements of this same reduction.

The lesson is broader than celestial mechanics. Every reduction in these notes
has the same shape: use a symmetry to find a conserved quantity, use the
conserved quantity to remove a variable, and then interpret the remaining
equation on the reduced space. Central forces are the cleanest place to see the
pattern because the geometry is visible in ordinary Euclidean space.

Let a particle of mass \(m\) move in a central potential \(V(r)\), where
\[
  r=\norm{x},
  \qquad
  x\in\R^3\setminus\{0\}.
\]
The Lagrangian is
\[
  L=\frac12m\norm{\dot x}^2-V(r).
\]
Rotational invariance gives angular momentum
\[
  J=x\times p=m x\times\dot x.
\]
Since
\[
  \frac{\dd J}{\dd t}
  =
  \dot x\times m\dot x+x\times m\ddot x
  =
  x\times\left(-V'(r)\frac{x}{r}\right)=0,
\]
the vector \(J\) is constant. Because \(J\cdot x=0\) and
\(J\cdot\dot x=0\), the motion lies in the plane perpendicular to \(J\). Use
polar coordinates \((r,\theta)\) in this plane:
\[
  L=\frac12m(\dot r^2+r^2\dot\theta^2)-V(r).
\]
The angle \(\theta\) is cyclic, so
\[
  \ell=\frac{\partial L}{\partial\dot\theta}=mr^2\dot\theta
\]
is conserved. The energy is
\[
  E=\frac12m\dot r^2+\frac{\ell^2}{2mr^2}+V(r).
\]
Thus the radial motion is one-dimensional motion in
\[
  V_{\mathrm{eff}}(r)
  =
  V(r)+\frac{\ell^2}{2mr^2}.
\]
Here \(\ell\) is the massful angular momentum. To compare with the
specific-energy Kepler convention used in
Section~\ref{sec:kepler-perturbation}, divide the Hamiltonian, energy, and
angular momentum by the particle mass; then
\(\ell_{\mathrm{spec}}=\ell/m=r^2\dot\theta\). The orbit equation is unchanged
after this bookkeeping conversion, provided \(k/m\) is read as the
gravitational parameter \(GM\).

\subsection{Binet Equation}
\label{subsec:detailed-binet}

Let
\[
  u(\theta)=\frac1r.
\]
Since
\[
  \dot\theta=\frac{\ell}{mr^2}=\frac{\ell}{m}u^2,
\]
we have
\[
  \frac{\dd}{\dd t}
  =
  \frac{\ell}{m}u^2\frac{\dd}{\dd\theta}.
\]
Also
\[
  \dot r
  =
  \frac{\dd}{\dd t}(u^{-1})
  =
  -u^{-2}\dot u
  =
  -\frac{\ell}{m}u',
\]
where a prime denotes \(\dd/\dd\theta\). Differentiating once more,
\[
  \ddot r
  =
  -\frac{\ell}{m}\frac{\dd u'}{\dd t}
  =
  -\frac{\ell^2}{m^2}u^2u''.
\]
The radial equation is
\[
  m(\ddot r-r\dot\theta^2)=F(r),
  \qquad
  F(r)=-V'(r).
\]
Substitute the formulas above:
\[
  m\left(
    -\frac{\ell^2}{m^2}u^2u''
    -
    \frac1u\frac{\ell^2}{m^2}u^4
  \right)
  =
  -\frac{\ell^2}{m}u^2(u''+u)
  =
  F(1/u).
\]
Therefore
\[
  \boxed{
  u''+u
  =
  -\frac{m}{\ell^2u^2}F(1/u).
  }
\]
This is Binet's equation. It turns a central-force problem into a differential
equation for the orbit shape \(r(\theta)\).

\subsection{Kepler Conics And The Laplace-Runge-Lenz Vector}
\label{subsec:detailed-kepler-lrl}

For the Kepler potential
\[
  V(r)=-\frac{k}{r},
  \qquad
  k>0,
\]
the force is
\[
  F(r)=-\frac{k}{r^2}.
\]
Binet's equation becomes
\[
  u''+u=\frac{mk}{\ell^2}.
\]
The solution is
\[
  u(\theta)=\frac{mk}{\ell^2}\left[1+e\cos(\theta-\theta_0)\right],
\]
or
\[
  r(\theta)=\frac{\ell^2/(mk)}
  {1+e\cos(\theta-\theta_0)}.
\]
This is a conic section with focus at the force center.

The extra conserved vector is
\[
  A=p\times J-mk\frac{x}{r}.
\]
To verify conservation, write \(\hat r=x/r\). Since \(J\) is constant,
\[
  \dot A=\dot p\times J-mk\dot{\hat r}.
\]
The force law gives
\[
  \dot p=-\frac{k}{r^2}\hat r.
\]
Use
\[
  J=mx\times\dot x=mr^2\dot\theta\,\hat z=\ell\hat z
\]
in the orbital plane. Then
\[
  \dot{\hat r}=\dot\theta\,\hat\theta,
\]
and
\[
  \dot p\times J
  =
  -\frac{k}{r^2}\hat r\times(\ell\hat z)
  =
  \frac{k\ell}{r^2}\hat\theta
  =
  mk\dot\theta\,\hat\theta
  =
  mk\dot{\hat r}.
\]
Hence
\[
  \dot A=0.
\]
The vector \(A\) points toward periapse. Taking the dot product with \(x\),
\[
  A\cdot x
  =
  (p\times J)\cdot x-mkr.
\]
By cyclic permutation,
\[
  (p\times J)\cdot x=J\cdot(x\times p)=J\cdot J=\ell^2.
\]
Therefore
\[
  A r\cos(\theta-\theta_0)=\ell^2-mkr,
\]
which rearranges to
\[
  r=\frac{\ell^2/(mk)}{1+(A/(mk))\cos(\theta-\theta_0)}.
\]
Thus the eccentricity is
\[
  e=\frac{\norm A}{mk}.
\]
The conic equation and the hidden conservation law are the same fact written
in two languages.

\subsection[Small Periapse Precession]{Small Periapse Precession From A
\texorpdfstring{\(1/r^3\)}{1/r^3} Perturbation}
\label{subsec:detailed-precession}

A useful perturbation of the Kepler potential is
\[
  V(r)=-\frac{k}{r}-\frac{\epsilon}{r^3},
  \qquad
  0<|\epsilon|\ll1.
\]
The radial force is
\[
  F(r)=-V'(r)=-\frac{k}{r^2}-\frac{3\epsilon}{r^4}.
\]
Binet's equation gives
\[
  u''+u
  =
  \frac{mk}{\ell^2}
  +
  \frac{3m\epsilon}{\ell^2}u^2.
\]
For a nearly circular orbit, write
\[
  u(\theta)=u_0+\eta(\theta),
  \qquad
  |\eta|\ll u_0.
\]
The constant \(u_0\) solves
\[
  u_0=\frac{mk}{\ell^2}+\frac{3m\epsilon}{\ell^2}u_0^2.
\]
Linearizing in \(\eta\),
\[
  \eta''+\eta
  =
  \frac{6m\epsilon u_0}{\ell^2}\eta,
\]
or
\[
  \eta''+\kappa^2\eta=0,
  \qquad
  \kappa^2=1-\frac{6m\epsilon u_0}{\ell^2}.
\]
The radial oscillation has angular period
\[
  \Delta\theta_r=\frac{2\pi}{\kappa}.
\]
The periapse advance per radial cycle is the excess over \(2\pi\):
\[
  \Delta\varpi
  =
  2\pi\left(\frac1\kappa-1\right)
  \approx
  \frac{6\pi m\epsilon u_0}{\ell^2}.
\]
Thus the \(1/r^3\) correction breaks the closed Kepler ellipse and produces
precession. This is the simplest analytic counterpart to the numerical
precession demo.

\section{Pendulum Phase Space And The Separatrix}
\label{sec:detailed-pendulum}

The simple pendulum is the canonical one-degree-of-freedom Hamiltonian example.
Let \(\theta\) be the angle from the downward vertical. The Lagrangian is
\[
  L=\frac12m\ell^2\dot\theta^2+mg\ell\cos\theta.
\]
The momentum and Hamiltonian are
\[
  p_\theta=m\ell^2\dot\theta,
  \qquad
  H(\theta,p_\theta)
  =
  \frac{p_\theta^2}{2m\ell^2}-mg\ell\cos\theta.
\]
Hamilton's equations are
\[
  \dot\theta=\frac{p_\theta}{m\ell^2},
  \qquad
  \dot p_\theta=-mg\ell\sin\theta.
\]
The equilibria are
\[
  (\theta,p_\theta)=(0,0)
  \quad\text{and}\quad
  (\pi,0)
\]
modulo \(2\pi\). The downward equilibrium is elliptic; the upward equilibrium
is hyperbolic.

The separatrix is the energy level through the hyperbolic point:
\[
  H=mg\ell.
\]
On this level,
\[
  \frac{p_\theta^2}{2m\ell^2}
  =
  mg\ell(1+\cos\theta)
  =
  2mg\ell\cos^2\frac{\theta}{2}.
\]
Therefore
\[
  p_\theta
  =
  \pm2m\ell\sqrt{g\ell}\cos\frac{\theta}{2}
\]
for \(-\pi\leq\theta\leq\pi\). The separatrix divides libration orbits, which
oscillate around \(\theta=0\), from rotation orbits, which wind around the
circle. This picture is the one-degree-of-freedom ancestor of separatrix
splitting in perturbed Hamiltonian systems.

\section[Oscillator Actions In Polar Coordinates]{Action Variables For The
Isotropic Oscillator In Polar Coordinates}
\label{sec:detailed-oscillator-actions}

A canonical test case for action-angle coordinates is the two-dimensional
isotropic oscillator
\[
  H=
  \frac{p_r^2}{2m}
  +
  \frac{p_\theta^2}{2mr^2}
  +
  \frac12m\omega^2r^2.
\]
The angular momentum
\[
  L=p_\theta
\]
is conserved. The angular action is immediate:
\[
  I_\theta=\frac1{2\pi}\oint p_\theta\,\dd\theta=L.
\]
The radial action is
\[
  I_r
  =
  \frac1{\pi}
  \int_{r_-}^{r_+}
  \sqrt{
    2mE-\frac{L^2}{r^2}-m^2\omega^2r^2
  }\,\dd r.
\]
Rather than evaluate this integral by brute force, use a structural argument.
For the Cartesian oscillator,
\[
  H=\omega(I_x+I_y).
\]
The angular momentum generates rotations in the plane. For nonnegative \(L\),
the circular orbit has all its action in rotation and no radial oscillation.
On a circular orbit,
\[
  I_r=0,
  \qquad
  E=\omega L.
\]
Radial oscillation has twice the angular oscillator frequency: one full radial
cycle occurs while the Cartesian oscillator runs through half of its ellipse.
Therefore the Hamiltonian in polar actions must be
\[
  \boxed{
  H=\omega(2I_r+|I_\theta|).
  }
\]
Solving for the radial action,
\[
  \boxed{
  I_r=\frac{E}{2\omega}-\frac{|L|}{2}.
  }
\]
This formula is a useful reminder: action variables depend on the chosen basis
of cycles. Cartesian actions and polar actions are related by an integral
linear change of torus cycles, not by an arbitrary smooth coordinate change.

\section{Symmetric Free Top: Actions And Dense Motion}
\label{sec:detailed-symmetric-top-actions}

For a symmetric free rigid body with
\[
  I_1=I_2>I_3,
\]
use Euler angles \((\phi,\theta,\psi)\). The conserved quantities are
\[
  H,\qquad J^2,\qquad J_z.
\]
The momentum conjugate to \(\psi\) is the body-axis component of angular
momentum:
\[
  p_\psi=M_3.
\]
The momentum conjugate to \(\phi\) is the spatial vertical component:
\[
  p_\phi=J_z.
\]
For the symmetric top,
\[
  H
  =
  \frac{J^2-M_3^2}{2I_1}
  +
  \frac{M_3^2}{2I_3}.
\]
Solving for \(M_3^2\) gives
\[
  M_3^2
  =
  \frac{2I_1I_3}{I_1-I_3}
  \left(
    H-\frac{J^2}{2I_1}
  \right).
\]
Thus one action is
\[
  I_\psi=M_3
  =
  \sqrt{
  \frac{2I_1I_3}{I_1-I_3}
  \left(
    H-\frac{J^2}{2I_1}
  \right)
  },
\]
with sign chosen by the orientation of the \(\psi\)-cycle. Another action is
\[
  I_\phi=J_z.
\]
The remaining nutation action is
\[
  \boxed{
  I_\theta=J-I_\psi.
  }
\]
This compact result has a geometric explanation. The angular momentum vector
has fixed length \(J\). The body symmetry axis precesses around \(J\), and the
spin about the body axis contributes \(I_\psi\). The residual part of the
angular-momentum sphere swept by the nutation cycle is \(J-I_\psi\).

The angle equations are
\[
  \dot\alpha_j=\frac{\partial H}{\partial I_j}.
\]
Unless the three frequencies are rationally related, a trajectory is dense on
the corresponding invariant three-torus. The rigid body therefore gives a
physical example in which dense quasi-periodic motion occurs without chaos.

\section{Driven One-Angle Hamiltonian And Resonant Normal Form}
\label{sec:detailed-driven-resonance}

Consider the one-angle time-periodic perturbation model
\[
  H(I,\theta,t)
  =
  \frac12I^2+\epsilon\cos(n\theta-\nu t).
\]
The unperturbed frequency is
\[
  \dot\theta=I.
\]
Resonance occurs when the phase
\[
  \psi=n\theta-\nu t
\]
is slow, i.e.
\[
  nI-\nu\approx0.
\]
Thus
\[
  I_{\mathrm{res}}=\frac{\nu}{n}.
\]
To make the system autonomous, introduce the extended phase-space pair
\((t,E_t)\), with Hamiltonian
\[
  K=E_t+\frac12I^2+\epsilon\cos(n\theta-\nu t).
\]
Use the canonical slow angle
\[
  \psi=n\theta-\nu t
\]
and an action \(P\) measuring displacement from resonance. Locally set
\[
  I=\frac{\nu}{n}+P.
\]
Dropping constants and terms independent of the slow dynamics, the resonant
Hamiltonian has the pendulum form
\[
  \boxed{
  K_{\mathrm{res}}(P,\psi)
  =
  \frac12P^2+\epsilon\cos\psi.
  }
\]
Up to scale factors depending on the precise canonical normalization, this is
the universal local model of a first-order resonance. Its phase portrait has
elliptic islands, hyperbolic fixed points, and a separatrix. This is why
planetary mean-motion resonances, the standard map, and the driven pendulum all
share the same local geometry.

\section{Henon-Heiles Perturbation And The First Small Divisor}
\label{sec:detailed-henon-small-divisor}

The explicit Henon-Heiles calculation is carried out in
Section~\ref{sec:lab-henon-heiles}. The point needed here is the general
normal-form lesson: when the unperturbed oscillator has
\(\omega_0=(1,1)\), Fourier modes with \(k_1+k_2=0\) are resonant, so the
homological equation cannot remove them by division. They must be retained in
the resonant normal form, where they organize the islands and separatrices seen
in Poincare sections.

\section{Pressure As A Lagrange Multiplier In Incompressible Flow}
\label{sec:detailed-pressure-multiplier}

For incompressible Navier-Stokes,
\[
  \partial_tu+u\cdot\nabla u
  =
  -\frac1\rho\nabla p+\nu\Delta u,
  \qquad
  \nabla\cdot u=0.
\]
The pressure is not determined by a separate thermodynamic equation in the
incompressible model. It is the Lagrange multiplier enforcing
\(\nabla\cdot u=0\).

Take the divergence of the momentum equation. Since
\[
  \nabla\cdot\partial_tu=\partial_t(\nabla\cdot u)=0
  \quad\text{and}\quad
  \nabla\cdot\Delta u=\Delta(\nabla\cdot u)=0,
\]
we obtain the pressure Poisson equation
\[
  \boxed{
  -\frac1\rho\Delta p
  =
  \nabla\cdot(u\cdot\nabla u)
  =
  \partial_i u_j\,\partial_j u_i.
  }
\]
The last equality follows by expanding:
\[
  \partial_i(u_j\partial_j u_i)
  =
  (\partial_i u_j)(\partial_j u_i)
  +
  u_j\partial_j(\partial_i u_i),
\]
and the second term vanishes by incompressibility. Thus \(p\) is whatever
spatial field is needed to project the nonlinear acceleration back onto the
space of divergence-free vector fields.

\section{Orr-Sommerfeld Equation With Boundary Conditions}
\label{sec:detailed-orr-sommerfeld}

For a parallel base flow
\[
  U(y)e_x,
\]
linearized incompressible perturbations satisfy
\[
  \partial_t\tilde u
  +
  U\partial_x\tilde u
  +
  \tilde v U'e_x
  =
  -\nabla\tilde p+\frac1{\mathrm{Re}}\Delta\tilde u,
  \qquad
  \nabla\cdot\tilde u=0.
\]
For two-dimensional perturbations, introduce a streamfunction
\[
  \tilde u=\partial_y\psi,
  \qquad
  \tilde v=-\partial_x\psi.
\]
The perturbation vorticity is
\[
  \tilde\omega_z
  =
  \partial_x\tilde v-\partial_y\tilde u
  =
  -\Delta\psi.
\]
Taking the curl of the linearized momentum equation gives
\[
  \left(\partial_t+U\partial_x\right)\Delta\psi
  -
  U''\partial_x\psi
  =
  \frac1{\mathrm{Re}}\Delta^2\psi.
\]
Insert the normal mode
\[
  \psi(x,y,t)=\phi(y)e^{i\alpha x+\sigma t}.
\]
Then
\[
  \Delta\psi=(D^2-\alpha^2)\phi\,e^{i\alpha x+\sigma t},
\]
and the eigenvalue equation becomes
\[
  \boxed{
  (\sigma+i\alpha U)(D^2-\alpha^2)\phi
  -
  i\alpha U''\phi
  =
  \frac1{\mathrm{Re}}(D^2-\alpha^2)^2\phi.
  }
\]
No penetration means
\[
  \tilde v=-\partial_x\psi=0
  \quad\Rightarrow\quad
  \phi=0
\]
at a wall. No slip also requires
\[
  \tilde u=\partial_y\psi=0
  \quad\Rightarrow\quad
  D\phi=0.
\]
Thus the Orr-Sommerfeld equation is fourth order and needs the four wall
conditions \(\phi=D\phi=0\) at the two boundaries.

\section{Mechanical Connection: Curvature From A Small Shape Loop}
\label{sec:detailed-connection-curvature}

Let \(s^1,s^2\) be two shape coordinates and let the zero-momentum
reconstruction equation be
\[
  g^{-1}\dot g
  =
  -\mathcal A_\alpha(s)\dot s^\alpha.
\]
For a tiny rectangular loop based at \(s_0\), first move by
\(\Delta s^1\), then \(\Delta s^2\), then back by \(-\Delta s^1\), then back
by \(-\Delta s^2\). To second order in the side lengths, the net group element
is
\[
  g_{\mathrm{loop}}
  =
  \exp(-\mathcal A_2\Delta s^2)
  \exp(-\mathcal A_1\Delta s^1)
  \exp(\mathcal A_2\Delta s^2)
  \exp(\mathcal A_1\Delta s^1)
\]
plus the correction from the variation of \(\mathcal A_\alpha(s)\) across the
loop. The Baker-Campbell-Hausdorff formula gives the commutator contribution,
while the spatial variation gives derivative terms. The result is
\[
  g_{\mathrm{loop}}
  =
  \exp\left(
    -\mathcal F_{12}\Delta s^1\Delta s^2
    +
    O(|\Delta s|^3)
  \right),
\]
where
\[
  \boxed{
  \mathcal F_{12}
  =
  \partial_1\mathcal A_2-\partial_2\mathcal A_1
  +
  [\mathcal A_1,\mathcal A_2].
  }
\]
This is why curvature is the local density of geometric phase. Even if the
connection components are constant, a nonabelian structure group can give
nonzero holonomy through the commutator term. The helicopter rotor model in
Chapter~\ref{chap:worked-problem-labs} is exactly of this type.

\section{What These Examples Teach}
\label{sec:detailed-examples-summary}

\begin{itemize}
\item Central-force problems are solved by symmetry reduction before any
special orbit equation is guessed.
\item The Kepler problem is special because its conics and its
Laplace-Runge-Lenz vector are two expressions of the same hidden symmetry.
\item Separatrices are not just special curves in a picture; they organize
nearby motion and become the templates for chaotic transport after perturbation.
\item Action variables are geometric period integrals. They are useful because
they survive coordinate changes better than local phase coordinates do.
\item Small perturbations of integrable systems do not merely change numbers;
they change the geometry of phase space through resonances and separatrices.
\item Pressure in incompressible flow and multipliers in constrained mechanics
play the same structural role: they enforce constraints.
\item Gauge curvature measures the failure of local reconstruction rules to be
path independent.
\item A worked derivation is complete only when the final equation is tied back
to the original assumptions: which terms were neglected, which quantities are
held fixed, and which variables are observable.
\end{itemize}


\backmatter
\chapter*{Further Directions}

Natural continuations include Dirac constraints and symplectic reduction,
Delaunay variables for celestial mechanics, resonance normal forms, Arnold
diffusion, heavy-top dynamics, numerical Poincare-section construction, and
structure-preserving integrators for fluids and deforming bodies. Each addition
should follow the same standard used here: state the model, define the symbols,
derive the equations, draw the geometry, and connect the mathematics to a
reproducible computation.

\chapter*{Suggested Further Reading}

For Lagrangian and Hamiltonian mechanics, Arnold's \emph{Mathematical Methods
of Classical Mechanics} and Goldstein, Poole, and Safko's \emph{Classical
Mechanics} are natural companions. For symplectic geometry, reduction, and
momentum maps, see Abraham and Marsden's \emph{Foundations of Mechanics} and
Marsden and Ratiu's \emph{Introduction to Mechanics and Symmetry}. For
perturbation theory and KAM context, Arnold, Kozlov, and Neishtadt's
\emph{Mathematical Aspects of Classical and Celestial Mechanics} and Holmes's
\emph{Introduction to Perturbation Methods} are useful next steps. For fluids
and elasticity, standard references include Batchelor's \emph{An Introduction
to Fluid Dynamics}, Drazin and Reid's \emph{Hydrodynamic Stability}, and
Marsden and Hughes's \emph{Mathematical Foundations of Elasticity}. For
geometric phases and deforming bodies, read Shapere and Wilczek's original
gauge-kinematics papers alongside modern discussions of mechanical
connections.

\end{document}
